%%
%% Thesis: Retrieval Augmented Generation in Financial QA: A Systematic Study
%% Amine Kobeissi, Universite de Montreal, 2026
%%

\documentclass[12pt,twoside,maitrise]{dms}
\usepackage{multirow}
\usepackage{booktabs}
% \usepackage{booktabs}
\usepackage{array}
\newcounter{example}
\newenvironment{example}[1]{%
   \refstepcounter{example}%
   \noindent\rule{\linewidth}{0.4pt}\\[2pt]%
   \noindent\textbf{Example \theexample\quad #1}\\[4pt]%
 }{%
   \noindent\rule{\linewidth}{0.4pt}\par\medskip%
 }
\usepackage{adjustbox} 
\usepackage{graphicx}
\usepackage{placeins}
\usepackage[labelfont=bf, width=\linewidth]{caption}
\usepackage{subcaption}
\usepackage{float}
\usepackage{makecell} % For line breaks within cells
\usepackage{tabularx}
\usepackage[utf8]{inputenc}
\usepackage[T1]{fontenc}
\usepackage{lmodern}
\DeclareSymbolFont{largesymbols}{OMX}{cmex}{m}{n}
\anglais

\def\sloppy{%
  \tolerance 500%
  \emergencystretch 3em%
  \hfuzz .5pt
  \vfuzz\hfuzz}
\sloppy
\usepackage[authoryear]{natbib}
\usepackage{graphicx,amssymb,icomma}
\usepackage{amsmath}
\usepackage{algorithm}
\usepackage{algorithmic}
\usepackage{booktabs}
\usepackage{xcolor}
\usepackage{tcolorbox}
\tcbuselibrary{breakable,skins}
\usepackage{hyperref}
\hypersetup{colorlinks=true,allcolors=black}
\usepackage{hypcap}
\usepackage{bookmark}

\newtheorem{cor}{\corollaryname}[section]
\newtheorem{theo}[cor]{\theoremname}
\newtheorem{prop}[cor]{Proposition}
\newtheorem{lem}[cor]{\lemmaname}
\theoremstyle{definition}
\newtheorem{deff}[cor]{\definitionname}
\newtheorem{ex}[cor]{\examplename}
\newtheorem{rem}[cor]{\remarkname}
\newtheorem{algo}[cor]{\algoname}

\numberwithin{equation}{section}
\numberwithin{table}{chapter}
\numberwithin{figure}{chapter}
\usepackage{mdframed}
\usepackage{enumitem}

\newcounter{failex}
\newenvironment{failex}[1]{%
  \refstepcounter{failex}%
  \begin{mdframed}[
    linewidth=0.8pt,
    linecolor=black,
    innertopmargin=10pt,
    innerbottommargin=10pt,
    innerleftmargin=14pt,
    innerrightmargin=14pt,
    skipabove=\baselineskip,
    skipbelow=\baselineskip,
  ]
  \noindent\textbf{Example \thefailex\enspace|\enspace #1}
  \par\medskip
}{%
  \end{mdframed}%
}

\renewcommand{\baselinestretch}{1.286}

%%%%%%%%%%%%%%%%%%%%%%%%%%%%%%%%%%%%%%%%%%%%%%%%%%%%%%%%%%%%
\begin{document}

\version{1}
\pagenumbering{roman}

\title{Retrieval Augmented Generation in Financial QA: A Systematic Study}
\author{Amine Kobeissi}
\copyrightyear{2026}
\department{Département d’informatique et de recherche opérationnelle}
\date{\today}
\sujet{Intelligence Artificielle}
\president{Jian-Yun Nie}
\directeur{Philippe Langlais}
\membrejury{Bang Liu}

\pdfbookmark[chapter]{Couverture}{PageUn}
\maketitle
\maketitle

\francais
\chapter*{R\'esum\'e}
Cette th\`ese \'etudie la recherche d'information dans les syst\`emes de
g\'en\'eration augment\'ee par r\'ecup\'eration appliqu\'es aux questions
financi\`eres. Dans ce domaine, les r\'eponses doivent \^etre fond\'ees sur des
documents primaires, souvent des rapports SEC longs, r\'ep\'etitifs et riches en
tableaux. Le travail montre que la principale difficult\'e n'est pas seulement de
retrouver le bon document, mais de trouver la bonne page \`a l'int\'erieur de ce
document. Nous proposons d'abord un cadre d'\'evaluation avec conditions
oracles qui s\'epare les erreurs au niveau du document, de la page et du
passage. Appliqu\'e \`a FinanceBench, ce cadre montre un \'ecart important entre
la reconnaissance du document et la reconnaissance de la page.

Pour r\'eduire cet \'ecart, la th\`ese introduit un pipeline hi\'erarchique qui
classe d'abord les pages avec un bi-encodeur \`a param\`etres partag\'es, puis
effectue la recherche de passages seulement dans les pages les mieux class\'ees.
Le mod\`ele est adapt\'e au domaine financier et am\'eliore le rappel des pages
ainsi que la qualit\'e des r\'eponses num\'eriques. La th\`ese pr\'esente aussi
une extension de benchmark bas\'ee sur les questions LoFIN/HiREC provenant de
FinQA et SEC-QA. Seule la partie FinQA a pu \^etre align\'ee \`a des pages PDF,
car les donn\'ees SEC-QA disponibles ne contenaient pas de texte d'\'evidence
suffisant pour faire la correspondance avec une page de r\'ef\'erence. Enfin,
nous analysons l'effet du d\'ecoupage en segments et comparons l'extraction de
texte avec une approche visuelle de type ColPali.

\textbf{Mots-cl\'es :} g\'en\'eration augment\'ee par r\'ecup\'eration,
question-r\'eponse financi\`ere, FinanceBench, FinQA, recherche hi\'erarchique,
rapports SEC, PDF.

\anglais
\chapter*{Abstract}
This thesis studies retrieval-augmented generation for financial question
answering over financial documents. In this setting, answers must be grounded in
primary-source documents, and the main practical challenge is often retrieval. The system must find a specific page or table inside a long, repetitive filing
before a language model can answer reliably. The thesis first introduces an
oracle-based evaluation framework that decomposes retrieval failure into
document-level, page-level, and chunk-level components. Applied to FinanceBench,
the analysis shows that current systems can often identify the correct filing
but still miss the evidence page, making within-document retrieval the main
bottleneck.

To address this bottleneck, the thesis proposes a hierarchical page-then-chunk
retrieval pipeline. A fine-tuned shared-parameter bi-encoder first scores pages
for relevance, and standard dense chunk retrieval is then restricted to the
highest-scoring pages. This design reduces noise from irrelevant sections and
improves page recall and downstream numerical answer accuracy. The thesis also
documents a benchmark-extension pipeline based on the LoFIN/HiREC benchmark,
which is composed of FinQA and SEC-QA. The FinQA portion was aligned to PDF
pages using fuzzy evidence matching, while the SEC-QA portion could not be added
because no gold page text was available to map questions to gold PDF pages. In
addition, the thesis evaluates chunking choices, compares PDF text extraction
methods, and explores visual page retrieval with ColPali as an alternative to
text-only retrieval.

\textbf{Keywords:} retrieval-augmented generation, financial question answering,
FinanceBench, FinQA, hierarchical retrieval, SEC filings, PDF retrieval.

\francais
\cleardoublepage
\pdfbookmark[chapter]{\contentsname}{toc}
\tableofcontents
\cleardoublepage
\phantomsection
\listoftables
\cleardoublepage
\phantomsection
\listoffigures

\anglais
\chapter*{List of acronyms and abbreviations}
\begin{twocolumnlist}{.2\textwidth}{.7\textwidth}
  BCE   & Binary Cross-Entropy loss \\
  BM25  & Best Match 25, a sparse term-based retrieval scoring function \\
  DPR   & Dense Passage Retrieval \\
  EDGAR & Electronic Data Gathering, Analysis, and Retrieval system of the SEC \\
  FT    & Fine-Tuned \\
  HyDE  & Hypothetical Document Embeddings \\
  LLM   & Large Language Model \\
  MNR   & Multiple Negatives Ranking loss \\
  NLP   & Natural Language Processing \\
  QA    & Question Answering \\
  RAG   & Retrieval-Augmented Generation \\
  RRF   & Reciprocal Rank Fusion \\
  SEC   & U.S.\ Securities and Exchange Commission \\
  VLM   & Vision-Language Model \\
\end{twocolumnlist}

\chapter*{Acknowledgements}
First and foremost, I would like to express my deepest gratitude to my supervisor, Philippe Langlais, for his support and guidance throughout this journey. His flexibility, patience, and insightful advice were key in bringing this work to life.

My time at the University of Montréal would not have been the same without the vibrant academic environment of the RALI lab. To all the lab members: thank you for the stimulating discussions, the shared insights, and the camaraderie that made the daily research process so much more rewarding. 

\NoChapterPageNumber
\cleardoublepage
\pagenumbering{arabic}

%%%%%%%%%%%%%%%%%%%%%%%%%%%%%%%%%%%%%%%%%%%%%%%%%%%%%
%%                  INTRODUCTION                    %%
%%%%%%%%%%%%%%%%%%%%%%%%%%%%%%%%%%%%%%%%%%%%%%%%%%%%%

\chapter*{Introduction}

The financial industry runs on documents. Each quarter, thousands of publicly traded
companies file detailed disclosures, accumulating a vast record of
annual reports, quarterly filings, earnings transcripts, and special notices that
constitutes the primary evidence base for analysts, investors, and compliance
professionals. Answering precise questions about these filings (documents) is not only useful but
essential to work in finance. Such questions range from what a company's
revenue was in a given fiscal year, to how management characterized a particular risk
exposure, to what the terms of a recent acquisition were. Until recently, answering such
questions required a human reader to navigate these documents manually, a process that is
time-consuming, expensive, and prone to the errors that come with searching through
hundreds of dense pages under pressure.

Large language models have recently introduced a new possibility, namely a system that
reads a question in natural language and produces an answer drawn directly from
the relevant context. The paradigm that makes this practical is retrieval-augmented
generation (RAG), which connects the reasoning and fluency of a generative model with a
retrieval component that identifies the most relevant passages from an external document
corpus. The appeal for financial applications is considerable. Such a system can ground
its answers in primary sources, support citation and auditability, and remain current as
new filings are published without requiring the underlying language model to be retrained.
Yet the gap between this description and reliable system in practice is
substantial, and understanding it requires examining the retrieval component as well as the generation component.

The difficulty is not primarily one of language generation. Modern generative models are
remarkably capable of synthesizing financial information when given the right context, and we believe that the primary challenge is retrieval. The system must first locate, within a corpus that may
contain hundreds of lengthy filings, the specific page and passage containing the
evidence needed to answer a given question. Financial filings are not only long (an
annual report can exceed two hundred pages) but structurally repetitive. Revenue is discussed
in the management commentary, in the income statement, in segment breakdowns, and in the
revenue recognition footnote. The same section headers appear across companies and across
years. Dense tables present critical numerical facts in layouts that standard text
extraction handles poorly. A retrieval system must navigate these challenges to surface
not just a topically relevant passage but the specific context a particular question
requires to generate an adequate answer.

This thesis takes the retrieval challenge as its central object of study. Its starting
point is an observation that has received limited (systematic) attention in the literature.
Namely, retrieval failure in long-document financial question answering is not a single
phenomenon. A pipeline can fail because it does not retrieve the correct document, or
because it retrieves the correct document but misses the relevant page within it, or
because it locates the correct page but still fails to surface the most informative
passage. Each failure mode has a different cause and different implications for how a
system should be improved.


The first contribution of this thesis is an oracle-based evaluation framework that
decomposes retrieval failure across multiple levels of granularity. By restricting the
search space to the correct filing (oracle document) or to the correct pages within it
(oracle page), the framework establishes empirical upper bounds at each level and measures
the gap between current systems and those bounds. The framework further quantifies generative performance by measuring the maximum answer accuracy gain that better page retrieval alone can deliver, with the generator held fixed. Applied to
FinanceBench~\citep{islam2023financebench}, the primary benchmark for PDF-native financial
question answering, this analysis reveals a consistent finding across retrieval methods
spanning sparse, dense, hybrid, and hierarchical approaches. While current systems achieve
reasonably high success at identifying the correct filing, their ability to find the
relevant page within that filing falls substantially short of the theoretical ceiling. Thus the
dominant bottleneck is within-document page retrieval, not document discovery.
 
The second contribution attempts to address this bottleneck directly. This thesis introduces a
hierarchical page-then-chunk retrieval pipeline centred on a learned page scorer. In Stage 1,
a bi-encoder fine-tuned via cross-validation on FinanceBench page-level relevance selects the top-$P$ most relevant pages from the full corpus. In Stage 2,
standard dense chunk retrieval is restricted to chunks from those pages, reducing noise from irrelevant filing sections while preserving fine-grained selection
within the most promising regions. The thesis further introduces a fine-tuned reranker, initialized from a general-purpose reranking checkpoint and
trained on Out-Of-Distribution (OOD) FinQA \citep{chen2021finqa} financial page data, which is then applied zero-shot to the held-out
FinanceBench benchmark. Together, these two components produce the strongest retrieval configuration evaluated in this thesis.
Within this same pipeline architecture, visual page scoring with ColPali \citep{faysse2024colpali} is explored as an alternative Stage 1 ranker, replacing the text-based bi-encoder with a vision-language model that represents pages as patch embeddings rather than extracted text tokens. Visual retrieval with VLMs remains largely unstudied in the financial domain, and this experiment is included to assess whether bypassing the text extraction bottleneck offers any advantage for SEC filings.
 
The third contribution is a systematic evaluation of retrieval and generation performance
across multiple dimensions. On the retrieval side, all configurations are evaluated at
multiple levels of granularity, document recall, page recall, and context quality, and
broken down by question type and filing type to surface where different methods succeed
and fail. A qualitative failure analysis then examines the recurring reasons behind
retrieval failures. On the generation side,
the thesis evaluates answer quality across generator sizes and question types, and
analyses the residual errors that persist even under perfect retrieval conditions.
 
Finally, the fourth contribution broadens the scope of evaluation through a benchmark extension pipeline. The thesis develops a procedure for aligning questions from the LoFIN/HiREC benchmark~\citep{choe2025hirec}, which is composed of two existing financial question answering (QA) datasets, to gold PDF pages in SEC filings retrieved from EDGAR. The aligned result serves as a starting point toward a larger evaluation set that goes beyond FinanceBench's 150 questions, and offers a preliminary test of whether retrieval performance holds on a separate question source. The scope of the current alignment is however limited by practical constraints discussed in the body of the thesis, and does not yet achieve the document and question type diversity that motivates the extension.



%%%%%%%%%%%%%%%%%%%%%%%%%%%%%%%%%%%%%%%%%%%%%%%%%%%%%
%%              CHAPTER 1: LITERATURE REVIEW        %%
%%%%%%%%%%%%%%%%%%%%%%%%%%%%%%%%%%%%%%%%%%%%%%%%%%%%%

\chapter{Literature Review}
\label{ch:lit}

\section{Introduction}

This chapter surveys the research landscape underpinning the contributions of this
thesis. The narrative follows the logical progression of the problem itself, beginning
with retrieval-augmented generation as a paradigm, tracing where it came from, how it
has evolved, and where it currently stands. It then turns to the information retrieval
methods that form the technical foundation of retrieval-augmented generation systems, examining their strengths and
failure modes with the financial retrieval problem in mind. From there the focus narrows
to financial question answering as an application domain, characterizing the properties
of financial documents that make retrieval especially challenging and surveying the
benchmarks that the field has developed to measure progress. The chapter then examines strategies for retrieval over long, structured documents,
surveying hierarchical retrieval approaches from general-purpose multi-stage indexing
systems to domain-specific financial pipelines, and treating chunking as a deliberate
design variable rather than a fixed preprocessing choice. It closes with treatments of
domain adaptation for financial retrieval and visual document retrieval as an emerging
alternative paradigm.


\section{Retrieval-Augmented Generation}

\subsection{From Parametric Knowledge to Grounded Generation}

Large language models encode a remarkable breadth of world knowledge in their parameters,
acquired through exposure to vast text corpora during pre-training \citep{brown2020language}. Yet this parametric
approach to knowledge carries inherent limitations that become acute in high-stakes
application domains. A model's knowledge is frozen at the time of its training and grows
stale as the world changes. It cannot easily cite the specific sources that underlie any
given claim, making it difficult to verify or audit, and because the model must compress
everything it has learned into a fixed parameter count, it is prone to confidently
generating plausible but factually incorrect statements, a failure mode commonly called
hallucination, particularly when asked about facts that are rare, recent, or
domain-specific \citep{brown2020language, ji2023survey}. In professional domains such
as finance, where users need answers they can verify against primary sources and where a
subtly incorrect number can have material consequences, these limitations make RAG necessary in high-stakes domains.

Retrieval-augmented generation (RAG) addresses these limitations by separating the language model's role as a generator and reasoner from the task of knowledge storage and information retrieval. Figure~\ref{fig:rag_pipeline} illustrates the standard RAG pipeline and highlights the central role of retrieval, whose effectiveness ultimately constrains the quality of generation. Rather than asking the model to recall facts from its parameters, a RAG system maintains an external corpus of documents that can be updated, versioned, and inspected independently. At inference time, a retrieval component identifies the most relevant passages from this corpus and supplies them to the language model as context, after which the model synthesizes those passages into a coherent answer. This architecture reduces hallucination by anchoring generation in retrieved evidence, enables provenance by allowing the retrieved passages to be surfaced alongside the answer, and supports knowledge currency by allowing the document corpus to be updated without retraining \citep{lewis2020retrieval}.
\begin{figure}[t]
    \centering
    \includegraphics[width=0.9\linewidth]{figures/rag_pipeline.png}
    \caption{Standard retrieval-augmented generation (RAG) pipeline. A query is used to retrieve relevant document fragments from an external corpus, which are then provided as context to a generative model to produce an answer.}
    \label{fig:rag_pipeline}
\end{figure}
The formal framework can be expressed concisely. Let $\mathcal{D} = \{d_1, \ldots, d_{|\mathcal{D}|}\}$ be a corpus of documents,
and let $\mathcal{C}$ denote the set of all textual chunks obtained by segmenting
the documents in $\mathcal{D}$. Given a query $q$, a retrieval function $R_k$ selects the $k$ most
relevant chunks according to a scoring function $s$,
\begin{equation}
  R_k(q) = \operatorname{Top}\text{-}k_{c \in \mathcal{C}}\; s(q,\, c),
\end{equation}
and a generator $G_\theta$ produces an answer conditioned on the query and the retrieved
context,
\begin{equation}
  \hat{y} = G_\theta\!\bigl(q,\, R_k(q)\bigr).
\end{equation}
In this formulation, the quality of $\hat{y}$ is bounded by the quality of $R_k(q)$.
If the relevant evidence is not retrieved, even the most capable generator cannot produce
a correct, grounded answer. This tight coupling between retrieval quality and generation
quality is the central observation motivating the present thesis.

% [FIGURE NOTE: A diagram of the RAG pipeline is well-placed here. It should show
% a document corpus on the left decomposed into pages and chunks, a vector index in
% the centre, a user query flowing into the retriever to select top-k chunks, and
% those chunks concatenated with the query entering a generative model on the right
% to produce a response. Annotate the two main failure modes: document retrieval
% failure (wrong filing retrieved) and page/chunk retrieval failure (correct filing,
% wrong page or chunk).]

\subsection{The Origins of RAG and Open-Domain QA}

The formal introduction of the RAG framework is typically attributed to
\citet{lewis2020retrieval}, who combined a pre-trained neural retriever with a
sequence-to-sequence generator and demonstrated substantial improvements over
closed-book baselines on knowledge-intensive natural language processing tasks. Their
formulation treats retrieved passages as a latent variable that the generator marginalizes
over, establishing a probabilistic grounding for the retrieve-then-generate pipeline that
subsequent work has largely preserved in spirit even while modifying its implementation
substantially.

This architecture built directly on the work of
open-domain question answering research that had already demonstrated the value of
retrieval-based approaches \citep{chen2017reading, lee2019latent}. What the integration
with large pre-trained language models added was a qualitative shift in generation
quality. Rather than extracting a span from a single retrieved passage, the model could
synthesize information from multiple passages, paraphrase, aggregate, and reason in ways
that earlier extractive systems could not. Fusion-in-Decoder
\citep{izacard2021leveraging} extended this insight by encoding each passage separately
and fusing them in the decoder's attention, enabling more robust multi-passage reasoning
and providing a strong generation backbone that many subsequent RAG systems have built
upon.

\subsection{The Evolution of RAG Paradigms}

The simplest implementations of retrieval-augmented generation, which the survey paper from
\citet{gao2023retrieval} characterise as Naive RAG, follow a three-stage pipeline.
Documents are indexed offline as dense vector representations, a query is encoded at
inference time and matched to the nearest indexed vectors, and the top-$k$ retrieved
passages are concatenated with the query and passed to the language model for generation.
This design is appealing for its simplicity and modularity, and it achieves strong
results on general-domain benchmarks where queries are short, documents are compact, and
relevant passages are relatively easy to distinguish from irrelevant ones.

In more demanding settings, the limitations of this approach become apparent. Chunking
documents into fixed-size windows can fragment evidence that spans multiple passages,
leaving no single chunk containing all the information the generator needs. Retrieving
by vector similarity favours passages that are topically related to the query but may
not contain the specific evidence it requires, a problem that grows more severe as
documents become longer and more internally repetitive. And because the retriever and
generator are trained independently, the retriever has no direct signal about which
passages are actually useful for downstream generation. These limitations set the
research grounds for what \citet{gao2023retrieval} call Advanced RAG.

Advanced RAG systems introduced targeted interventions at both the beginning and the end
of the retrieval process. Pre-retrieval optimization focuses on improving the match
between the query and the indexed passages before any retrieval takes place. Query
reformulation methods expand or decompose the original question, they can generate a
hypothetical answer document whose embedding may align better with relevant passages, or
break a complex question into simpler sub-questions that can each be matched to a
specific passage \citep{gao2023precise, ma2023query, karpukhin2020dense}. Index optimization techniques
restructure the document representation to make retrieval more reliable, for instance by
preserving table structure, encoding section metadata, or building hierarchical indices
that capture document organization \citep{gao2023retrieval}.

Post-retrieval optimization intervenes after the initial candidate set has been
assembled. Reranking passes retrieved candidates through a more powerful but
computationally expensive scoring model that can evaluate query-passage relevance more
accurately than the first-stage retriever \citep{nogueira2019document}. Context
compression removes passages judged redundant or irrelevant, allowing the generator to
focus on the most pertinent evidence \citep{xu2023recomp}.

More recent developments have moved toward what \citet{gao2023retrieval} describe as
Modular RAG, in which retrieval and generation components are decomposed into
independently configurable modules that can be orchestrated in diverse patterns depending
on the task. Iterative retrieval allows the system to perform multiple rounds of
retrieval interleaved with partial generation, using the model's intermediate outputs
to guide subsequent queries \citep{shao2023enhancing}. Adaptive retrieval enables the
system to decide based on the model's assessed uncertainty whether a given question
requires retrieval at all or can be answered from parametric knowledge
\citep{jiang2023active}. The current state-of-the-art is one in which the core
retrieve-then-generate pipeline is well understood and research attention has shifted
toward specialised domains, improved diagnostic evaluation, and the identification of
the specific failure modes that aggregate metrics conceal.


\section{Information Retrieval Methods for RAG}

The practical performance of a RAG system is determined in large part by the quality of
its retrieval component. This section surveys the main classes of retrieval method that
appear in the RAG literature and in the experimental evaluation of this thesis, examining
their respective strengths and failure modes with the financial retrieval problem
specifically in view.

\subsection{Sparse Retrieval}

Sparse retrieval methods represent documents and queries as high-dimensional sparse
vectors in vocabulary space, where each dimension corresponds to a term in the vocabulary
and the vast majority of values are zero. The most widely deployed sparse retrieval
function is BM25 \citep{robertson2009probabilistic}, which scores a document $D$ against
a query $Q$ as
\begin{equation}
  \text{BM25}(D, Q) = \sum_{t \in Q} \text{IDF}(t) \cdot
  \frac{f(t,D)\,(k_1+1)}{f(t,D) + k_1\!\left(1 - b + b\,\dfrac{|D|}{\text{avgdl}}\right)},
\end{equation}
where $f(t, D)$ is the frequency of term $t$ in document $D$, $|D|$ is the document
length, avgdl is the average document length in the collection, and $k_1$ and $b$ are
tuning parameters controlling term frequency saturation and length normalization
respectively. The inverse document frequency term $\text{IDF}(t)$ rewards terms that
are distinctive across the corpus, ensuring that rare terms carry more weight than
common ones.

BM25 remains a competitive baseline in many specialised retrieval settings because its
reliance on exact term matching is precisely what makes it effective when the query and
the relevant document share distinctive surface forms. Named entities, technical
identifiers, ticker symbols, and fiscal year markers are all retrieved reliably if the
query contains them verbatim. BM25 is also transparent (its scores decompose directly
to show which query terms contributed most to a document's ranking) and efficient,
relying on inverted indices that support fast retrieval over large corpora
\citep{lin2021pyserini}. Its limitation is that the
method cannot account for semantic equivalence, so a query about ``automobile
manufacturing revenue'' will not reliably retrieve a passage written in terms of
``vehicle production sales'' unless there is meaningful lexical overlap.

Learned sparse retrieval methods such as SPLADE \citep{formal2021splade} attempt to
address vocabulary mismatch while preserving the practical advantages of sparse
indexing. Using a transformer-based encoder, SPLADE assigns vocabulary-space weights to
queries and documents based on contextualized representations rather than raw term
frequencies, effectively performing implicit query expansion. The result is a sparse
representation that retains compatibility with inverted index retrieval while capturing
some semantic generalization that pure lexical methods lack.

\subsection{Dense Retrieval and the Bi-Encoder Architecture}

Dense retrieval methods represent queries and documents as compact, continuous vectors
in a shared embedding space, enabling retrieval based on semantic similarity rather than
surface-level term overlap. The dominant architecture is the bi-encoder, which encodes
queries and documents independently into fixed-dimensional dense representations and
scores relevance as their cosine similarity. In many modern retrieval models, including
the setup used in this thesis, the same parameterized encoder is applied
to both the query and chunk inputs,
\begin{equation}
  s(q,\, c) = \frac{\langle \mathbf{h}_q,\; \mathbf{h}_c \rangle}
                   {\|\mathbf{h}_q\|\;\|\mathbf{h}_c\|},
  \quad
  \mathbf{h}_q = f_\theta(q),
  \quad
  \mathbf{h}_c = f_\theta(c).
  \label{eq:cosine}
\end{equation}
Because document embeddings can be precomputed and indexed offline, retrieval reduces to
a nearest-neighbour search in embedding space, which is fast and scalable even over large
corpora. Models such as DPR \citep{karpukhin2020dense} and BGE-M3 \citep{chen2024bge}
demonstrated that pre-training on large-scale data followed by contrastive fine-tuning
on question-passage pairs produces embeddings substantially more effective for question
answering retrieval than sparse baselines on general-domain benchmarks.

The key advantage of the formulation in Equation~\eqref{eq:cosine} is its ability to
capture semantic similarity without requiring shared vocabulary. A query about
``profitability'' can match a passage discussing ``net income margins'' even if neither
term appears in the other. This semantic generalization makes dense retrievers
particularly effective for queries phrased in natural conversational language.

The contrastive training objective that underlies most dense retrieval models teaches
the encoder to place queries close to their relevant passages and far from irrelevant
ones in embedding space \citep{izacard2022unsupervised}. Negative examples play a critical role in shaping what the
model learns. Easy negatives drawn at random provide weak signal, while hard negatives
drawn from topically similar but answer-irrelevant passages teach the model
finer-grained distinctions. This distinction between easy and hard negatives becomes
especially important for domain adaptation, as discussed in
Section~\ref{sec:domain_adaptation}.

\subsection{Hybrid Retrieval and Reciprocal Rank Fusion}

Sparse and dense methods exhibit complementary failure modes. Sparse methods fail when
query and document vocabulary diverge, but succeed when queries contain rare or
distinctive terms that must appear verbatim. Dense methods succeed at semantic matching
but can be distracted by topically related passages that do not contain the specific
evidence required. Hybrid retrieval combines both to exploit their complementarity,
typically by running sparse and dense retrievers independently and fusing their ranked
lists \citep{luan2021sparse}.

Reciprocal Rank Fusion \citep{cormack2009reciprocal} provides a simple and robust fusion
strategy that focuses on rank positions rather than absolute scores, which are on
incomparable scales across retrievers,
\begin{equation}
  \text{RRF}(d) = \sum_{r \in \text{rankers}} \frac{1}{k_{\text{rrf}} + \text{rank}_r(d)},
\end{equation}
where $\text{rank}_r(d)$ is the position of document $d$ in the ranking from retriever
$r$, and $k_{\text{rrf}}$ (typically set to 60) prevents small rank differences from
dominating the score. In financial document retrieval specifically, hybrid approaches
are particularly valuable because sparse retrieval reliably captures company names,
ticker symbols, and fiscal year markers, while dense retrieval handles semantic
variations in question phrasing.

% [FIGURE NOTE: A diagram showing two parallel retrieval tracks, a BM25 sparse
% retriever on the left and a dense bi-encoder on the right, each producing a
% candidate ranked list, which are then fused via RRF into a single merged ranking.
% Include approximate corpus sizes at the top (full corpus) and after fusion
% (top-20 candidates) to illustrate the efficiency-precision trade-off.]

\subsection{Query Reformulation and Hypothetical Document Embeddings}

A structural challenge in retrieval for question answering is the asymmetry between the
short, high-level phrasing of a natural language question and the long, detailed phrasing
of a relevant passage. A query of the form ``What was the company's FY2022 net revenue?''
may share few tokens with the income statement row it is trying to locate. Query
reformulation methods address this asymmetry by transforming the query before retrieval
to better match the style and vocabulary of relevant passages.

Hypothetical Document Embeddings (HyDE) \citep{gao2023precise} address this problem
directly. Rather than embedding the question itself, HyDE prompts a language model to
generate a hypothetical passage that would answer the question, then uses that
hypothetical passage's embedding for retrieval. The generated passage does not need to
contain accurate facts, it simply needs to be phrased in a style and vocabulary that aligns with the
documents in the corpus. For financial queries, this means the hypothetical passage is
written in the tone of a financial report, with appropriate terminology and sentence
structures, which can substantially improve the match between the query embedding and
the embeddings of genuine evidence passages. Multi-HyDE, a variant evaluated in this
thesis, averages the embeddings of multiple independently sampled hypothetical passages
to reduce variance and improve retrieval stability, as expressed by
\begin{equation}
  \mathbf{h}_{\text{Multi-HyDE}}(q) = \frac{1}{N}\sum_{n=1}^{N} f_\theta\!\bigl(\hat{d}_n\bigr),
\end{equation}
where $\hat{d}_1, \ldots, \hat{d}_N$ are independently sampled hypothetical passages
generated by the language model given query $q$, and $f_\theta$ is the shared encoder
used by the bi-encoder.

\subsection{Reranking with Cross-Encoders}

First-stage retrieval with bi-encoders is fast and scalable but limited in its ability
to model fine-grained query-document interactions, because each input is encoded
independently. Cross-encoder rerankers address this by jointly encoding the query and
candidate passage through a transformer that attends to both simultaneously. For a
cross-encoder with parameters $\phi$, the relevance score for a query $q$ and candidate
passage $c$ is
\begin{equation}
  r_\phi(q, c) = f_\phi\!\bigl([\texttt{CLS}]\; q\; [\texttt{SEP}]\; c\; [\texttt{SEP}]\bigr),
\end{equation}
where $f_\phi$ is a transformer that operates over the full concatenated sequence and
produces a scalar relevance estimate. This enables the model to capture nuanced relevance
signals, such as matching a specific fiscal year in the query to a year in the passage
header, or recognizing that a passage discusses a related metric rather than the specific
one asked about, in ways that a bi-encoder cannot easily learn. Because cross-encoders
require a full forward pass for each query-document pair, they are too expensive for
first-stage retrieval over large corpora but well-suited as a second stage that reranks
a small set of candidates produced by a fast first-stage retriever \citep{nogueira2019document}.

In financial documents, rerankers are particularly valuable for resolving near-duplicate
sections. Different fiscal years' financial statements follow identical formats with
nearly identical language, and cross-encoder rerankers can exploit subtle contextual
cues to make the fine-grained distinctions that bi-encoders miss.

\subsection{Late Interaction Models}

A middle path between the efficiency of bi-encoders and the expressiveness of
jointly-encoding cross-encoders is represented by late interaction architectures, of
which ColBERT \citep{khattab2020colbert} is the most prominent. Rather than condensing
the entire document into a single vector, ColBERT produces a sequence of contextualized
token embeddings for each document. Relevance is then computed via the MaxSim operator,
where for each query token embedding $\mathbf{e}^q_i$ the maximum cosine similarity to
any document token embedding $\mathbf{e}^d_j$ is taken, and these per-token maxima are
summed over all query tokens,
\begin{equation}
  S_{\text{ColBERT}}(q,\, d) =
  \sum_{i=1}^{|q|} \max_{j=1}^{|d|} \cos\!\bigl(\mathbf{e}^q_i,\, \mathbf{e}^d_j\bigr).
  \label{eq:maxsim}
\end{equation}
Document token embeddings can be precomputed offline, so retrieval remains efficient
despite the richer interaction model. Equation~\eqref{eq:maxsim} is the conceptual
foundation for ColPali, the visual retrieval model discussed in
Section~\ref{sec:vlm_retrieval}.

% [FIGURE NOTE: A two-stage retrieval pipeline diagram showing fast bi-encoder
% first-stage retrieval over the full corpus producing N=20 candidates, followed
% by a cross-encoder reranker that reorders and retains the top-k=5. Annotate
% the approximate time complexity of each stage to illustrate why cross-encoders
% cannot be used as a first stage.]


\section{Financial Question Answering}
\label{sec:finqa}

\subsection{The Distinctive Challenges of Financial Documents}

Financial question answering is a demanding application domain for retrieval
systems because it requires searching over documents that exhibit multiple
challenging properties, both individually and in combination. A typical 10-K annual report spans between one hundred and three
hundred pages and follows a standardized structure mandated by regulation, covering the
business description, risk factors, financial statements, management's discussion and
analysis, and a collection of exhibits \citep{islam2023financebench}. This standardization means that similar section
headers appear across every company and every fiscal year, creating a retrieval challenge
distinct from that of a heterogeneous document corpus. A query about operating expenses
is relevant to the income statement, but the income statement in one year's 10-K is
structurally nearly identical to the income statement in every other year's 10-K from
every other company. Retrieval systems relying on topical similarity must therefore make
fine-grained distinctions that first-stage bi-encoders frequently fail to make reliably.

Financial documents are also numerically dense in a way that matters for question
answering. Many of the most important questions that analysts care about, including
revenue figures, earnings per share, capital expenditures, and debt ratios, require
extracting or computing specific numerical values from dense tabular contexts \citep{chen2021finqa}. These
values appear as entries in tables whose structure encodes essential context.

Finally, financial questions require precision in a sense that general-domain questions
do not. A question asking for FY2022 revenue must be answered with the value from the
2022 filing, not the 2021 or 2023 filing, even though all three filings are from the
same company and discuss the same concept in nearly identical language. This temporal
specificity imposes a retrieval requirement that goes beyond topical relevance and
extends to exact contextual matching of year references, quarter identifiers, and other
temporal markers. Combined with the structural repetitiveness described above, this creates a
retrieval setting in which passages from different reporting periods are often
semantically and structurally near-identical. As a result, retrieval is not
simply a matter of identifying the correct topic or year, but of resolving fine-grained distinctions within highly similar content, making accurate ranking
non-trivial.

\subsection{The Evolution of Financial QA Benchmarks}

The benchmark landscape for financial question answering reflects a progression from
closed-book settings, where the relevant document or passage is provided to the model,
toward open-book settings where the model must retrieve its own evidence from a realistic
document corpus. This progression matters for interpreting what different benchmarks
measure and what limitations they impose.

The earliest influential benchmarks in the space, including FinQA \citep{chen2021finqa}
and TAT-QA \citep{zhu2021tatqa}, focused primarily on numerical reasoning capability.
Both datasets provide question-answer pairs grounded in short document contexts that
contain the necessary evidence, essentially supplying the retrieval result to the model
and asking only whether it can reason over what has been handed to it. FinQA in
particular annotates reasoning programs that trace the multi-step calculations required
to derive each answer, making it a valuable resource for studying arithmetic reasoning.
ConvFinQA \citep{chen2022convfinqa} extended this framework to conversational settings,
introducing multi-turn dependencies where each question builds on the previous one.
These benchmarks established that language models can perform sophisticated numerical
reasoning over financial tables and narrative text, but because they supply the evidence
rather than requiring its retrieval, they do not measure the retrieval component of
financial QA.

More recent benchmarks have introduced retrieval requirements in varying forms. SEC-QA
\citep{lai2024secqa} provides a systematic evaluation corpus for financial QA over SEC
filings with an emphasis on cross-company comparisons, though its use of template-based
question generation produces questions that may not capture the heterogeneous,
narrative-driven nature of genuine analyst queries. LOFin \citep{choe2025hirec} offers a
large-scale open-domain benchmark with QA pairs spanning a very large number of SEC
documents, enabling evaluation under conditions where the system must identify the relevant company and filing period before locating specific evidence.
% Fin-RATE
% \citep{wang2025finrate} introduces fine-grained evaluation beyond answer-level
% correctness, assessing cross-company comparison, cross-year tracking, and multi-document
% reasoning with step-by-step reasoning traces that enable more diagnostic evaluation of
% failure patterns.

\subsection{FinanceBench as the Primary Evaluation Standard}

Among existing benchmarks, FinanceBench \citep{islam2023financebench} stands out as the
most directly relevant to the work of this thesis. As far as we know, it is the only publicly available open-source benchmark that requires end-to-end retrieval over real PDF filings with explicit page-level evidence annotations,
making it possible to evaluate retrieval quality at the granularity this thesis is most
concerned with.

FinanceBench comprises question-answer-evidence triplets based on real SEC filings from
forty U.S.\ publicly traded companies spanning nine GICS sectors. Each sample includes a
natural language question, a reference answer, a verbatim evidence string extracted from
the source document, and metadata identifying the source filing and the page number on
which the evidence appears. This page-level annotation is what makes FinanceBench
uniquely valuable, it provides ground-truth labels for page-level retrieval evaluation,
enabling the oracle analysis framework described in Chapter~\ref{chap:methodology} to
operate at the granularity where the critical failure mode lies.

The questions in FinanceBench fall into three types, with fifty questions of each kind
in the open-source subset. Domain-relevant questions ask about facts broadly applicable
across companies, such as dividend policies or executive compensation structures, and
often require locating narrative statements that may appear in varying sections of
different filings. Metrics-generated questions require extracting or computing specific
financial indicators from dense tabular contexts, these are especially demanding for
retrieval because the answer lies in a specific table cell, and topically similar passages
from adjacent pages or adjacent fiscal years serve as near-perfect retrieval decoys.
Novel-generated questions are company-specific and require reasoning beyond simple
extraction. The document types in the benchmark are dominated by 10-K annual reports at
approximately 75\%, with the remainder drawn from 10-Q quarterly reports, 8-K current
reports, and earnings call transcripts.

\subsection{Limitations of Existing Benchmarks and the Motivation for Extension}

While FinanceBench provides an invaluable structured testbed, its 150-question
open-source subset imposes real constraints on what can be concluded from experiments
conducted on it alone. The strong skew toward 10-K filings means that findings about
retrieval difficulty, chunking effects, or page scorer performance on annual reports may
not generalize to quarterly reports, earnings transcripts, or more specialized filing
types. The limited number of open-source questions means that stratified analysis by question type
or filing type is based on small subgroups where individual question characteristics can
have an outsized effect on aggregate metrics. Finally, the lack of separate train, dev,
and test sets limits the generalization claims of domain fine-tuning for PDF-native
benchmarks.

These limitations directly motivate the benchmark extension described in
Section~\ref{sec:method_benchmark}. At a broader level, they reflect a genuine gap in
the financial QA literature. Most benchmarks that require retrieval either operate in a
closed-book setting, assume that text has already been extracted from the original
document, have limited or no publically available data, or restrict evaluation to a single narrow filing type. A benchmark that stresses
end-to-end PDF retrieval at scale across diverse filing types and question styles remains
an open contribution, and the extension work of this thesis is a step toward filling it.

\subsection{Recent RAG Systems for Financial QA}

Several recent systems have pushed state-of-the-art performance on FinanceBench and
related benchmarks, and understanding their contributions is important for contextualizing
the contributions of this thesis. FinSage \citep{wang2024finsage} introduces a
multi-aspect RAG framework that combines domain-specific preprocessing, multi-path
retrieval over different document representations, and a reranker fine-tuned using direct
preference optimization on financial query-document pairs. The preprocessing pipeline is
particularly thoughtful, preserving table structure, encoding section metadata, and
constructing semantically coherent chunks that respect the document's visual layout, an
approach that the authors show improves retrieval precision substantially over naive
fixed-size chunking.

\citet{kim2024optimizing} propose a three-phase pipeline integrating query
expansion, hybrid retrieval with fine-tuned embedders, and a direct preference optimization (DPO) trained reranker with
document selection, achieving strong improvements in retrieval precision and downstream
answer accuracy across several financial QA datasets. These systems demonstrate that
substantial performance gains are achievable through careful engineering of the full RAG
pipeline.

What neither system provides, however, is a systematic decomposition of where their
improvements come from. Neither study separates document-level retrieval performance
from page-level retrieval performance, nor establishes oracle bounds that quantify the
headroom remaining at each level of the retrieval hierarchy. It is therefore difficult
to know from their reported results whether further investment in retrieval quality would
yield significant returns, or whether the remaining performance gap is attributable to
generation limitations that better retrieval cannot address. The oracle-based analysis
framework introduced in this thesis is designed to fill precisely this diagnostic gap.


\section{Hierarchical Retrieval Strategies for Long Documents}
\label{sec:hierarchical_retrieval}

% \subsection{The Granularity Tension in Long-Document Retrieval}

A fundamental challenge in building retrieval systems for long documents is selecting the
appropriate granularity for the units that will be indexed and retrieved. Retrieving at
the level of full documents preserves all context but makes it impractical for the
retrieval system to identify which part of a long filing contains the relevant evidence,
forcing the generator to process an enormous amount of irrelevant text. Retrieving small
chunks enables precise localization but risks fragmenting evidence that spans multiple
passages and loses the surrounding context that helps the generator interpret what it has
found. Hierarchical retrieval approaches address this tension by organizing documents
into multi-level structures, using coarser levels to quickly identify the most promising
regions of the document space and finer levels to precisely localize answer-bearing
evidence within those regions. This section surveys the main families of hierarchical
retrieval method, from general-purpose recursive indexing schemes to specialized
pipelines designed for long financial documents.

\subsection{Multi-Stage Retrieval and Hierarchical Indexing}

A prominent family of approaches addresses the granularity tension by constructing
hierarchical document representations offline, which are then searched in a
coarse-to-fine manner at retrieval time.

RAPTOR \citep{sarthi2024raptor} constructs a tree of abstractive summaries by
recursively clustering document chunks and generating summary nodes at increasing levels
of abstraction. A leaf node represents a short passage, its parent summarizes a cluster
of related passages, and the root node represents a high-level synopsis of the full
document. At retrieval time, the system can be queried at any level of the tree or across
multiple levels simultaneously, allowing high-level questions to be matched to high-level
summaries and detailed questions to be matched to fine-grained passages. On long-document
benchmarks, this design substantially improves recall on questions requiring the
synthesis of information from multiple disparate sections, because the summarization step
aggregates evidence that no individual passage contains. The principal limitation is
that recursive summarization can introduce information loss and hallucination at higher
levels of abstraction, making the quality of the intermediate summaries a critical
determinant of overall retrieval performance.

More generally, two-stage coarse-to-fine pipelines interleave a cheap first-stage
retriever over coarse units (documents or sections) with a more expensive second-stage
ranker over fine-grained units (passages or chunks) within the top-$M$ coarse candidates.
This design concentrates the second stage on a small, pre-filtered candidate pool,
allowing a computationally expensive ranker to operate efficiently while covering the
full document corpus at the first stage. The architecture introduces a new failure mode, where any relevant passage whose parent document or section is not retrieved at the
first stage cannot be recovered, making the coarse retrieval component a hard ceiling on
overall recall.

\subsection{Parent-Child Structures and Decoupled Retrieval Granularity}

A closely related pattern is the parent-child document structure, in which fine-grained
child chunks are indexed for retrieval precision while coarser parent chunks are returned
to the generator to provide broader context \citep{gao2023retrieval}. After retrieving
the top-$k$ child chunks by similarity to the query, the system maps each child to its
parent, deduplicates, and returns the unique parent chunks as the context for generation.
This decoupling of retrieval granularity from generation granularity addresses a challenge, the optimal unit size for retrieval is often not the optimal unit size for
generation, because retrieval benefits from short, discriminative passages while
generation benefits from longer passages that supply complete context. In practice, a
common instantiation uses sentences or short fixed-size spans as children and paragraphs
or pages as parents, though the appropriate granularity depends heavily on the document
type and the nature of the questions being asked.

The parent-child pattern is relevant to long financial documents because the evidence
for a single question often occupies a coherent physical unit such as a financial
statement page or a management discussion subsection, while precise retrieval requires
matching the query against specific sentences or table rows within that unit. Indexing
at both granularities, without committing to either, is a design principle that recurs
across several recent financial retrieval systems.

\subsection{Chunking as a Granularity Decision}
\label{sec:chunking_lit}

Chunking sits at the boundary between document processing and retrieval design, and is
most naturally understood as the decision that defines the leaf nodes of the retrieval
hierarchy. Its importance is widely acknowledged, but it is rarely treated as a primary
variable of study. Most financial RAG papers fix a chunking configuration early in system design
and optimize other components without varying it, so the literature offers many data
points about what retrieval performance is achievable with a given strategy but little
evidence about the interactions between chunking choices and retrieval architecture.

The most common approach, fixed-size chunking, divides document text into overlapping
windows of a fixed token count. Given a document of length $n$ tokens, a sliding window of size $w$ with overlap $o$ (and stride $s = w - o$) produces $\left\lceil \frac{n - o}{w - o} \right\rceil$ chunks, where each chunk begins $s$ tokens after the previous one. This approach requires no document understanding and
is simple to scale, but it treats documents as undifferentiated character
streams. In financial filings, where a table broken at a row boundary loses the
correspondence between values and their column headers, these distortions can directly
degrade retrieval quality.

Structure-aware chunking strategies attempt to respect the document's natural
organization. Section-boundary preservation prevents chunks from crossing the mandated
section divisions of SEC filings. Table-aware splitting treats tabular regions as atomic
units, keeping column headers co-located with the data rows they describe. Boilerplate
detection identifies legal disclaimers and standardized caveats that can dilute retrieval
quality if indexed without downweighting. \citet{wang2024finsage} demonstrate that
structure-aware preprocessing substantially improves retrieval precision on FinanceBench
relative to naive fixed-size chunking, with the largest gains on metrics questions where
precise table retrieval is decisive. \citet{jimeno2024financial} similarly find that
preserving table structure within chunk text improves downstream numeric answer accuracy.

The interaction between chunk size and retrieval architecture is the aspect of chunking
that has received least systematic attention. Larger chunks preserve more co-located
context, smaller chunks produce more discriminative embeddings that better distinguish
between adjacent passages. This trade-off also depends on whether a hierarchical stage has already narrowed the candidate pool. When retrieval is restricted to a small set of pre-selected pages, the optimal within-page chunk size may differ from what is optimal for global search over an entire corpus.

% [FIGURE NOTE: A single document page shown three ways side by side: as it appears
% in the PDF (with its visual table structure intact), after full-page text
% extraction (column structure collapsed), and split into 256-token vs 1024-token
% chunks. This figure would make concrete the trade-off between chunk size and
% evidence completeness, and the information loss from text extraction.]

\subsection{Hierarchical Retrieval for Long Financial Documents}

Financial QA introduces additional structure that specialized hierarchical systems can
exploit. A typical end-to-end pipeline must navigate from a large corpus of filings
(thousands of documents) to a single filing, then from that filing's many pages to the
specific page containing the answer, and finally to the precise passage or table row
within that page. Treating these as separate retrieval problems rather than one
allows each stage to be optimized for its specific discriminative challenge.

The LoFIN benchmark and HiREC system  provide a direct illustration of this principle for open-domain financial QA. LoFIN is a large-scale evaluation corpus that requires a system to identify the relevant company and filing period from a collection of over 145,000 SEC documents before locating specific evidence. HiREC addresses this challenge with a hierarchical retrieval pipeline that first identifies related documents to resolve company and period ambiguity, and then selects the most relevant passages within those specific documents. This is followed by an evidence curation step that filters out irrelevant data and generates complementary queries if the initial retrieval is insufficient. The system’s explicit separation of the document retrieval process from the within-document passage selection provides a principled way to reduce confusion caused by the repetitive boilerplate content and similar table structures inherent in standardized financial reports.

Within financial documents, the page provides a practical intermediate unit between
the full filing and individual text chunks. Pages often correspond to coherent
units of presentation, such as complete financial statements or self-contained
subsections, meaning that their boundaries tend to align with meaningful semantic
structure. In contrast, section boundaries require additional parsing and are not
consistently defined across filings, while fixed-size chunking can arbitrarily
fragment structured content. Because page boundaries are explicitly defined by the document format and
consistently available across filings, they offer a stable and interpretable
granularity for retrieval. More broadly, this reflects a general
observation in retrieval systems that the choice of representation granularity
can significantly impact retrieval effectiveness \citep{izacard2021leveraging, liu2024lost}.

\section{Visual Document Retrieval}
\label{sec:vlm_retrieval}

\subsection{The Text Extraction Bottleneck}

All the retrieval methods described in the preceding sections share a fundamental
assumption, namely that documents can be adequately represented as sequences of
text tokens. For financial filings, this introduces a systematic source of
information loss. In practice, important aspects of these documents are encoded
through visual structure rather than purely textual content. Layout plays a
functional role, column alignment defines relationships between entities and
values, and formatting cues such as headings and spacing convey hierarchical
organization. Standard text extraction and parsing pipelines recover most narrative text
reliably, but they can fail on complex multi-column tables, sometimes producing
linearized representations in which values are separated from the headers that
give them meaning \citep{xu2020layoutlm}. For metrics-generated questions, whose answers often lie in
these tabular structures, such extraction errors can propagate directly to
retrieval, limiting the effectiveness of downstream models.

\subsection{Vision-Language Models and the ColPali Architecture}

Vision-language models offer an alternative approach that bypasses text extraction
entirely by processing document pages as images. A vision encoder, typically a Vision
Transformer that divides each page image into a grid of patches, represents the visual
content of the page as a sequence of patch embeddings, preserving spatial structure
rather than linearizing it \citep{dosovitskiy2020image}. These patch embeddings can then
be processed by a language model backbone to produce contextualized representations
that integrate visual layout with semantic content.

ColPali \citep{faysse2024colpali} is the most directly relevant vision-language model
for document retrieval. It combines the PaLIGemma architecture
\citep{beyer2024paligemma} with the ColBERT late-interaction retrieval formulation,
producing page-level retrieval by scoring relevance as
\begin{equation}
  S_{\text{ColPali}}(q,\, p) =
  \sum_{i} \max_{j} \cos\!\bigl(\mathbf{e}^q_i,\, \mathbf{e}^p_j\bigr),
  \label{eq:colpali_background}
\end{equation}
where $\mathbf{e}^p_j$ are the patch embeddings of page $p$ produced by the vision
encoder, and $\mathbf{e}^q_i$ are the query token embeddings. This formulation allows
individual query terms to align with specific visual regions of the page, enabling
a query about ``capital expenditures'' to align with the patches covering the relevant
table row even if OCR of that region would be imperfect. Because patch embeddings can
be precomputed offline, indexing a document corpus with ColPali is efficient and
retrieval at query time requires only the MaxSim computation over pre-indexed patch
embeddings.

On the DocVQA and ViDoRe benchmarks, ColPali substantially
outperforms text-extraction-based baselines on pages with complex table layouts and
embedded figures \citep{faysse2024colpali}, with the largest gains precisely on the content types that most
challenge PDF parsing.

\subsection{Recent Advances and Remaining Challenges}

A limitation of the flat patch-aggregation approach in ColPali is that all patches of a
page are indexed at a single level of granularity, regardless of their visual complexity
or their relevance to typical queries. Recent work has explored multi-resolution indexing
strategies that extend the patch-based paradigm to multiple spatial scales. \citet{fang2025pyramid}
propose a three-pass pyramid indexing framework that constructs separate visual
representations at page, region, and patch levels, allowing the retrieval system to match
queries at whichever visual granularity is most discriminative. On document-heavy
benchmarks, this multi-scale approach improves recall on questions that require matching
fine-grained visual elements such as specific table cells or column labels, where
full-page patch aggregation tends to dilute the relevant signal across too many patches.

Despite these advances, visual retrieval methods face a set of challenges that remain
partially unresolved. The resolution at which VLMs process document pages is constrained
by the input size of the vision encoder, many financial tables use small fonts and dense
layouts that are difficult to read at the resolutions typically used for training. Domain
mismatch is also a concern, since most VLM-based retrieval models are trained on general
document corpora and may not have seen sufficient examples of the standardized typography
and table structures used in SEC filings to develop discriminative representations for
that domain. Finally, the computational cost of storing full patch embedding sequences
for large document corpora is substantially higher than storing single-vector chunk
embeddings, limiting the scalability of late-interaction visual retrieval. Whether
domain-specific fine-tuning, higher-resolution encoders, or hybrid approaches combining
text and visual signals can address these gaps is an open and active area of research.


\section{Domain Adaptation for Financial Retrieval}
\label{sec:domain_adaptation}

Embedding models trained on general-domain corpora face a systematic mismatch between
the distributional properties of their training data and the properties of financial
documents and queries. Financial reports use a specialized vocabulary in which terms such
as ``EBITDA,'' ``net tangible assets,'' and ``diluted earnings per share'' carry precise
technical meanings that diverge from their general-language senses or have no
general-language counterpart. Queries about financial metrics are typically short,
formally phrased, and saturated with abbreviations, while the relevant passages are long, structurally repetitive, and express the same concepts using standardized financial language of
SEC filings. This query-passage asymmetry is quite prominent in the financial domain, and it is compounded by the specificity of financial
questions. A model that has not seen many examples of matching FY2022 queries to FY2022
evidence, rather than to structurally identical FY2021 or FY2023 passages, has no
learned basis for that distinction.

% The standard approach to bridging this domain gap is contrastive fine-tuning. Given
% domain-specific query-passage pairs, the embedding model is trained to score relevant
% pairs higher than irrelevant ones. The Multiple Negatives Ranking loss is particularly
% effective in this setting. For a minibatch of $B$ positive query-passage pairs
% $\{(q_i, c_i^+)\}_{i=1}^B$, it treats all passages in the batch other than the paired
% one as negatives for each query and minimizes
% \begin{equation}
%   \mathcal{L}_{\text{MNR}} = -\frac{1}{B}\sum_{i=1}^{B}
%   \left[
%     s(q_i, c_i^+)
%     - \log\sum_{j=1}^{B} \exp s(q_i, c_j^+)
%   \right],
%   \label{eq:mnr_background}
% \end{equation}
% where $s(\cdot,\cdot)$ is the cosine similarity from Equation~\eqref{eq:cosine}.
% Equation~\eqref{eq:mnr_background} is the training objective used for the learned page scorer in
% this thesis, as detailed in Chapter~\ref{chap:methodology}. The loss encourages the
% model to rank the gold page above all other pages in the minibatch, including the gold
% pages of other questions in the same batch. Because different questions in the batch come
% from different companies and fiscal years, this provides hard negatives that teach the
% model to distinguish company-specific numerical context from generic financial language.


% The domain adaptation approach adopted in this thesis applies the same principle but at
% a different level of the retrieval hierarchy. Rather than fine-tuning a general
% chunk-level retriever, it fine-tunes a page-level bi-encoder specifically for the task
% of identifying which pages within a financial filing are most likely to contain the
% evidence for a given question. This targeted training is the core methodological
% innovation, and its full design is presented in Chapter~\ref{chap:methodology}.

The standard approach to bridging this domain gap is fine-tuning. Given
domain-specific query-passage pairs, the embedding model is trained to score relevant
pairs higher than irrelevant ones. For bi-encoder fine-tuning, the Multiple Negatives
Ranking (MNR) loss is commonly used. It treats all other passages in the
minibatch as implicit negatives for each query, creating a computationally efficient
approximation to listwise ranking. For cross-encoder rerankers, binary cross-entropy
(BCE) loss over explicit positive and sampled-negative pairs is preferred, as it
calibrates a pointwise relevance score rather than requiring exhaustive listwise
normalisation. Hard negatives, drawn from topically similar but answer-irrelevant
passages rather than at random, play a critical role in shaping what either model
learns, particularly in the financial domain where boilerplate language makes random
negatives trivially easy to dismiss. \citet{kim2024optimizing} demonstrate that fine-tuning BGE-M3 on financial query-passage
pairs improves FinanceBench retrieval substantially relative to the zero-shot model, with
the fine-tuned model showing improved ability to match financial terminology and temporal
expressions. FinSage \citep{wang2024finsage} applies direct preference optimization (DPO) on
financial query-document pairs to train a reranker that learns fine-grained distinctions,
including between passages from the same document that discuss the same topic in different
fiscal years. The specific loss formulations, negative sampling
strategies, and hyperparameter choices for both the bi-encoder page scorer and the
fine-tuned cross-encoder reranker used in this thesis are presented in detail in
Chapter~\ref{chap:methodology}.


\section{Summary}

This chapter has traced the evolution of retrieval-augmented generation from its
foundational contributions through the current state of the field, examined the
information retrieval methods that form the technical core of RAG systems, and situated
financial question answering as a domain that stresses these methods in ways that
general-domain evaluations do not reveal. The survey of benchmarks established
FinanceBench as the primary evaluation standard for this thesis while motivating its
extension to address limitations in document diversity and coverage.

The treatment of hierarchical retrieval surveyed the landscape from general-purpose
recursive indexing systems such as RAPTOR to domain-specific two-stage pipelines for
financial documents. It treated chunking as an integral granularity decision within that
hierarchy rather than a separate preprocessing step, identifying the interaction between
chunk size and retrieval architecture as an understudied variable that this thesis
addresses experimentally. Domain adaptation
established the contrastive fine-tuning approaches that allow general embedding models to
acquire the financial retrieval competence needed to distinguish temporal near-duplicates
and match domain-specific terminology. The closing section on visual document retrieval
introduced the ColPali architecture and recent multi-scale extensions as an emerging
complementary paradigm whose applicability to financial documents is evaluated in the
methodology that follows.

The literature reviewed in this chapter shows that financial RAG poses a retrieval problem that is more difficult than standard end-to-end evaluation. Long, repetitive, and table-heavy filings create challenges that motivate more fine-grained retrieval analysis and specialized retrieval architectures. The next chapters therefore turn to the datasets, indexing decisions, and evaluation setup used to study these challenges.


%%%%%%%%%%%%%%%%%%%%%%%%%%%%%%%%%%%%%%%%%%%%%%%%%%%%%
%%          CHAPTER 2: EXPERIMENTAL SETUP           %%
%%%%%%%%%%%%%%%%%%%%%%%%%%%%%%%%%%%%%%%%%%%%%%%%%%%%%

\chapter{Experimental Setup}
\label{chap:setup}

\section{Problem Formulation}
\label{sec:prob_form}
Let $\mathcal{D} = \{d_1, \ldots, d_{|\mathcal{D}|}\}$ be a corpus of financial filings,
where each document $d \in \mathcal{D}$ consists of a set of pages
$\mathcal{P}(d) = \{p^d_1, \ldots, p^d_{n_d}\}$, each page being a sequence of
extracted text. The full corpus is segmented into a chunk corpus
\begin{equation}
  \mathcal{C} = \{(c,\, d,\, p) \mid c \in \text{chunks}(p),\, p \in \mathcal{P}(d),\,
  d \in \mathcal{D}\},
\end{equation}
where each chunk $c$ carries document identifier $d$ and page index $p$ as metadata.
Given a natural language question $q$, a retriever selects the top-$k$ chunks
\begin{equation}
  R_k(q) = \operatorname{Top\text{-}}k_{(c,d,p)\in\mathcal{C}}\; s(q,\, c)
  = \{(c_i,\, d_i,\, p_i)\}_{i=1}^{k},
  \label{eq:retriever}
\end{equation}
and a generator produces an answer $\hat{y} = G_\theta(q,\, R_k(q))$. FinanceBench
provides a gold filing $d^\star(q)$ and a set of gold page indices $P^\star(q)$ for
each question $q$, enabling evaluation at both the document and page level.

%--------------------------------------------------------------------
\section{Dataset}
\label{sec:dataset}
\label{sec:pipeline}

% ======================================================================
\subsection{FinanceBench}
\label{sec:financebench_dataset}

Our primary evaluation dataset is the open-source subset of FinanceBench
\citep{islam2023financebench}, comprising 150 question-answer pairs over 84 unique
financial PDF documents spanning 32 companies and fiscal years from 2015 to 2024.
Each sample includes a natural-language question, a reference answer, a verbatim
evidence string, and gold metadata identifying the source document and evidence page
number. This page-level annotation provides the ground-truth labels for document
recall, page recall, and oracle retrieval conditions.

% ------------------------------------------------------------------
\subsubsection{Question and Document Types}
\label{sec:financebench_types}

The 150 questions are evenly split across three types with 50 questions each.
\emph{Domain-relevant} questions require locating narrative disclosures in text-heavy
filing sections.
\emph{Metrics-generated} questions require extracting or computing specific financial
indicators from tabular contexts such as income statements or balance sheets.
\emph{Novel-generated} questions test company-specific reasoning beyond simple
extraction, often targeting events disclosed in non-standard filing sections.

The underlying documents span four filing types, summarised in
Table~\ref{tab:filing_types}.

\begin{table}[ht]
  \centering
  \small
  \caption{Filing types in the FinanceBench corpus (84 documents, 150 questions).}
  \label{tab:filing_types}
  \begin{tabularx}{\linewidth}{@{}l c c X@{}}
    \toprule
    \textbf{Filing type} & \textbf{Docs} & \textbf{Questions} & \textbf{Description} \\
    \midrule
    Annual report (10-K)
      & 64 & 112 (74.7\%)
      & Long-form annual disclosure covering risk factors, MD\&A, audited
        financial statements, and footnotes. \\[4pt]
    Quarterly report (10-Q)
      & 8 & 15 (10.0\%)
      & Structurally similar to 10-Ks but shorter and unaudited, covering a
        single fiscal quarter. \\[4pt]
    Earnings call transcript
      & 6 & 14 (9.3\%)
      & Unstructured conversational documents where evidence is spread across
        analyst Q\&A exchanges rather than tables. \\[4pt]
    Current report (8-K)
      & 6 & 9 (6.0\%)
      & Brief event-driven filings where evidence is typically concentrated
        on a single page. \\
    \bottomrule
  \end{tabularx}
\end{table}

% ------------------------------------------------------------------
\subsubsection{Illustrative Examples by Question Type}
\label{sec:financebench_examples}

To make the three question types concrete, Examples~\ref{ex:domainrel}
through~\ref{ex:novelgen} present one representative example from each category.
Each example displays the natural-language question, the reference answer, the verbatim
gold-evidence string, and the source filing metadata. These examples illustrate the
qualitative differences in evidence location, answer type, and required reasoning that
motivate the per-type analysis in Sections~\ref{sec:retrieval_failure}
and~\ref{sec:gen_failure}.

%% Colour palette
\definecolor{domRelColor}{RGB}{31,73,125}
\definecolor{metGenColor}{RGB}{0,102,68}
\definecolor{novGenColor}{RGB}{140,60,10}
\definecolor{errBoxColor}{RGB}{160,30,30}

%% Shared tcolorbox style for dataset examples
\tcbset{
  examplebox/.style={
    enhanced, breakable,
    colframe=#1, colback=#1!5!white,
    fonttitle=\bfseries\small,
    coltitle=white, attach boxed title to top left={yshift=-2mm, xshift=4mm},
    boxed title style={colback=#1, sharp corners},
    sharp corners, boxrule=0.6pt,
    left=6pt, right=6pt, top=10pt, bottom=6pt,
    before upper={\setlength{\parskip}{3pt}}
  }
}

\newcommand{\exrow}[2]{%
  \noindent\textbf{#1}\quad #2\par\medskip
}

\begin{failex}{Domain-Relevant Question}
\label{ex:domainrel}

\noindent\textbf{Question:}\enspace
Among operating, investing, and financing activities, which brought in the most
(or lost the least) cash flow for Nike in FY2023?

\medskip
\noindent\textbf{Expected answer:}\enspace
Cash flow from operations was the highest in FY2023.

\medskip
\noindent\textbf{Evidence:}\enspace
``NIKE, INC.\ CONSOLIDATED STATEMENTS OF CASH FLOWS.\ YEAR ENDED MAY 31 (Dollars in millions).\ 2023 \quad 2022 \quad 2021.\
Cash provided (used) by operations: Net income \$5{,}070 \quad \$6{,}046 \quad \$5{,}727.\
...
\textbf{Cash provided (used) by operations 5{,}841 \quad 5{,}188 \quad 6{,}657.}\
...
\textbf{Cash provided (used) by investing activities 564 \quad (1{,}524) \quad (3{,}800).}\
...
\textbf{Cash provided (used) by financing activities (7{,}447) \quad (4{,}836) \quad (1{,}459).}\
...''

\medskip
\noindent\textbf{Document Name:} NIKE\_2023\_10K \quad \textbf{Filing:} 10-K \quad \textbf{Gold Page:} p.~61
\end{failex}


\begin{failex}{Metrics-Generated Question}
\label{ex:metricsgen}

\noindent\textbf{Question:}\enspace
What is Amcor's year end FY2020 net AR (in USD millions)? Address the question by
adopting the perspective of a financial analyst who can only use the details shown
within the balance sheet.

\medskip
\noindent\textbf{Expected answer:}\enspace
\$1616.00.

\medskip
\noindent\textbf{Evidence:}\enspace
``Amcor plc and Subsidiaries Consolidated Balance Sheet (in millions) As of June 30,
2020 2019 Assets Current assets: Cash and cash equivalents \$ 742.6 \$ 601.6; Trade
receivables, net \textbf{1,615.9} 1,864.3; Inventories...''

\medskip
\noindent\textbf{Document Name:} AMCOR\_2020\_10K \quad \textbf{Filing:} 10-K \quad \textbf{Gold Page:} p.~49
\end{failex}


\begin{failex}{Novel-Generated Question}
\label{ex:novelgen}

\noindent\textbf{Question:}\enspace
By how much did Pepsico increase its unsecured five year revolving credit agreement
on May 26, 2023?

\medskip
\noindent\textbf{Expected answer:}\enspace
\$400,000,000 increase.

\medskip
\noindent\textbf{Evidence:}\enspace
``Effective May 26, 2023, PepsiCo terminated the
\textbf{\$3{,}800{,}000{,}000 five year unsecured revolving credit agreement}
(2022 Five Year Credit Agreement).
...
On May 26, 2023, PepsiCo entered into a new
\textbf{\$4{,}200{,}000{,}000 five year unsecured revolving credit agreement}
(2023 Five Year Credit Agreement) among PepsiCo, as borrower, the lenders party
thereto, and Citibank, N.A., as administrative agent.
...''

\medskip
\noindent\textbf{Document Name:} PEPSICO\_2023\_8K \quad \textbf{Filing:} 8K \quad \textbf{Gold Page:} p.~1
\end{failex}
\vspace{1em}

The PDF documents are downloaded from the public FinanceBench
release.\footnote{\url{https://github.com/patronus-ai/financebench}} Page text is
extracted using PyMuPDF from LangChain \citep{Chase_LangChain_2022}, preserving page
boundaries as metadata. All page numbers correspond directly to the physical page
indices in the PDF documents, validated against the gold annotations.

% ======================================================================
\subsection{Benchmark Extension: FinQA-503}
\label{sec:extended_sources}

FinanceBench's 150-question open-source subset covers a limited range of companies,
years, and document types. To study generalisation beyond this distribution, we
construct an extended evaluation set from LoFIN, the benchmark introduced with HiREC
\citep{choe2025hirec}, which combines the FinQA test set \citep{chen2021finqa} and the
SEC-QA test set \citep{lai2024secqa}.

FinQA \citep{chen2021finqa} contains 8,281 question-answer pairs sourced from 10-K
filings of S\&P~500 companies, emphasising multi-step numerical reasoning over tables
and surrounding text. Each record includes a question, a structured numerical answer,
and \texttt{pre\_text}/\texttt{post\_text} passages bracketing the relevant table
(example in Figure~\ref{fig:finqa_ex} in the Appendix). Because FinQA annotations
reference source passages rather than PDF page numbers, linking them to our corpus
requires an alignment step described in Section~\ref{sec:method_benchmark}.

Of the 1,147 test-set questions, 530 could be matched to downloadable PDFs and 503
were successfully aligned to gold PDF pages via our evidence-matching procedure,
forming the FinQA-503 evaluation subset. SEC-QA questions could not be
included as their records lacked gold evidence text suitable for page alignment.
Unlike FinanceBench, the FinQA-503 corpus is homogeneous in document type, with all
115 documents being annual 10-K reports spanning fiscal years 2008--2019. Answers are
almost exclusively single numerical values, with a mean answer length of approximately
1 token.

% ======================================================================
\subsection{Corpus Statistics}
\label{sec:corpus_stats}

Table~\ref{tab:dataset_stats} reports summary statistics for both datasets.
Token counts use whitespace tokenisation, and chunks are produced with
\texttt{RecursiveCharacterTextSplitter} at a chunk size of 1024 tokens 128 overlap.

\begin{table}[ht]
\centering
\small
\caption{Summary statistics for FinanceBench and FinQA-503.}
\label{tab:dataset_stats}
\begin{tabular}{lcc}
\toprule
\textbf{Metric} & \textbf{FinanceBench} & \textbf{FinQA-503} \\
\midrule
\textbf{Questions} & 150 & 503 \\
\textbf{Documents} & 84 & 115 \\
\textbf{Total pages} & 12,013 & 16,177 \\
\textbf{Pages / document} & $143.0 \pm 98.7$ & $140.7 \pm 59.3$ \\
\textbf{Total chunks} & 15,171 & 22,473 \\
\midrule
\textbf{Query tokens} & $40 \pm 24$ & $36 \pm 9$ \\
\textbf{Answer tokens} & $19.5 \pm 24.1$ & $1.1 \pm 0.7$ \\
\midrule
\textbf{Evidence tokens} & $428 \pm 378$ & $923 \pm 371$ \\
\bottomrule
\end{tabular}
\end{table}

The two corpora are comparable in scale at the document and page level but differ
in question density and answer style. FinanceBench averages 1.8 questions per document
versus 4.5 for FinQA-503, and its answers are free-form natural language while FinQA answers are almost exclusively single numerical values. The evidence token counts are not directly comparable across the two datasets
as they reflect different annotation formats. FinanceBench provides a verbatim snippet selected from the gold page (median 167 tokens, well below the full-page
median of 312 tokens), whereas FinQA-503 supplies the full
\texttt{pre\_text}/\texttt{post\_text} context spanning multiple paragraphs around the
relevant table (mean 923 tokens). Gold evidence is concentrated on a single page in
the large majority of questions in both datasets, with 115 of 150 FinanceBench
questions having a single gold page (35 spanning two or more) and 494 of 503 FinQA-503
questions having exactly one.

% ======================================================================
\subsection{Query-Corpus Lexical Asymmetry}
\label{sec:lexical_asymmetry}

To characterise the vocabulary gap between queries and their source documents, we
compute four complementary statistics reported in Table~\ref{tab:lexical_asymmetry}.
All text is lowercased, tokenised with a regex pattern, stripped of leading and
trailing punctuation. The result is a bag of content tokens per query and per evidence page,
where each query's document is formed by concatenating all its gold evidence pages.

Let $Q$ and $D$ denote the sets of unique word types in the query and document
respectively, and let $p_Q$, $p_D$ denote their normalised unigram frequency
distributions. Term coverage and its complement the out-of-vocabulary (OOV) rate,
\begin{equation}
  \text{Term Coverage} = \frac{|Q \cap D|}{|Q|}, \qquad
  \text{OOV Rate} = 1 - \text{Term Coverage},
\end{equation}
measure the fraction of unique query words present in (or absent from) the document
vocabulary. A high OOV rate indicates that most query terms have no verbatim match in
the source document, which limits the effectiveness of lexical retrieval methods. Jaccard similarity \citep{jaccard1901etude, manning2008introduction}
similarly compares the two type sets but additionally penalises documents with large
vocabularies, since a filing that introduces many terms irrelevant to the query will
push the denominator up even when all query words are covered,
\begin{equation}
  \text{Jaccard}(Q, D) = \frac{|Q \cap D|}{|Q \cup D|}.
\end{equation}
Jensen--Shannon divergence \citep{lin1991divergence} moves beyond set overlap to
capture distributional mismatch, measuring how differently query and document weight
their shared vocabulary,
\begin{equation}
  \text{JSD}(p_Q \| p_D) = \frac{1}{2} D_{\mathrm{KL}}(p_Q \| m)
                          + \frac{1}{2} D_{\mathrm{KL}}(p_D \| m), \quad
  m = \frac{p_Q + p_D}{2},
\end{equation}
where $D_{\mathrm{KL}}$ is the Kullback--Leibler divergence
\citep{kullback1951information}, computed with base-2 logarithms so that JSD is
bounded in $[0,1]$. These measures have been used in prior work to characterise
query--document lexical asymmetry \citep{Kohail2017MatchingRA, luetal2020diverging}.

\begin{table}[ht]
  \centering
  \caption{Query--corpus lexical asymmetry for FinanceBench ($n{=}150$) and FinQA
  ($n{=}503$). OOV rate and JSD indicate mismatch ($\uparrow$ = harder for lexical
  retrieval), Jaccard and term coverage indicate overlap ($\downarrow$ = harder).}
  \label{tab:lexical_asymmetry}
  \resizebox{\linewidth}{!}{%
  \begin{tabular}{ll cccc}
    \toprule
    & & \textbf{OOV Rate} $(\uparrow)$
      & \textbf{Jaccard} $(\downarrow)$
      & \textbf{Term Cov.} $(\downarrow)$
      & \textbf{JSD} $(\uparrow)$ \\
    \midrule
    \multirow{2}{*}{\textbf{FinanceBench}}
      & Mean $\pm$ Std       & $0.68 \pm 0.19$ & $0.03 \pm 0.03$ & $0.32 \pm 0.19$ & $0.94 \pm 0.05$ \\
      & Median {[}Q1, Q3{]}  & $0.70$ [0.57, 0.83] & $0.03$ [0.01, 0.05] & $0.30$ [0.17, 0.43] & $0.95$ [0.91, 0.98] \\
    \midrule
    \multirow{2}{*}{\textbf{FinQA}}
      & Mean $\pm$ Std       & $0.46 \pm 0.19$ & $0.05 \pm 0.04$ & $0.54 \pm 0.19$ & $0.91 \pm 0.05$ \\
      & Median {[}Q1, Q3{]}  & $0.46$ [0.33, 0.60] & $0.04$ [0.03, 0.06] & $0.55$ [0.40, 0.67] & $0.92$ [0.89, 0.95] \\
    \bottomrule
  \end{tabular}}
\end{table}

As shown in Table~\ref{tab:lexical_asymmetry}, FinanceBench exhibits severe lexical
asymmetry, with a median OOV rate of 0.70 indicating that approximately seven in ten
query terms do not appear verbatim in the source filing, and a median JSD of 0.95
confirming that query and document unigram distributions are nearly disjoint. FinQA
shows substantially lower mismatch, with a median OOV rate of 0.46 and a JSD of 0.92,
consistent with its questions being written directly from filing tables and therefore
sharing more verbatim tokens with the source text. This contrast has a direct
consequence for retrieval, as BM25 relies entirely on lexical overlap and collapses
under the high asymmetry of FinanceBench yet remains competitive on FinQA where that
overlap exists. This method inversion is revisited in
Section~\ref{sec:extension_analysis}.

%--------------------------------------------------------------------
\section{PDF Acquisition Pipeline}
\label{sec:pdf_pipeline}

A part of this work involved constructing a reliable PDF corpus from the EDGAR free
online database. FinanceBench already provides PDFs, but the LoFIN/HiREC extension
required recovering filings from EDGAR identifiers since the PDFs were not provided in
the benchmark. The acquisition procedure first normalises company identifiers, resolves
each ticker to its corresponding unique identifier (CIK), queries the EDGAR submissions
metadata for the target accession number, and parses the filing index to select the
primary filing document. Native PDF filings are downloaded directly while HTML filings are
converted into page-preserving PDFs before indexing.

\subsection{HTML-to-PDF Conversion}
\label{sec:html_pdf_conversion}

Many EDGAR filings, particularly older 10-K and 10-Q submissions, are submitted as HTML
documents rather than native PDFs. For these filings, the HTML is rendered to PDF using
a headless browser. Before rendering, the document is made self-contained by embedding
linked images, stylesheets, and background resources. This avoids repeated network calls
during rendering and improves reproducibility across machines.

Before rendering, the conversion procedure detects likely page boundaries in the live
document structure. It uses explicit page containers, horizontal separators,
page-number footers, roman numerals, appendix-style page markers, and related layout
cues. Native page-break rules are then replaced with controlled page breaks, and
oversized sections are scaled to reduce blank or truncated pages. The rendered PDF uses
Letter-sized pages, printed backgrounds, minimal browser margins, and print-media
emulation. Generated PDFs are validated through minimum file-size checks, PDF signature
checks, and page counting when a readable page structure is available. To ensure page
boundaries were respected, a manual check was done on each retrieved PDF.


%--------------------------------------------------------------------
\section{Document Indexing and Retrieval Infrastructure}
\label{sec:indexing}

\subsection{Text Extraction and Page Representation}
\label{text_ext}
For each filing in the corpus, page text is extracted using PyMuPDF via LangChain's
\texttt{PyMuPDFLoader} interface, preserving page boundaries as metadata tuples
$(d, p)$, where $d$ is the document identifier and $p$ is the 0-indexed page number.
All retrieval methods share this same extractor.

PyMuPDF was chosen as the default extractor because it is fast, stable, directly
supported by LangChain's PDF loading interface, and because FinanceBench
\citep{islam2023financebench} adopts it as its reference extractor. A comparison of
extraction methods on the FinanceBench corpus is reported in
Appendix~\ref{app:extraction-comparison}. Overall, the differences across extractors
are small (at most 0.03 at page recall at $k=5$ and 0.04 at page recall at $k=10$),
suggesting that extractor choice is not a primary driver of retrieval performance on
this corpus. Although, future work should more systematically evaluate extraction quality, as the relatively small differences observed here may reflect the limitations of current extraction methods rather than their true equivalence.
%--------------------------------------------------------------------
\subsection{Chunk Construction and Chunking Strategies}
\label{sec:chunking_strategies}

We build the chunk corpus $\mathcal{C}$ by splitting page text into overlapping spans.
Each chunk is represented as a tuple $(c, d, p)$ carrying the chunk text, the source
document identifier, and the page index. The default configuration uses a recursive character splitter with a window of 1024 tokens and an overlap of 128 tokens, following Yepes et al. [2024]. To verify the suitability of this 1024-token configuration for the FinanceBench corpus and to ensure our baseline was optimized, we conducted a preliminary ablation over varying chunk sizes and strategies. The results and configurations of these variants are detailed in Appendix ~\ref{app:chunking_ablation}.

\paragraph{Fixed-size naive chunking.} Text is split at fixed character or token
boundaries without regard to sentence or paragraph structure; included as a baseline to
verify that recursive splitting outperforms naive splitting on financial text.

\paragraph{Recursive character splitting.} The splitter attempts to split on a
priority-ordered list of separators (paragraph breaks, newlines, periods, spaces),
recursively falling back to shorter separators when a chunk exceeds the target size.
This is the default configuration (1024-token window, 128-token overlap).

\paragraph{Semantic chunking.} Consecutive sentences are grouped into chunks by
detecting BGE-M3 embedding similarity drops between adjacent sentences, with a new
chunk started whenever cosine similarity falls below a configurable threshold
(default 0.6).

\paragraph{Parent-child chunking.} Two levels of chunks are constructed: coarse parent
chunks of 1024 tokens serving as generation context, and fine child chunks of 256 tokens
indexed for retrieval and mapped back to their parent at generation time.

\paragraph{Page-aligned chunking.} Chunk boundaries are constrained to fall within a
single page, preventing chunks from spanning adjacent page boundaries and ensuring
unambiguous page metadata for the page scorer pipeline.

\paragraph{Late chunking.} A long document passage is encoded as a single sequence
using BGE-M3, and span-level embeddings are extracted by mean-pooling token-level
representations within each chunk's token span \citep{gunther2024jina}. Page assignment
uses character-level offset tracking and the tokenizer's
\texttt{return\_offsets\_mapping} flag to recover the originating page for each span.

\subsection{Vector Store: ChromaDB}
\label{sec:vectorstore}

All dense retrieval experiments use ChromaDB \citep{chromadb2023}  as the vector store. ChromaDB was selected
over FAISS \citep{douze2025faiss}  primarily for its native support for structured metadata filtering. Each chunk
in the index carries \texttt{doc\_name} and \texttt{page\_num} as metadata fields, and
ChromaDB's \texttt{where} filter allows oracle document and oracle page retrieval
conditions to be implemented as first-class query constraints rather than as
post-retrieval filtering heuristics. 

A single ChromaDB collection is built once per
$(\text{embedding model},\, \text{chunk size},\, \text{chunk overlap})$ configuration
and reused across all retrieval experiments sharing that configuration. This ensures that
comparisons across retrieval methods reflect genuine method differences rather than index
construction variance.

%--------------------------------------------------------------------
\section{Retrieval Systems}
\label{sec:retrieval_systems}

We evaluate multiple retrieval configurations, all operating over the same shared
multi-document index unless otherwise noted. This shared-index design controls for
implementation differences and ensures that metric differences across methods reflect
genuine retrieval quality rather than index construction artefacts.

\subsection{Baseline Retrievers}

Dense retrieval maps both the query and each chunk to a shared vector space via a
bi-encoder and ranks chunks by cosine similarity \citep{karpukhin2020dense}. We
evaluate three dense backbones. BAAI/BGE-M3 \citep{chen2024bge} is a multilingual,
multi-granularity model producing 1024-dimensional vectors with strong performance
across retrieval benchmarks; it serves as the primary dense retriever and as the base
model for the learned page scorer. MPNet is a smaller general-purpose encoder producing
768-dimensional vectors, included as a lighter alternative to separate the effect of
model capacity from pipeline design. \texttt{FinLang/finance-investopedia-embeddings} \citep{finlang_investopedia_2024}
is a domain-specific model pre-trained on Investopedia text (768-dimensional vectors),
included to test whether financial-domain vocabulary exposure alone, without
retrieval-specific contrastive training, is sufficient to outperform general-purpose
encoders on SEC filing retrieval.

For sparse retrieval, BM25 \citep{robertson2009probabilistic} is implemented using the \texttt{rank-bm25} library with
a finance-aware tokenizer that preserves currency symbols, percentage signs, and
formatted numbers. SPLADE \citep{formal2021splade} expands query and document
representations via a masked language model into sparse term-weighted vectors over the
full vocabulary. We use the \texttt{naver/splade-cocondenser-ensembledistil} checkpoint \citep{https://doi.org/10.48550/arxiv.2205.04733},
truncating inputs to 512 tokens and retaining the top 256 weighted terms per
representation for index efficiency.

The hybrid baseline combines BM25 and BGE-M3 rankings using Reciprocal Rank Fusion
(RRF) \citep{cormack2009reciprocal}. For a chunk at rank position $r$, its RRF score is
\begin{equation}
  \text{RRF}(r) = \frac{1}{k_{\text{rrf}} + r},
\end{equation}
with $k_{\text{rrf}} = 60$ following standard practice. Scores from both retrievers are
summed and chunks are re-ranked by the combined score. We evaluate three mixing ratios
(sparse-heavy 25/75, balanced 50/50, and dense-heavy 75/25) to establish the best
hybrid configuration.

\subsection{Query Reformulation}

A deterministic preprocessing step replaces common financial abbreviations
(e.g., \texttt{EPS}, \texttt{EBITDA}, \texttt{FY}) with their full-form equivalents
and normalises fiscal year references to a canonical format (e.g., ``FY2022'' to
``fiscal year 2022''), improving lexical overlap with filing text without requiring a
generative model.

Hypothetical Document Embeddings (HyDE) \citep{gao2023precise} generate a hypothetical
passage that would answer the query and embed that passage rather than the query itself.
We generate one hypothetical passage per query using Qwen-2.5-7B-Instruct at temperature
$0.7$ and embed it with BGE-M3. Multi-HyDE \citep{george2025enhancing} reduces variance
by sampling four independent hypothetical passages per query and averaging their BGE-M3
embeddings into a single query representation; the prompt template is reproduced in
Appendix~\ref{fig:mHprompt}.

\subsection{Reranking}
\label{sec:reranking}

We first implement a zero-shot cross-encoder via \texttt{BAAI/bge-reranker-v2-m3} serving as a reranker. The
dense retriever first produces $N=20$ candidate chunks, the cross-encoder then assigns a
scalar relevance score to each (query, chunk) pair via full cross-attention, and the
top $k=5$ chunks by reranker score are passed to the generator
\citep{nogueira2019document}. BGE-M3 + Multi-HyDE + ReRanker, which applies query
augmentation before retrieval and reranks the resulting candidates, constitutes the
strongest non-learned baseline.

We additionally fine-tune \texttt{BAAI/bge-reranker-v2-m3} on FinQA financial
page-level data to adapt the reranker to the domain. The training procedure and loss
function are described in Section~\ref{sec:method_ft_reranker}; at inference time it
operates identically to the zero-shot variant.

As an alternative second-stage model, we evaluate a late interaction reranker using
ColBERT-style MaxSim scoring \citep{khattab2020colbert}. The query is encoded as a
sequence of token-level embeddings $\{\mathbf{q}_i\}$ and each candidate chunk as a
sequence $\{\mathbf{d}_j\}$; relevance is computed as
\begin{equation}
  s_{\text{MaxSim}}(q, c) = \sum_{i} \max_{j}\; \mathbf{q}_i^\top \mathbf{d}_j.
\end{equation}
We use the \texttt{colbert-ir/colbertv2.0} checkpoint for this configuration.

%--------------------------------------------------------------------
\section{Generation Setup}
\label{sec:generation_setup}

The generator model, prompt template, and decoding settings are held fixed across all
retrieval methods to ensure that observed differences in answer quality reflect
retrieval differences alone. We use Qwen-2.5-7B-Instruct \citep{qwen2025qwen25} as our
primary generator, and use its 3B, 7B, and 14B variants to analyse the impact of
scaling. The prompt follows a system-user template that presents the retrieved context
and question, instructing the model to provide a concise, evidence-grounded answer in
the style of a financial analyst. Greedy decoding is used with a maximum of 256 new
tokens. The model is loaded in 8-bit quantization via \texttt{bitsandbytes} to reduce
GPU memory consumption from approximately 16~GB to 9~GB, enabling experiments on a
single consumer GPU. We additionally evaluate Llama-3.1-8B-Instruct \citep{meta2024llama} as a baseline from
a different generator family. The generator prompt template is given in full in Appendix~\ref{app:gprom}. 

%--------------------------------------------------------------------
\section{Evaluation Metrics}
\label{sec:metrics}

\subsection{Retrieval Metrics}


We begin by defining the core retrieval metrics used throughout this work, document and page recall. These metrics
measure whether a system successfully retrieves the correct document and localizes the
relevant evidence within it, reflecting the two primary stages at which retrieval can fail. 

For each question $q$ with gold filing $d^\star(q)$ and gold pages $P^\star(q)$, let
the top-$k$ retrieved chunks be $R_k(q) = \{(c_i, d_i, p_i)\}_{i=1}^k$. Document
recall at $k$ is the indicator that the gold filing appears among the retrieved chunks'
source documents,
\begin{equation}
  \text{DocRec@}k(q) = \mathbb{I}\!\bigl[d^\star(q) \in \{d_i\}_{i=1}^k\bigr].
  \label{eq:docrecall}
\end{equation}
In the FinanceBench dataset, all questions can be answered by a single document. Page recall at $k$ is the fraction of gold pages covered by chunks retrieved from the
correct filing,
\begin{equation}
  \text{PageRec@}k(q) =
  \frac{\bigl|P^\star(q) \cap \{p_i \mid d_i = d^\star(q)\}_{i=1}^k\bigr|}
       {\,|P^\star(q)|\,}.
  \label{eq:pagerecall}
\end{equation}
There can be more than one needed page to retrieve the full context required to answer a question. We report both metrics averaged over all 150 questions at $k=5$.  We prioritize recall because if the
correct evidence is not retrieved, the generator cannot produce a correct answer. Missing even a single relevant page constitutes a failure mode that cannot be corrected
downstream.

To measure how well the retrieved context captures the gold evidence independently of
page-index accuracy, we also report the maximum chunk-level BLEU \citep{papineni2002bleu}
and ROUGE-L \citep{lin2004rouge} scores between any retrieved chunk and the gold
evidence string $e^\star(q)$,
\begin{equation}
  \text{MaxBLEU@}k(q)    = \max_{i=1}^{k}\;\text{BLEU}(c_i,\; e^\star(q)),
  \quad
  \text{MaxROUGE-L@}k(q) = \max_{i=1}^{k}\;\text{ROUGE-L}(c_i,\; e^\star(q)).
\end{equation}
These metrics complement page-level recall since a method might achieve high page recall yet
still produce low context quality if chunk boundaries do not align well with the
evidence text. Note, context quality can be high even when page recall is
imperfect, in particular if the retrieved chunks happen to overlap the gold evidence despite originating
from a different page index. This situation is noted in
Section~\ref{sec:retrieval_failure}, where a retrieved chunk contains the correct
information but comes from a different page than the annotation.


\subsection{Generative Metrics}

For downstream generation quality, we report ROUGE-L, numeric match accuracy, and BERTScore F1 \citep{zhang2019bertscore}. Numeric match extracts numeric tokens from the generated answer $\hat{y}$ and the reference $y^\star$ by stripping commas and currency symbols, and marks a question as correct if any pair $(n^\star, \hat{n}) \in N^\star \times \hat{N}$ satisfies
\begin{equation}
  \bigl|\hat{n} - n^\star\bigr| \leq \varepsilon_{\text{abs}}
  \;\;\text{or}\;\;
  \frac{\bigl|\hat{n} - n^\star\bigr|}{n^\star} \leq \varepsilon_{\text{rel}},
\end{equation}
with $\varepsilon_{\text{abs}} = \varepsilon_{\text{rel}} = 0.03$, following \citet{islam2023financebench}. BERTScore F1 embeds generated and reference tokens using a frozen contextual encoder and computes precision, recall, and their harmonic mean via pairwise cosine similarity, capturing semantic overlap beyond surface-level n-gram matching.

These metrics are applied selectively by question type. Metrics-generated questions always produce numeric outputs, so they are evaluated using numeric match only, and BERTScore is excluded for these questions because embedding-based similarity is unreliable over short scalar outputs and can penalise numerically correct answers while rewarding numerically incorrect but semantically related ones. BERTScore F1 and ROUGE-L are instead reported for domain-relevant and novel-generated question types, where semantic similarity is a meaningful signal.

\subsubsection{LLM-Based and Human Evaluation}
\label{sec:llm_eval}
LLM-as-a-judge frameworks are not used as primary metrics in this thesis. LLM-based evaluation is stochastic, sensitive to prompt design and model choice, and subject to well-documented systematic biases including positional and verbosity bias \citep{zheng2023judging}. That said, LLM-based scoring is applied to one model configuration as an auxiliary comparison point to situate results relative to the human evaluation and figures reported in prior literature. The primary generative results otherwise rely on the deterministic metrics described above combined with manual human evaluation.

The 150-question FinanceBench subset is small enough to make comprehensive manual evaluation feasible. The evaluator has completed relevant undergraduate coursework in finance and economics, supporting the assessment of financial reasoning and terminology. Grading follows the FinanceBench binary protocol: answers are marked correct or incorrect with no partial credit, and the validity of the reasoning process is prioritised over the final output alone, so a correct answer reached through flawed reasoning or by accident is marked incorrect. A limitation of this setup is that recruiting additional annotators was not feasible within the scope of this project, and the subjectivity of single-annotator evaluation should be kept in mind when interpreting the results.

%--------------------------------------------------------------------
% \section{Chunking Ablation Design}
% \label{sec:chunking_design}

% To isolate the effect of chunking from other design choices, we hold the retrieval
% method fixed at BGE-M3 dense retrieval with $k=5$ and vary chunk size across 256, 512,
% and 1024 tokens with overlaps of 0 and 128 tokens. We additionally compare a
% page-boundary-aligned variant in which chunk boundaries are constrained to fall within
% a single page, preventing chunks from spanning adjacent pages. This yields five
% configurations in total.

% We measure the impact of chunking independently on standalone dense retrieval and on the
% learned page scorer pipeline. The interaction between chunking and the two-stage
% architecture is of particular interest: when the page scorer has already filtered the
% search space to the most relevant pages, the optimal chunk size within that filtered
% pool may differ from the optimal chunk size for global corpus search. We additionally
% compute evidence coverage as a function of chunk size, measuring the fraction of the
% gold evidence string recovered by the top-5 retrieved chunks, to assess how frequently
% gold evidence is fragmented across chunk boundaries. Results are reported in
% Section~\ref{sec:chunking_ablation}.

%--------------------------------------------------------------------
\section{Computational Infrastructure}
\label{sec:infrastructure}

Each job was allocated one NVIDIA RTX 3090 GPU (24~GB VRAM), 8 CPU cores, and between
32~GB and 64~GB of RAM depending on the experiment. Generator models
(Qwen-2.5-7B-Instruct and Llama-3.1-8B-Instruct) are loaded in 8-bit precision using
\texttt{bitsandbytes} \citep{dettmers2022llmint8}, which reduces memory consumption by
approximately half relative to 16-bit loading. Embedding models (BGE-M3, BGE-base, and
the Investopedia encoder) are loaded in full 16-bit precision.


%%%%%%%%%%%%%%%%%%%%%%%%%%%%%%%%%%%%%%%%%%%%%%%%%%%%%
%%              CHAPTER 3: METHODOLOGY              %%
%%%%%%%%%%%%%%%%%%%%%%%%%%%%%%%%%%%%%%%%%%%%%%%%%%%%%

\chapter{Methodology}
\label{chap:methodology}


\section{Overview}
This chapter presents the methodological contributions of this thesis. The first is an oracle-based evaluation framework that decomposes retrieval failures into document-level and page-level components, providing empirical upper bounds at each granularity and a diagnostic framework for attributing performance deficits to their underlying causes. The second is a hierarchical page-then-chunk
retrieval pipeline which leverages a learned page scorer by domain fine-tuning a bi-encoder
based on BAAI/BGE-M3 that treats pages as the main retrieval unit. The third is a domain fine-tuned cross-encoder reranker based on
BGE-reranker-v2-m3, trained on FinQA data and evaluated zero-shot on
FinanceBench, enabling a controlled comparison between zero-shot and domain-adapted
reranking. The fourth evaluates ColPali visual page scoring as an alternative Stage~1
ranker within the same pipeline architecture, allowing a direct comparison between
text-based and image-based page representations on financial documents. The fifth describes the benchmark extension FinQA alignment pipeline used to extend FinanceBench with verified page-level evidence annotations from EDGAR filings, including the fuzzy-matching procedure and
the exclusion criteria applied to unverifiable examples. 
% \vspace{-0.9em}

\section{Oracle-Based Evaluation Framework}
\label{sec:oracle_eval}
A RAG pipeline can fail at two distinct stages. The retriever may fail to surface the
relevant evidence, or the generator may fail to answer correctly even when the relevant
evidence is present.  Measuring end-to-end accuracy alone conflates these two sources
and makes it impossible to identify the primary bottleneck.  To separate them, we
construct oracle retrieval conditions that progressively restrict the candidate pool to
the gold evidence, yielding empirical upper bounds on what retrieval improvements alone
could achieve.

% ─────────────────────────────────────────────────────────────────────────────
\subsection{Oracle Retrieval Conditions}
\label{sec:oracle_conditions}

We denote the full chunk corpus by $\mathcal{C}$, where each chunk is a tuple
$(c,\,d,\,p)$ of text, document identifier, and page number, as defined in
Section~\ref{sec:prob_form}.  For each question $q$, FinanceBench provides a gold
document $d^{\star}$ and a set of gold pages $P^{\star}$.  We define three candidate
sets by successively narrowing the search space.

First is standard retrieval, where the retriever searches the full corpus $\mathcal{C}$.  This is the baseline setting.

Second, Oracle-Document retrieval is defined as a setting in which the candidate set is restricted to chunks from the gold filing.
\begin{equation}
  \mathcal{C}^{\mathrm{doc}}(d^{\star})
    \;=\; \bigl\{\,(c,\,d,\,p)\in\mathcal{C} \;\bigm|\; d = d^{\star}\,\bigr\}.
  \label{eq:oracle_doc}
\end{equation}
This eliminates cross-document retrieval errors entirely.  Any remaining failure is
attributable to within-document page or chunk selection.

Finally, Oracle page retrieval is when the candidate set is further restricted to chunks from the gold pages from the gold document,
\begin{equation}
  \mathcal{C}^{\mathrm{page}}(d^{\star},\,P^{\star})
    \;=\; \bigl\{\,(c,\,d,\,p)\in\mathcal{C} \;\bigm|\; d = d^{\star},\;\;
                  p \in P^{\star}\,\bigr\}.
  \label{eq:oracle_page}
\end{equation}
The gold evidence is guaranteed to appear in the candidate pool, the only remaining
source of retrieval failure is intra-page chunk quality and ranking.

The three candidate sets satisfy strict containment,
\begin{equation}
  \mathcal{C}^{\mathrm{page}}(d^{\star},\,P^{\star})
    \;\subset\;
  \mathcal{C}^{\mathrm{doc}}(d^{\star})
    \;\subset\;
  \mathcal{C}.
  \label{eq:containment}
\end{equation}
Both oracle conditions are implemented via ChromaDB's metadata filtering API, which
restricts nearest-neighbour search to chunks whose document name matches
$d^{\star}$ and, under oracle page conditions, whose page numbers belong to
$P^{\star}$.  This ensures that, within a fixed retrieval model, the only variable across conditions is the candidate set, while the embedding model and similarity function remain identical.

% ─────────────────────────────────────────────────────────────────────────────
\subsection{Oracle Conditions as Empirical Upper Bounds}
\label{sec:oracle_bounds}

Because each oracle candidate set is a subset of the previous one that is guaranteed
to contain the gold evidence, restricting the search space reduces distractors and typically improves retrieval quality by increasing the likelihood that gold evidence appears in the top-$k$ set.  Let $\mathcal{R}_k(q;\,\mathcal{X})$
denote the top-$k$ chunks retrieved when searching candidate set $\mathcal{X}$, and
let $G_\theta(q,\,\mathcal{R}_k)$ denote the generator output given those chunks.  For
any robust retrieval metric $\mathrm{R\text{-}Score}$ and any robust generation metric
$\mathrm{G\text{-}Score}$, the containment in Equation~\eqref{eq:containment} implies,
in expectation over the evaluation set,
\begin{align}
  \mathrm{R\text{-}Score}\!\bigl(\mathcal{R}_k(q;\,\mathcal{C})\bigr)
    &\;\leq\;
  \mathrm{R\text{-}Score}\!\bigl(\mathcal{R}_k(q;\,\mathcal{C}^{\mathrm{doc}})\bigr)
    \;\leq\;
  \mathrm{R\text{-}Score}\!\bigl(\mathcal{R}_k(q;\,\mathcal{C}^{\mathrm{page}})\bigr),
  \label{eq:r_bound}\\[4pt]
  \mathrm{G\text{-}Score}\!\bigl(G_\theta(q,\,\mathcal{R}_k(q;\,\mathcal{C}))\bigr)
    &\;\leq\;
  \mathrm{G\text{-}Score}\!\bigl(G_\theta(q,\,\mathcal{R}_k(q;\,\mathcal{C}^{\mathrm{doc}}))\bigr)
    \;\leq\;
  \mathrm{G\text{-}Score}\!\bigl(G_\theta(q,\,\mathcal{R}_k(q;\,\mathcal{C}^{\mathrm{page}}))\bigr).
  \label{eq:g_bound}
\end{align}
Here, a robust metric refers to one that is monotonic with respect to the inclusion of gold evidence and aligned with task correctness. On the retrieval side, metrics such as DocRec@k and PageRec@k satisfy this property, as they depend only on whether gold documents or pages are retrieved and therefore cannot decrease when distractors are removed. On the generation side, metrics such as human evaluation are more robust, as they directly reflect factual correctness. In contrast, surface-level similarity metrics such as ROUGE-L, BLEU, or embedding-based scores may not be monotonic, as they can be influenced by irrelevant context or lexical variation.

These bounds are defined with respect to a fixed retrieval model and therefore characterize the performance ceiling achievable under that specific embedding and similarity function. When evaluating optimized retrieval models, we compare their performance to these oracle results as empirical reference points under a fixed retriever, rather than as direct upper bounds across different embedding models. The retrieval bounds in Equation~\eqref{eq:r_bound} hold because removing non-gold
documents and pages eliminates distractors that would otherwise displace gold chunks
from the top-$k$ list. The generation bound in Equation~\eqref{eq:g_bound} follows from the fact that the generator $G_\theta$ is fixed across conditions, so improvements in retrieved evidence quality propagate to downstream answer quality.

Equation~\eqref{eq:g_bound} provides the diagnostic core of this framework.  Any gap
between end-to-end performance under standard retrieval and under oracle page retrieval
is attributable purely to retrieval failure, since the generator is identical across
conditions and differs only in the evidence it receives.  These oracle conditions should
be read as empirical diagnostic bounds rather than mathematical guarantees for every
possible model. They restrict the candidate set but do not directly optimize the final
ranking, so a learned reranker or a domain-tuned retriever can exceed an
oracle-document score by using ranking signals that the fixed oracle retriever does not
use. This is why the experiments interpret oracle results as reference points for
localizing failure, not as absolute ceilings on all architectures.

% ─────────────────────────────────────────────────────────────────────────────
\subsection{Retrieval Gap Decomposition}
\label{sec:oracle_gaps}

The oracle conditions enable a decomposition of the performance gap between standard retrieval and oracle-page retrieval into two components.

First, cross-document retrieval gap.
\begin{equation}
  \Delta_{\mathrm{page}}
    \;=\;
  \mathrm{PageRec@}k\!\bigl(\mathcal{C}^{\mathrm{doc}}\bigr)
    \;-\;
  \mathrm{PageRec@}k\!\bigl(\mathcal{C}\bigr).
  \label{eq:delta_page}
\end{equation}
This quantity measures how much page recall improves when document-level retrieval is assumed to be perfect. A large $\Delta_{\mathrm{page}}$ indicates that many relevant pages are missed under standard retrieval because chunks from incorrect documents are ranked ahead of them. In this case, the model is capable of identifying the correct pages once the search is restricted, and the primary bottleneck lies in document-level retrieval. A small $\Delta_{\mathrm{page}}$ indicates that page recall does not improve even when the correct document is given. This suggests that the model fails to rank the answer-bearing pages highly within the document itself, pointing to a within-document retrieval bottleneck. Thus, $\Delta_{\mathrm{page}}$ distinguishes whether errors arise from document discovery or from page-level context retrieval.

Next, the generation gap,
\begin{equation}
  \Delta_{\mathrm{gen}}
    \;=\;
  \mathrm{NumMatch}\!\bigl(\mathcal{C}^{\mathrm{page}}\bigr)
    \;-\;
  \mathrm{NumMatch}\!\bigl(\mathcal{C}^{\mathrm{doc}}\bigr).
  \label{eq:delta_gen}
\end{equation}
This quantity measures how much answer accuracy improves when the correct pages are provided, given that the correct document has already been retrieved. A large $\Delta_{\mathrm{gen}}$ indicates that errors are primarily due to imprecise page-level evidence, suggesting a page-level retrieval bottleneck.

Together, $\Delta_{\mathrm{page}}$ and $\Delta_{\mathrm{gen}}$ decompose the gap between standard retrieval and oracle-page performance into two components. If $\Delta_{\mathrm{page}}$ is large, improvements should focus on cross-document retrieval. If $\Delta_{\mathrm{gen}}$ is large, improvements should focus on finer-grained evidence localization within the document.
\section{Hierarchical Page-then-Chunk Retrieval}
\label{sec:method_page_scorer}

\subsection{Motivation and Architecture}

Standard chunk retrieval searches $\mathcal{C}$ directly, competing against all chunks
in the corpus simultaneously. For long financial filings with repeated templates, dense
tables, and structurally similar sections across years, this means that the retriever
must simultaneously distinguish the correct filing from thousands of similar passages
and locate the correct page within hundreds of pages in that filing. These two tasks
have different natural difficulty levels, and conflating them in a single retrieval step
means that the method optimized for document-level discrimination may not be optimal for
page-level localization within the correct document.

The proposed hierarchical pipeline separates these tasks explicitly. In Stage~1, a
learned page scorer ranks all pages in the corpus and selects the top-$P$ pages. In
Stage~2, standard dense chunk retrieval is applied, but restricted to chunks from those
$P$ pages. This design reduces the effective search space for chunk retrieval from
$|\mathcal{C}|$ to approximately $P \cdot \bar{n}$ chunks, where $\bar{n}$ is the
average number of chunks per page. With $P=20$, this represents a reduction of roughly
two orders of magnitude for a corpus of the FinanceBench scale. The filtering also
concentrates chunk retrieval on the most semantically relevant regions of the corpus,
reducing noise from unrelated sections while preserving fine-grained selection within the
selected pages.

\subsection{The Learned Page Scorer}

For each filing $d$ and page $p$, we extract the normalized page text
$\tilde{t}_{d,p} = \mathcal{N}(t_{d,p})$ as described in Section~\ref{sec:indexing} and
encode it with a fine-tuned bi-encoder $E_\theta^{\text{page}}$. The page relevance
score for query $q$ is the cosine similarity between the query embedding and the page
embedding,
\begin{equation}
  s_{\text{page}}(q,\, d,\, p) =
  \cos\!\bigl(E_\theta^{\text{page}}(q),\;\; E_\theta^{\text{page}}(\tilde{t}_{d,p})\bigr).
  \label{eq:page_score}
\end{equation}
The same encoder is used for both the query and the page text; the model learns to
represent queries and pages in a shared space where relevant pairs are close. The top-$P$
pages selected by Stage~1 are
\begin{equation}
  \mathcal{P}_P(q) = \operatorname{Top\text{-}}P_{(d,p)\in\mathcal{P}} \; s_{\text{page}}(q,\, d,\, p),
\end{equation}
and the filtered chunk corpus is
\begin{equation}
  \mathcal{C}_{\text{filt}}(q) = \{(c,\, d,\, p) \in \mathcal{C} \mid (d,\, p) \in \mathcal{P}_P(q)\}.
\end{equation}
Stage~2 then runs the standard dense chunk retriever over $\mathcal{C}_{\text{filt}}(q)$
and returns the top-$k$ chunks for generation. The full routine is given in
Algorithm~\ref{alg:page_chunk}.

\begin{algorithm}[ht]
\caption{Page-then-chunk retrieval with the learned page scorer.}
\label{alg:page_chunk}
\begin{algorithmic}[1]
\REQUIRE Query $q$; fine-tuned page encoder $E_\theta^{\text{page}}$; chunk encoder
$E^{\text{chunk}}$; page index $\mathcal{I}_{\text{page}}$ storing $(d,p,\mathbf{v}^{\text{page}}_{d,p})$
with $\mathbf{v}^{\text{page}}_{d,p} = E_\theta^{\text{page}}(\tilde{t}_{d,p})$ precomputed
offline; chunk index $\mathcal{I}_{\text{chunk}}$ storing $(c, \mathbf{v}^{\text{chunk}}_c, d_c, p_c)$
with $\mathbf{v}^{\text{chunk}}_c = E^{\text{chunk}}(c)$ precomputed offline; number of pages $P$;
number of chunks $k$.
\STATE \textbf{Stage 1 (page retrieval):}
\STATE $\mathbf{e}^{\text{page}}_q \leftarrow E_\theta^{\text{page}}(q)$
\STATE $\mathcal{P}_P(q) \leftarrow \text{SearchCos}(\mathcal{I}_{\text{page}},\; \mathbf{e}^{\text{page}}_q,\; P)$
\COMMENT{nearest-neighbour search over stored $\mathbf{v}^{\text{page}}_{d,p}$}
\STATE \textbf{Stage 2 (chunk retrieval on the filtered search space):}
\STATE $\mathbf{e}^{\text{chunk}}_q \leftarrow E^{\text{chunk}}(q)$
\STATE $\mathcal{C}_{\text{filt}} \leftarrow \{c \in \mathcal{I}_{\text{chunk}} \mid (d_c,\, p_c) \in \mathcal{P}_P(q)\}$
\STATE $R_k(q) \leftarrow \text{SearchCos}(\mathcal{C}_{\text{filt}},\; \mathbf{e}^{\text{chunk}}_q,\; k)$
\COMMENT{search over stored $\mathbf{v}^{\text{chunk}}_c$ in the filtered pool}
\RETURN $R_k(q)$
\end{algorithmic}
\end{algorithm}

The ablation in Appendix~\ref{app:page_ablation} shows that $P=20$ achieves the best
balance between recall and noise reduction, with little gain from larger values of $P$.

\subsection{Training Procedure}
\label{sec:training_procedure}

The page scorer is trained by fine-tuning BAAI/BGE-M3 as a bi-encoder for page-level
relevance on financial filings. The model is initialized from the pretrained BGE-M3
checkpoint and optimized using the Multiple Negatives Ranking (MNR) loss
\citep{henderson2017efficient}. For a minibatch of $B$ positive (query, page) pairs
$\{(q_i, \tilde{t}_{d_i^\star, p_i^+})\}_{i=1}^B$, the loss is
\begin{equation}
  \mathcal{L}_{\text{MNR}} = -\frac{1}{B} \sum_{i=1}^{B}
  \left[
    s_{\text{page}}(q_i, d_i^\star, p_i^+)
    - \log \sum_{j=1}^{B} \exp\, s_{\text{page}}(q_i, d_j^\star, p_j^+)
  \right].
  \label{eq:mnr}
\end{equation}
For each query $q_i$, its paired gold page is the positive target and all other gold
pages in the minibatch serve as in-batch negatives. Because questions in the same
minibatch come from different companies and fiscal years, the negatives are
\emph{semantically hard}: they describe the same types of financial information (income
statements, balance sheets) but from different entities, forcing the model to learn
company- and year-specific numerical grounding rather than relying on generic financial
vocabulary.

Training data is constructed from the 150 FinanceBench annotations. For each question,
we pair it with every annotated gold page in the gold filing, yielding 170 positive pairs
(some questions have multiple gold pages). To prevent data leakage, questions are
partitioned by source document so that no document appears in both training and evaluation
folds. We perform 5-fold cross-validation, training on 80\% of documents and evaluating
on the remaining 20\% in each fold.

Optimization uses AdamW with a learning rate of $2 \times 10^{-5}$, linear warmup over
the first 10\% of training steps, and cosine decay thereafter. Each fold is trained for
15 epochs with a batch size of 16, maximum sequence length of 512 tokens, and early
stopping with patience of 3 evaluations (evaluated every 3 epochs).

% \subsection{Query Augmentation}

% To close the distributional gap between short financial queries and long structured page
% text, each query is augmented with the company name and fiscal year extracted from the
% gold document identifier before training and inference. For example, a query about
% ``net accounts receivable'' from \texttt{ZBH\_2008\_10K} becomes ``ZBH net accounts
% receivable FY2008.'' This augmentation is applied deterministically and ensures that the
% model can exploit entity-specific context during page scoring, which is particularly
% helpful for metrics questions that reference company-specific line items without naming
% the company.

\subsection{Key Design Choices}

Several design choices distinguish the learned page scorer from prior hierarchical
retrieval methods.

The choice of \emph{pages} as the retrieval unit, rather than arbitrary passages or
sections, is motivated by their unambiguous boundaries, their correspondence to coherent
presentation units in financial filings, and their direct annotation in FinanceBench,
which provides supervision unavailable for sub-page passage boundaries.

The choice to \emph{fine-tune a separate page encoder} rather than reusing the zero-shot
chunk-level retriever explicitly teaches the model the asymmetry between short analytical
queries and long, structured financial page text.

The use of \emph{document-level cross-validation splits} prevents the model from
memorizing filing-specific patterns and provides a cleaner estimate of generalization.
Using question-level splits instead would allow questions about the same filing to appear
in both training and test folds, inflating reported performance.

Finally, using \emph{in-batch negatives} from the gold pages of other questions
provides hard negatives without requiring explicit offline negative mining, which would
require a separate retrieval pass and introduces risk of false negatives (pages that are
actually relevant but not annotated as gold).

%--------------------------------------------------------------------
\section{Fine-Tuned Cross-Encoder Reranker}
\label{sec:method_ft_reranker}

\subsection{Motivation}

The zero-shot cross-encoder reranker (Section~\ref{sec:reranking}) applies a
general-purpose relevance model to financial passages without any domain-specific
adaptation.  While this is a strong baseline, financial filings have a distinctive
vocabulary and structure that a general model may not score optimally.  To quantify how
much domain fine-tuning helps on the reranker, we train a second cross-encoder on FinQA financial
page data and evaluate it on the held-out FinanceBench benchmark.

\subsection{Architecture}

The fine-tuned reranker is initialized from \texttt{BAAI/bge-reranker-v2-m3}, a
multilingual cross-encoder pre-trained for passage reranking.  The model takes as input
the concatenation $[q;\, \texttt{[SEP]};\, d]$ of a query $q$ and a passage $d$, passes
the sequence through the transformer stack, and applies a linear projection to the
\texttt{[CLS]} representation to produce a scalar relevance score
\begin{equation}
  s_\theta(q, d) = \mathbf{w}^\top \text{CLS}\bigl(E_\theta([q;\, \texttt{[SEP]};\, d])\bigr).
\end{equation}
The maximum input length is capped at 512 tokens, which accommodates both the short
financial query and a substantial portion of the page text.

\subsection{Training Data}

Training data is constructed from the FinQA training set \citep{chen2021finqa}, which
contains 6,251 financial question-answer pairs annotated against pages from SEC annual
report filings.  For each FinQA training sample, we build the gold text
by concatenating the pre-text, formatted table, and post-text fields (Appendix ~\ref{fig:finqa_ex}) corresponding to the
annotated gold evidence:
\begin{align}
  t_{\text{page}} &= \text{pre\_text} \;\|\; \text{table} \;\|\; \text{post\_text}.
\end{align}
This produces one positive pair $(q, t_{\text{page}}, 1)$ per sample.

For hard negatives, we sample $m=3$ pages from other questions' gold pages,
specifically from questions whose source document differs from the current sample.  This
creates inter-document negatives, financially plausible pages (income statements, balance
sheets from other companies or years) that are wrong answers to the current question.  We
deliberately avoid intra-document negatives, passages retrieved from within the same
page as the gold evidence, because our retrieval task is inter-page (ranking different
pages against each other), and training on within-page negatives would teach the wrong
discrimination skill.  Each training batch therefore contains one positive and three
negative pairs per question, giving roughly $6{,}251 \times 4 = 25{,}004$ training
examples.

\subsection{Loss Function}

The reranker is trained with binary cross-entropy loss applied to the sigmoid-transformed
scalar score:
\begin{equation}
  \mathcal{L}_{\text{BCE}} = -\frac{1}{N}\sum_{i=1}^{N}
  \Bigl[y_i \log \sigma(s_\theta(q_i, d_i))
       + (1 - y_i) \log \bigl(1 - \sigma(s_\theta(q_i, d_i))\bigr)\Bigr],
  \label{eq:bce}
\end{equation}
where $y_i \in \{0, 1\}$ is the binary relevance label and $\sigma$ is the sigmoid
function.  This loss directly optimises the model to assign high scores to positive
(query, page) pairs and low scores to negative pairs, without requiring the
listwise normalisation required by MNR loss.  Binary cross-entropy is appropriate here
because the negatives are sampled rather than exhaustive, so the model learns a calibrated
relevance signal rather than a relative rank within a fixed candidate list.

\subsection{Training Procedure}

Training uses the AdamW optimiser with a learning rate of $2 \times 10^{-5}$, a linear
warmup over the first 10\% of steps, and no additional decay schedule.  We train for
3 epochs with a batch size of 32 pairs; at this batch size and a sequence length of
512 tokens under FP16 mixed precision, memory consumption is within the 46\,GB VRAM
budget of the RALI L40S GPU.  The best checkpoint is selected by MAP on the FinQA
development set, evaluated using the \texttt{CERerankingEvaluator} from
\texttt{sentence-transformers}. 

Table~\ref{tab:ft_reranker_training} summarises the hyperparameters used for both the
bi-encoder page scorer (Section~\ref{sec:training_procedure}) and the cross-encoder
reranker, enabling direct comparison of the two fine-tuning regimes.

\begin{table}[ht]
\centering
\caption{Hyperparameters for the bi-encoder page scorer and the fine-tuned cross-encoder
  reranker.}
\label{tab:ft_reranker_training}
\small
\begin{tabular}{lcc}
\toprule
\textbf{Hyperparameter} & \textbf{Bi-Encoder (Page Scorer)} & \textbf{Cross-Encoder (Reranker)} \\
\midrule
Base model          & BAAI/bge-m3              & BAAI/bge-reranker-v2-m3 \\
Training data       & FinanceBench (5-fold CV) & FinQA train (6,251 samples) \\
Loss function       & MNR (in-batch negatives) & BCE (sampled negatives) \\
Negatives per query & In-batch (15 per step)   & 3 inter-document \\
Learning rate       & $2 \times 10^{-5}$       & $2 \times 10^{-5}$ \\
Batch size          & 16                       & 32 \\
Epochs              & 15                       & 3 \\
Warmup              & 10\% of steps            & 10\% of steps \\
Max sequence length & 512 tokens               & 512 tokens \\
GPU                 & RTX 3090                 & L40S \\
\bottomrule
\end{tabular}
\end{table}

At inference time, the fine-tuned cross-encoder reranker operates identically to its
zero-shot counterpart. BGE-M3 dense retrieval first produces $N=20$ candidate chunks, the
reranker scores each (query, chunk) pair in batches of 64, and the top-5 chunks by
reranker score are returned to the generator.

%--------------------------------------------------------------------
\section{Visual Page Scoring with ColPali}
\label{sec:method_colpali}

As discussed in Section~\ref{sec:vlm_retrieval}, visual retrieval models offer an
alternative approach to text extraction entirely.
Rather than representing pages as sequences of extracted tokens, a vision-language model
encodes each page as a sequence of patch embeddings derived from the rendered page image.
This approach preserves the spatial layout of tables and multi-column structures that
text extraction frequently distorts, and may capture visual cues such as font weight
and column alignment that are semantically significant but invisible to text-based
encoders.

To evaluate whether visual page scoring can serve as a viable or complementary
alternative, we replace the fine-tuned text-based bi-encoder in Stage~1 from the learned page scorer \ref{sec:method_page_scorer} with ColPali
\citep{faysse2024colpali}. Under this variant, each page is represented offline as a set
of patch embeddings $\{\mathbf{e}^p_j\}$ produced by the PaLIGemma vision encoder \citep{beyer2024paligemma}. At
query time, the relevance score for page $p$ is computed using the ColBERT-style MaxSim
operator,
\begin{equation}
  s_{\text{ColPali}}(q, p) = \sum_{i} \max_{j}\; \mathbf{q}_i^\top \mathbf{e}^p_j,
  \label{eq:colpali}
\end{equation}
where $\{\mathbf{q}_i\}$ are the token-level query embeddings. The top-$P$ pages are
selected by this visual score, and Stage~2 chunk retrieval proceeds identically over the
filtered page set using the same BGE-M3 chunk encoder. This design isolates the
contribution of the Stage~1 scoring function and all downstream behaviour is held constant, only the page ranking criterion changes between the text-based and visual variants.

One practical difference between the two variants is indexing cost. Text-based page
embeddings are single 1024-dimensional vectors per page, while ColPali stores a full
sequence of patch embeddings per page (typically several hundred vectors), increasing
the Stage~1 index footprint by roughly two orders of magnitude and requiring a MaxSim
scan at retrieval time rather than a standard nearest-neighbour lookup. For the
FinanceBench corpus scale this remains tractable, but the difference is relevant when
considering deployment at scale.

%--------------------------------------------------------------------

%--------------------------------------------------------------------
\section{Benchmark Extension: FinQA Alignment Pipeline}
\label{sec:method_benchmark}

The benchmark extension used in this thesis is based on the FinQA labels provided by
LoFIN only \citep{choe2025hirec}. While HiREC releases the broader LoFIN benchmark, the
local SEC-QA records did not contain gold evidence text or page-level annotations that
could be reliably mapped to PDF pages. SEC-QA was therefore excluded from the page-level
benchmark, as downloadable PDFs alone are not sufficient for PageRec@k evaluation, which
requires verified gold pages.

Each LoFIN FinQA example contains a structured identifier encoding the company, fiscal
year, and source document. This identifier is matched to the corresponding original FinQA
record, from which the surrounding evidence text is extracted. To recover page-level
supervision, we construct an evidence anchor by concatenating a fixed number of sentences
from the end of the \texttt{pre\_text} field and the beginning of the \texttt{post\_text}
field, capturing the local context in which the answer appears.

For each candidate filing, page text is extracted from the corresponding PDF, as in Section \ref{text_ext}. Page-level alignment is then performed using token-set fuzzy
matching implemented via \texttt{rapidfuzz}~\citep{bachmann2023rapidfuzz}. Specifically,
we compute a similarity score between the anchor and each page using the
\texttt{token\_set\_ratio} function, which is robust to partial containment of the anchor
within a longer page. This is well-suited to our setting, as the anchor typically represents
a short excerpt of the full page content. The highest-scoring page is selected and accepted
as the gold page only if its similarity exceeds a threshold $\tau = 0.70$, otherwise the
example is excluded from page-level supervision.

From the LoFIN FinQA test split (1,102 rows), 530 examples were successfully matched to
FinQA records and had candidate PDF mappings identified via identifier matching and PDF
lookup. Of these, 503 (94.9\%) produced PDF-verified page evidence. The remaining 27
aligned examples were excluded from page-level supervision since 13 fell below the fuzzy-matching
threshold, and 14 did not have a valid PDF available in the local corpus at the time of
processing.


\chapter{Results and Analysis}
\label{chap:results}

This chapter presents and analyses all experimental results in four parts.
Section~\ref{sec:oracle_results} establishes oracle-based empirical upper bounds that
serve as reference points throughout.  Sections~\ref{sec:retrieval_results}
and~\ref{sec:gen_results} evaluate retrieval and generation respectively, each following
the same three-level structure: aggregate performance, stratified breakdowns by question
and document type, and analysis.  Finally, Section~\ref{sec:crossdomain} reports generalisation to the benchmark extension.  

% ══════════════════════════════════════════════════════════
\section{Oracle Analysis}
\label{sec:oracle_results}

Before evaluating any optimized retrieval system, we apply the oracle-based evaluation
framework introduced in Section~\ref{sec:oracle_eval} to quantify where retrieval failures
actually occur. By progressively restricting the candidate pool from the full
corpus~$\mathcal{C}$ to the gold filing~$\mathcal{C}^{\mathrm{doc}}$ and then to the
gold pages~$\mathcal{C}^{\mathrm{page}}$, while holding the embedding model and
generator fixed, we isolate cross-document errors from within-document localisation
errors. Table~\ref{tab:oracle} reports the results.

\begin{table}[ht]
\centering
\caption{Oracle retrieval results at $k=5$ using BGE-M3. Numeric match is drawn
  from the downstream generation evaluation in Table~\ref{tab:gen_overall}.}
\label{tab:oracle}

\resizebox{\linewidth}{!}{
\begin{tabular}{lccccc}
\toprule
\textbf{Setting} & \textbf{DocRec@5} & \textbf{PageRec@5}
  & \textbf{BLEU@5} & \textbf{ROUGE-L@5} & \textbf{Num.~Match} \\
\midrule
Dense BGE-M3 (standard) & 0.94 & 0.35 & 0.25 & 0.35 & 0.32 \\
Oracle-Document         & 1.00 & 0.60 & 0.25 & 0.42 & 0.44 \\
Oracle-Page             & 1.00 & 1.00 & 0.40 & 0.59 & 0.70 \\
\bottomrule
\end{tabular}
}
\end{table}

Standard retrieval achieves high document recall ($\text{DocRec@5} = 0.94$) but low page recall ($\text{PageRec@5} = 0.35$). Supplying the gold document
raises page recall to $0.60$, with a $\Delta_{\mathrm{page}} = 0.25$, meaning that chunks from incorrect filings competing for the top-$k$
slots, or cross-document interference, accounts for 0.25 points of lost page recall. However, even within the
correct filing, page recall reaches only 0.60 rather than 1.00, leaving a residual
within-document gap of $0.40$ points. This is the more consequential bottleneck, the
retriever consistently fails to retrieve the answer-bearing pages even once the search is
restricted to the right document.

This page localisation failure carries a substantial downstream cost. $\Delta_{\mathrm{gen}}$
measures how much numeric answer accuracy improves when the correct pages are provided rather
than merely the correct filing, with all else held equal. Moving from Oracle-Document
to Oracle-Page raises numeric match from $0.44$ to $0.70$, a gain of $0.26$ points
that reflects the quality of the evidence passed to the generator, not any change in
the generator itself. The implication is that imprecise page localisation, not
generation capacity, is the primary obstacle to correct answers.

Together, these results confirm that the primary bottleneck is within-document page
localisation rather than document discovery, motivating the approaches developed in the methodology. Throughout the chapter,
Oracle-Document serves as the reference bound for gains attributable to perfect document
selection, and Oracle-Page bounds the gains achievable from perfect page localisation
under a fixed generator.
% ══════════════════════════════════════════════════════════
\section{Retrieval Evaluation}
\label{sec:retrieval_results}
% ══════════════════════════════════════════════════════════

This section evaluates all retrieval configurations at three levels of granularity.
Aggregated results over the full 150-question benchmark establish the overall ranking
of methods and identify the dominant failure modes. A budget analysis examines how recall
evolves as the retrieval budget $k$ increases, separating structural retrieval failures
from cases where a larger candidate set would suffice. Then, stratified breakdowns by question
type and document type decompose these aggregate figures to show where each
method succeeds and where it fails, and why. The section closes with a
qualitative analysis that grounds the quantitative patterns in concrete
examples drawn from question types.

\subsection{Aggregated Retrieval Performance}
\label{sec:agg_retrieval}

Table~\ref{tab:retrieval} reports DocRec@5, PageRec@5, BLEU@5, and ROUGE-L@5 across all
configurations evaluated on the full 150 FinanceBench questions at $k=5$.

\begin{table}[ht]
\centering
\caption{Aggregated retrieval and context-quality performance over 150 FinanceBench
  questions at $k=5$.  BLEU and ROUGE-L measure context quality (retrieved chunks
  vs.\ gold evidence string).  Bold marks the best non-oracle result per column.}
\label{tab:retrieval}
\scriptsize
\renewcommand{\arraystretch}{0.86}
\setlength{\tabcolsep}{3pt}
\begin{tabularx}{\textwidth}{Xcccc}
\toprule
\textbf{Method} & \textbf{DocRec@5} & \textbf{PageRec@5}
  & \textbf{BLEU@5} & \textbf{ROUGE-L@5} \\
\midrule
\multicolumn{5}{l}{\textit{Sparse}} \\
BM25                                  & 0.55 & 0.11 & 0.07 & 0.15 \\
SPLADE                                & 0.81 & 0.31 & 0.24 & 0.36 \\
\midrule
\multicolumn{5}{l}{\textit{Dense}} \\
MPNet                                 & 0.77 & 0.33 & 0.24 & 0.34 \\
Finance-Investopedia                  & 0.52 & 0.10 & 0.09 & 0.19 \\
Dense BGE-M3                          & 0.94 & 0.35 & 0.25 & 0.35 \\
BGE-M3 + RR                           & 0.95 & 0.44 & 0.29 & 0.40 \\
\midrule
\multicolumn{5}{l}{\textit{Hybrid}} \\
RRF 25/75 (sparse-heavy)              & 0.62 & 0.11 & 0.07 & 0.15 \\
RRF 50/50                             & 0.87 & 0.27 & 0.19 & 0.28 \\
RRF 75/25 (dense-heavy)               & 0.93 & 0.32 & 0.22 & 0.32 \\
\midrule
\multicolumn{5}{l}{\textit{Hierarchical and query reformulation}} \\
Parent-Child                          & 0.91 & 0.32 & 0.23 & 0.34 \\
Query Expansion                       & 0.93 & 0.38 & 0.26 & 0.37 \\
BGE-M3 + HyDE                         & 0.91 & 0.40 & 0.26 & 0.37 \\
BGE-M3 + Multi-HyDE                   & 0.89 & 0.44 & 0.27 & 0.39 \\
BGE-M3 + Multi-HyDE + RR              & 0.93 & 0.48 & 0.32 & 0.43 \\
\midrule
\multicolumn{5}{l}{\textit{Proposed}} \\
Page-then-chunk (in-domain)           & 0.95 & 0.55 & 0.33 & 0.46 \\
ColQwen2 Page-then-chunk (VLM)           & 0.65 & 0.22 & 0.15 & 0.24 \\
BGE-M3 + Multi-HyDE + FT-RR          & 0.95 & 0.56 & 0.33  & 0.44 \\
\midrule
Oracle document                       & 1.00 & 0.60 & 0.25 & 0.42 \\
Oracle page                           & 1.00 & 1.00 & 0.40 & 0.59 \\
\bottomrule
\end{tabularx}
\end{table}

As a complementary check on the aggregate ranking, the FinanceBench baseline
retrieval outputs were also compared using paired bootstrap significance tests on
PageRec@5, with 10{,}000 resamples and Holm--Bonferroni correction for multiple
comparisons. This analysis is useful for determining whether observed differences
are supported consistently at the per-question level, rather than being driven by
a small number of examples. At the same time, it should be interpreted as a
limited robustness check. Because the bootstrap operates on a single realised run
for each method, it quantifies uncertainty due to finite sample composition, but
not variability induced by random initialization, training order, or other
stochastic effects in learned retrieval pipelines. In other words, it is not a
substitute for repeated evaluation across multiple random seeds.

For that reason, the bootstrap results reported for this section should be read as
supporting evidence rather than as a definitive measure of experimental
robustness. In particular, small differences between high-performing methods
should be interpreted cautiously, especially when they are not clearly separated
after multiple-comparison correction. A stronger validation protocol would repeat
training and evaluation over several seeds and then combine both seed-level and
question-level uncertainty. Such a study was not carried out here due to time and
compute constraints, and remains an important direction for future work.

A clear structural pattern emerges across all methods, in particular that document-level retrieval is largely strong, while page-level localisation remains the dominant bottleneck. High-performing dense retrievers consistently achieve DocRec@5 above 0.90, yet PageRec@5 remains substantially lower, even for the strongest non-oracle systems. This persistent gap reflects a fundamental limitation of similarity-based retrieval in long, internally repetitive documents. Financial filings contain many near-duplicate sections (e.g., repeated tables across years), so identifying the correct document is relatively easy, but distinguishing the specific evidence-bearing page requires fine-grained reasoning beyond coarse semantic similarity. 

The relative performance of sparse, dense, and hybrid methods is best understood through their treatment of query–document mismatch. Sparse retrieval, represented by BM25, performs poorly across all metrics, reflecting its reliance on exact lexical overlap. Queries in this benchmark are short and analytical, while evidence pages are expressed in formal, domain-specific language and often embedded within tables, limiting the effectiveness of term matching. SPLADE improves substantially over BM25 by learning vocabulary expansion, but remains below dense retrieval methods, as its representations are still constrained to the vocabulary space and cannot fully capture semantic equivalence \citep{formal2021splade}.

Dense retrieval models overcome this limitation by embedding queries and documents into a shared semantic space. General-purpose encoders such as MPNet already outperform sparse methods, while BGE-M3 achieves near-saturated document recall, indicating that semantic matching is sufficient for identifying the correct filing. However, the gap between BGE-M3 and its reranked variant highlights the limitation of independent encoding. Because bi-encoders compute representations separately, they cannot model fine-grained interactions between query and passage. Adding a cross-encoder reranker (BGE-M3 + RR) improves PageRec@5 substantially, consistent with prior findings that joint encoding enables more precise
relevance estimation in reranking settings \citep{nogueira2019document}.

Hybrid retrieval does not consistently improve performance on this benchmark. While hybrid methods are often effective when sparse and dense signals are complementary, the results here suggest that the lexical signal contributes limited additional benefit. As the sparse weight increases, performance tends to degrade, and even dense-heavy fusion does not consistently surpass the dense baseline. One plausible explanation is that, in this setting, queries are short and do not contain distinctive lexical anchors (e.g., exact table headers or entity names), reducing the usefulness of term-based matching. At the same time, the structured and repetitive nature of financial filings means that semantic similarity already captures much of the signal required for document-level retrieval. These findings suggest that, on this benchmark, dense representations provide the dominant retrieval signal, while the contribution of sparse features appears more limited.


The largest gains among non-supervised approaches arise from query reformulation.
Methods such as HyDE and Multi-HyDE consistently improve PageRec@5 relative to the dense baseline, indicating that transforming the query into a document-like representation reduces the mismatch between query and evidence distributions. This behaviour is consistent with the formulation of HyDE, which uses generated pseudo-documents to better align embedding spaces \citep{gao2023precise}. The additional improvement from Multi-HyDE suggests that averaging over multiple generations stabilises the representation and reduces variance introduced by any single synthetic query. When combined with reranking, this produces the strongest off-the-shelf baseline, highlighting that query-side improvements and ranking refinements
address complementary aspects of the retrieval problem.

The proposed methods further demonstrate that explicitly modelling the structure of the retrieval task yields substantial gains. The page-then-chunk pipeline achieves a large improvement in PageRec@5, closing most of the gap to the oracle document condition. This improvement arises from search space restriction, by first selecting a small set of high-probability pages, the system reduces noise and increases the density of relevant candidates before performing fine-grained retrieval. This coarse-to-fine strategy is consistent with prior work on hierarchical retrieval, which shows that multi-stage filtering improves performance in long-document settings \citep{gao2023retrieval}. The fine-tuned reranker (BGE-M3 + Multi-HyDE + FT-RR) achieves comparable performance through a different mechanism, learning domain-specific relevance signals that improve ranking without explicitly restricting the search space.

The visual variant of the proposed pipeline, based on ColQwen2, underperforms its
text-based counterpart across all metrics. While it retains moderate document-level recall, its PageRec@5 and context overlap scores remain substantially lower. This reflects a mismatch between the inductive biases of vision-language models and the structure of financial documents. Visual encoders are designed to capture layout and high-level semantic patterns, but in this domain the discriminative signal lies in precise numerical content and subtle textual differences between structurally similar tables. As a result, visual similarity tends to emphasise layout correspondence rather than factual correctness, leading to confusion between near-identical pages.

Finally, context quality (BLEU@5 and ROUGE-L@5) highlights the role of ranking beyond
document-level recall. Both proposed methods outperform all baselines and match or exceed
the oracle document condition, indicating that optimising the ordering of retrieved content
is as important as identifying the correct document. The page-then-chunk model achieves the highest ROUGE-L@5, benefiting from in-domain
training on FinanceBench and restriction to a small set of high-probability pages.
However, this approach relies on supervision from the target dataset and may not generalise
as effectively. The ColQwen2 variant further illustrates the limitations of alternative
representations, with substantially lower page recall and context overlap. In contrast,
the fine-tuned reranker (BGE-M3 + Multi-HyDE + FT-RR), trained out-of-domain on FinQA,
achieves comparable context quality, indicating that domain-adaptive reranking provides
a more robust and transferable improvement. The remaining gap to the oracle page condition
confirms that page-level retrieval remains the primary limitation.

% ──────────────────────────────────────────────────────────
\subsection{Effect of Retrieval Budget}
\label{sec:budget}
% ──────────────────────────────────────────────────────────

Figure~\ref{fig:recall_curves} shows how PageRec@k (left) and DocRec@k (right) evolve
as the retrieval budget $k$ increases for key methods.

\begin{figure}[ht]
    \centering
    \includegraphics[width=\linewidth]{figures/recall_curves.png}
    \caption{PageRec@$k$ (left) and DocRec@$k$ (right) at
      $k \in \{1, 3, 5, 10, 20\}$.}
    \label{fig:recall_curves}
\end{figure}

The PageRec@k curves (left panel) exhibit steady growth across all methods, with no clear
saturation within the evaluated range. This indicates that page-level recall remains
sensitive to retrieval budget, and that increasing $k$ continues to recover additional
relevant pages. However, the relative ordering of methods is largely preserved across
budgets, with approaches incorporating query reformulation and reranking consistently
outperforming standard dense retrieval. In particular, the fine-tuned reranker achieves
high page recall at relatively small $k$, reaching performance comparable to the best
baselines at larger budgets. This suggests that improvements in ranking quality can
partially substitute for increased retrieval depth, concentrating relevant evidence
earlier in the ranked list.

In contrast, the DocRec@k curves (right panel) saturate rapidly for dense methods.
Document-level recall reaches near-ceiling values at small $k$, and increasing the
retrieval budget yields only marginal gains. This behaviour indicates that document
discovery is largely solved under dense retrieval, and that additional budget does not
substantially improve performance at this level. The slower growth and lower ceiling of
BM25 further suggest that its limitations arise from lexical mismatch rather than
insufficient retrieval depth.

Taken together, these results reinforce the central finding that the primary challenge
lies in within-document (page) localisation rather than document retrieval. Increasing the
retrieval budget alone is insufficient to close this gap, methods that improve
ranking quality or focus retrieval on the most relevant regions of the document provide
more effective gains. 

% ──────────────────────────────────────────────────────────
\subsection{Performance by Question Type}
\label{sec:qtype_retrieval}
% ──────────────────────────────────────────────────────────

Table~\ref{tab:qtype} decomposes DocRec@5 and PageRec@5 by the three FinanceBench
question types for four configurations, together with the oracle document
bound.  Each question type contains 50 questions.  These four configurations were selected
to isolate the incremental contribution of each component (reranking, query reformulation,
and fine-tuning) relative to the dense baseline. The three question types exhibit different retrieval styles, and the relative
contribution of each pipeline component varies across them.


\begin{table}[ht]
\centering
\caption{DocRec@5 and PageRec@5 at $k=5$ by question type.
  MG = metrics-generated; DR = domain-relevant; NG = novel-generated.}
\label{tab:qtype}
\small
\begin{tabular}{lcccccc}
\toprule
\multirow{2}{*}{\textbf{Method}}
  & \multicolumn{2}{c}{\textbf{MG}} & \multicolumn{2}{c}{\textbf{DR}}
  & \multicolumn{2}{c}{\textbf{NG}} \\
\cmidrule(lr){2-3}\cmidrule(lr){4-5}\cmidrule(lr){6-7}
  & D@5 & P@5 & D@5 & P@5 & D@5 & P@5 \\
\midrule
Dense BGE-M3                       & 0.98 & 0.49 & 0.92 & 0.22 & 0.92 & 0.35 \\
BGE-M3 + ReRanker                  & 1.00 & 0.77 & 0.86 & 0.18 & 0.98 & 0.38 \\
BGE-M3 + Multi-HyDE + ReRanker    & 1.00 & 0.78 & 0.84 & 0.29 & 0.94 & 0.38 \\
BGE-M3 + Multi-HyDE + FT ReRanker & 0.98 & 0.80 & 0.90 & 0.32 & 0.96 &0.54 \\
\midrule
Oracle document                    & 1.00 & 0.68 & 1.00 & 0.53 & 1.00 & 0.58 \\
\bottomrule
\end{tabular}
\end{table}


For metrics-generated questions, all methods achieve strong performance, and reranking
provides a substantial improvement over the dense baseline. Both zero-shot and fine-tuned
rerankers surpass the oracle document condition, indicating that ranking alone is sufficient
to reliably surface the correct evidence. Notably, the gap between the zero-shot and
fine-tuned reranker is small, suggesting that additional domain-specific training provides
limited benefit for this question type.  This is because these questions typically ask for specific numerical values (e.g., revenue,
net income, or totals) that appear in well-defined parts of a document, such as tables or
summary sections. Once the correct document is identified, there are only a few plausible
locations where the answer can be found. As a result, a strong zero-shot reranker can already
identify these pages effectively, leaving little room for further improvement through
fine-tuning.

Domain-relevant questions are more challenging for all methods. Page-level performance
remains well below the oracle bound, and gains from reranking and fine-tuning are more
limited. These questions tend to use broader, descriptive language while also relying on
specialised financial terms, which creates a mismatch with how information is expressed
in the document. As a result, retrieval often returns multiple sections that are either
topically related or share similar terminology, but do not contain the exact target
evidence. This combination of ambiguity and domain-specific vocabulary makes it difficult
to reliably identify the correct page, and neither query reformulation nor reranking fully
resolves the issue.

For novel-generated questions, improvements from fine-tuning are more pronounced. The
fine-tuned reranker substantially improves page recall relative to both the dense baseline
and the zero-shot reranker, closing most of the gap to the oracle document condition.
This suggests that even when query phrasing differs from training data, learning
query–passage interactions can generalise to related evidence types. In contrast, query
reformulation and zero-shot reranking provide more modest gains, indicating that they are
less effective at capturing the implicit structure required to locate the correct page.

Overall, these results show that the effectiveness of each retrieval component depends
strongly on the structure of the query. When answers appear in predictable regions of a
document, ranking alone is often sufficient. When queries are ambiguous or rely on
domain-specific terminology, retrieval becomes more challenging, as it may return multiple
plausible matches or fail to retrieve the correct evidence altogether. For queries that
require more implicit reasoning or indirect connections to the evidence, fine-tuning
provides the most consistent improvements. Across all question types, however, the
remaining gap to the oracle page condition confirms that page-level localisation remains
the primary limitation.

% ──────────────────────────────────────────────────────────
\subsection{Performance by Document Type}
\label{sec:doctype_retrieval}
% ──────────────────────────────────────────────────────────


Figure~\ref{fig:doc_type} stratifies PageRec@5 and DocRec@5 by SEC filing type across
all evaluated configurations. The benchmark is dominated by 10-K annual reports
($n=112$), with substantially smaller samples for 10-Q quarterly reports ($n=15$),
8-K current reports ($n=9$), and earnings call transcripts ($n=14$). As a result,
the observations for non–10-K document types should be interpreted with caution,
and we avoid drawing strong comparative conclusions across document classes.
Future work should address this imbalance to enable a more statistically robust
evaluation.


\begin{figure}[ht]
    \centering
    \includegraphics[width=1\linewidth]{figures/butterfly_doc_type_ret.png}
    \caption{PageRec@5 (left bar, green) and DocRec@5 (right bar, purple) at $k=5$
      by filing type.}
    \label{fig:doc_type}
\end{figure}


Across all document types, document-level recall is near-perfect, including under the
oracle document condition, confirming that document retrieval is not the limiting factor.
Differences in performance are therefore driven by page-level localisation within each
document type.

On 10-K annual reports, the fine-tuned reranker achieves the strongest performance,
slightly exceeding both the best baseline and the oracle document condition for this
filing type. This is consistent with earlier results on metrics-style queries and reflects
the alignment between the training distribution and the structure of these documents.
Given the length and internal complexity of 10-K filings, improvements in page ranking
are critical, and both reranking and fine-tuning provide clear gains over standard dense
retrieval.

On 8-K current reports, all methods achieve high PageRec@5, including BM25. These
documents are short and focused, typically containing a small number of relevant sections,
which makes the target page easier to identify once the document is identified. In this
setting, even simple lexical matching is often sufficient, and the relative gains from
more advanced retrieval methods are limited.

On 10-Q quarterly reports, performance is more mixed. While the fine-tuned reranker
improves over the zero-shot baseline, it performs worse than the Multi-HyDE variant
without reranking, suggesting that the learned ranking signal does not fully align with
this document type. Although 10-Q reports share structural similarities with 10-K filings,
differences in content and presentation appear to reduce the effectiveness of the
fine-tuned model. In this setting, improving the query representation through
Multi-HyDE provides a stronger signal than applying a mismatched reranker, indicating
that query-side adjustments can be more robust than learned ranking when transfer is
imperfect.

Finally, earnings call transcripts are difficult for all methods. All configurations
remain below the oracle document bound, indicating that page-level localisation is still
challenging for this document type. Unlike other settings, BM25 performs on par with the
stronger neural methods, suggesting that simple lexical matching is stronger in
this more descriptive format. The zero-shot reranker achieves the best performance, while
the fine-tuned reranker performs noticeably worse, indicating that fine-tuning degrades
ranking in this case. This behaviour is likely due to a mismatch between the training data
and the target document type. The fine-tuned model is trained on structured documents with
clearly defined evidence patterns, whereas earnings call transcripts are more
conversational and less structured. As a result, the learned ranking signal does not align
well with this setting, and a more general zero-shot model provides a better signal. As in
other cases, document-level recall remains high, confirming that the failure is isolated
to page-level ranking.

Overall, these results suggest that retrieval performance is sensitive to document
structure and training distribution. Fine-tuning provides clear benefits in-domain,
particularly for long and structured documents, but can degrade performance when applied
out-of-distribution. Across all types, however, the remaining gap to the oracle page
condition confirms that page-level localisation remains the primary limitation.


% ──────────────────────────────────────────────────────────
\subsection{Retrieval Failure Analysis}
\label{sec:retrieval_failure}
% ──────────────────────────────────────────────────────────

The aggregate and stratified results reveal high-level structural patterns in where and
why retrieval fails. The following qualitative analysis provides a closer look at the
underlying mechanisms. Three failure modes are distinguished. A total failure
occurs when neither the gold document nor any of its gold pages appear in the top-$k$
results. A partial-document
failure occurs when the gold document is retrieved but the gold pages are missed. A partial-page failure occurs when the question has multiple gold pages and only a
subset are recovered. The subsequent analysis examines each question type through this lens,
grounded in concrete examples where they best illustrate the underlying mechanism.


\subsubsection{Metrics-Generated Questions}
\label{sec:fail_metgen}

% Metrics-generated questions require locating specific numerical values in financial
% tables. The gold evidence almost always resides in a standardised financial statement,
% making this category the most structurally predictable and the highest-recall category
% across all methods. Its failure modes are consequently more specific.

A dominant failure mode for metrics-generated questions under dense retrieval alone is
temporal confusion. In these cases, the system retrieves a structurally correct and plausible page
from the correct company, but from the wrong reporting period. This constitutes a
total retrieval error rather than a partial-document failure, as the gold
filing itself is not retrieved. Financial statements are highly standardised across
years, with nearly identical layouts, row labels, and column structures, differing
primarily in the reported values and year identifiers. Dense similarity scores collapse
this distinction, treating pages from different years as near-duplicates. Consider the following question from FinanceBench,

\begin{equation*}
\textit{What is the year-end FY2018 net PP\&E for 3M? Answer in USD billions.}
\end{equation*}

\noindent Table~\ref{tab:temporal_confusion_decomp} decomposes the retrieval outcome. As can be observed, every structural dimension is correct, the retrieved page is a balance sheet, from the
correct company, with the PP\&E row and the year-end column. The issue is that the filing year is
wrong. The system returns the FY2022 balance sheet instead of FY2018, yielding a
generated answer of \$9.17B rather than the correct \$8.70B. The side-by-side pages are
reproduced in Appendix~\ref{app:temp_conf}.

\begin{table}[H]
\centering
\small
\begin{tabular}{lcccc}
\toprule
\textbf{Component} & \textbf{Table} & \textbf{Document} & \textbf{Row} & \textbf{Column} \\
\midrule
\textbf{Ground truth}
  & Correct balance sheet & FY2018 & PP\&E (net) & Year-end (2018) \\
\textbf{Retrieved}
  & Correct balance sheet & \textbf{FY2022} & PP\&E (net) & Year-end (2022) \\
\bottomrule
\end{tabular}
\caption{Error decomposition for a temporal confusion failure. Table type, row, and
  column are all correct; the only error is the document-level filing year.}
\label{tab:temporal_confusion_decomp}
\end{table}

\noindent Once the cross-encoder reranker is introduced, this failure mode largely
disappears. The reranker performs joint scoring over the full query-page pair and attends
directly to the year tokens in both the question and the table header, demoting the
wrong-year page in favour of the temporally matching filing. This accounts for most of the PageRec@5 gain from 0.49 (dense BGE-M3) to 0.77 (BGE-M3 + ReRanker) on this question type.  The residual failures that persist after reranking are lexical gap failures of the same
kind described for novel-generated and domain-relevant questions in Section~\ref{sec:fail_novel}.

\subsubsection{Novel-Generated and Domain-Relevant Questions}
\label{sec:fail_novel}

Novel-generated and Domain-relevant questions require company-specific reasoning, synthesis across adjacent
passages, or computation that cannot be answered by locating a single sentence or table
cell.  The most consequential failure mode is a lexical gap between the query vocabulary
and the vocabulary of the answer-bearing page.  The following example illustrates this
directly (see Appendix ~\ref{app:lex_gap} for the full document pages and additional examples).

\begin{failex}{Lexical Gap / Domain Vocabulary Failure}
\label{ex:lexical}

\noindent\textbf{Question.}\enspace
If JPMorgan went bankrupt by the end of Q1 2021 and liquidated all of its
assets to pay shareholders, how much could each shareholder receive?

\medskip
\noindent\textbf{Expected page.}\enspace
p.\,5 (reports tangible book value per share (TBVPS) of
\textbf{\$66.56} as of March 31, 2021).

\noindent\textbf{Retrieved pages.}\enspace
p.\,81 (1st),\enspace p.\,79 (2nd),\enspace p.\,83 (3rd),\enspace
p.\,77 (4th),\enspace p.\,85 (5th).

\medskip
\noindent\textbf{Expected answer.}\enspace
\textbf{\$66.56} per share.

\medskip
\noindent\textbf{Generated answer.}\enspace
\textit{``Equity (assets minus liabilities) equals \$180{,}714 million.
With 4{,}104{,}933{,}895 shares outstanding, each shareholder would receive
approximately \textbf{\$43.99} per share.''}
\end{failex}
\vspace{1em}

The query uses informal language (``liquidate assets per shareholder'') that does not
lexically match the financial term \emph{TBVPS} reported in the
filing.  Retrieval consequently fails to surface the correct capital-metrics page and
instead ranks balance sheet pages containing assets and liabilities. While these pages
are semantically related, they require additional derivation and context to produce the answer.  The
generator performs a from-scratch calculation using retrieved values and produces a
numerically plausible but incorrect estimate, rather than reading off the directly
reported metric.  This failure mode is particularly resistant to reranking because the
reranker also lacks a signal connecting ``liquidate assets'' to ``tangible book value'',
it would require query expansion or domain glossary mappings to bridge the vocabulary
gap at retrieval time.

This example illustrates a broader failure pattern shared by both novel-generated and
domain-relevant questions. When the query uses language that does not align with how
the relevant concept is expressed in the document, retrieval surfaces pages that are
topically related but do not contain the exact evidence required. In such cases, the
correct page is often absent from the candidate set, limiting the effectiveness of
reranking. While fine-tuning improves performance by learning stronger query–passage
interactions, its gains are concentrated on questions where the target evidence follows
more familiar patterns. Queries requiring implicit mappings or less direct terminology
remain challenging, highlighting the importance of improving query representation or
vocabulary alignment at the retrieval stage.

\subsubsection{Annotation Limitation}
\label{sec:fail_annot}

A recurring source of false retrieval failures is benchmark annotation that designates a
single gold page per question even when the text required to answer appears on other
pages of the same filing.  Under the failure taxonomy, this example is classified as a partial-document failure, as the correct document is retrieved but the annotated gold page is not present within the top-$k$. Because PageRec@5 is computed against the annotated gold set,
retrieving an equally informative page that was not annotated counts as a miss,
artificially deflating recall while the downstream generator still produces the correct
answer.  The following example illustrates this directly.

\begin{failex}{Multi-Location Answer / Annotation Mismatch}
\label{ex:annot_mismatch}

\noindent\textbf{Question.}\enspace
What is the quantity of restructuring costs directly outlined in PepsiCo's income
statement for FY2022? If restructuring costs are not explicitly outlined, state 0.

\medskip
\noindent\textbf{Expected answer.}\enspace \textbf{\$411 million}.

\medskip
\noindent\textbf{Generated answer.}\enspace
\textit{``The quantity of restructuring costs directly outlined in PepsiCo's income
statement for FY2022 is \textbf{\$411 million}.''}

\medskip
\noindent\textbf{Annotated gold page (Appendix~\ref{fig:annot_gold}).}\enspace
The restructuring table presents \emph{Total restructuring and impairment charges}
of \$411M for FY2022 alongside prior-year comparatives.

\medskip
\noindent\textbf{Top-1 retrieved page (Appendix~\ref{fig:annot_ret}), not annotated.}\enspace
The Consolidated Statement of Cash Flows, ranked first by the retrieval system,
contains the line \emph{Restructuring and impairment charges} with the identical
figure of \$411M as a non-cash adjustment to net income.

\medskip
\noindent\textbf{Retrieval outcome.}\enspace
p.\,63 (1st, not annotated),\enspace
p.\,\textit{GT} (not retrieved within top-5).
\end{failex}


\vspace{1em}


The generator produces the correct answer from the retrieved cash-flow page because it
carries the identical monetary figure.  Yet the evaluation framework records PageRec@5
of zero because the gold annotation points exclusively to the restructuring-table page.
This pattern arises whenever a key figure is disclosed in multiple sections of the same
filing, which is common for summary metrics such as restructuring charges, depreciation,
and tax provisions.  It underscores why evaluation must span multiple granularities,
page-level recall alone cannot distinguish a genuine retrieval failure from an annotation
artefact, and only chunk-level and generative metrics can recover the true system utility
in such cases.

\subsubsection{Failure Mode Summary}
\label{sec:fail_summary}

Across question types, retrieval failures can be attributed to a small number of
recurring mechanisms. For metrics-generated questions, the primary issue is temporal
confusion, where structurally similar pages from different reporting periods are
indistinguishable under dense similarity. For domain-relevant questions, ambiguity and
domain-specific terminology lead to multiple plausible matches or missed evidence,
limiting the effectiveness of both reranking and query reformulation. For
novel-generated questions, lexical mismatch between query phrasing and document
terminology causes retrieval to surface related but incorrect pages, with fine-tuning
providing the most consistent improvements in this setting. In addition, annotation limitations introduce systematic discrepancies between measured
and actual performance, as valid evidence may appear on non-annotated pages. These
findings highlight that retrieval failures arise not only from recall limitations, but
from mismatches between query language, document structure, and evaluation assumptions.

This motivates the generation failure analysis in
Section~\ref{sec:gen_failure}, which examines the residual errors that persist even when
retrieval succeeds.


% ══════════════════════════════════════════════════════════
\section{Generative Evaluation}
\label{sec:gen_results}
% ══════════════════════════════════════════════════════════

This section evaluates downstream answer quality produced by the RAG pipeline under
three representative retrieval pipelines and the two oracle baselines. The primary
metrics are NumMatch on the 50 metrics-generated questions, as well as the full 150 question set, and human accuracy over all
150 questions, where human evaluation marks each generated answer as correct or incorrect.
ROUGE-L and BERTScore serve as surface-level consistency checks. Section~\ref{sec:gen_overall}
reports aggregate results and contextualises them against the oracle bounds established
in Section~\ref{sec:oracle_results}. Section~\ref{sec:gen_scaling} examines how answer
quality scales with generator size and compares our fine-tuned reranker within published
results on FinanceBench. Sections~\ref{sec:gen_qtype} and~\ref{sec:gen_doctype}
report breakdowns by question type and document type, incorporating both automatic
metrics and human accuracy to give a complete picture of where the system succeeds and
where it fails. Section~\ref{sec:gen_failure} analyses the residual generation failures
that persist even under oracle retrieval conditions.

% ──────────────────────────────────────────────────────────
\subsection{Overall Answer Quality}
\label{sec:gen_overall}
% ──────────────────────────────────────────────────────────

Table~\ref{tab:gen_overall} reports downstream answer quality for three retrieval
configurations and the two oracle baselines, using Qwen-2.5-7B-Instruct and
Llama-3.1-8B-Instruct with the fixed prompt template in Appendix~\ref{app:prompt}.
The proposed method refers to BGE-M3 dense retrieval combined with Multi-HyDE query
reformulation and the domain fine-tuned cross-encoder reranker (Multi-HyDE + FT-RR).
NumMatch~(m) is evaluated on the 50 metrics-generated questions, and NumMatch~(all) covers
all 150 questions.

\begin{table}[ht]
\centering
\small
\begin{adjustbox}{max width=\columnwidth}
\begin{tabular}{l l c c c c}
\toprule
\textbf{LLM} & \textbf{Retrieval} & \textbf{ROUGE-L} & \textbf{BERT-F1}
  & \textbf{NumMatch (m)} \\
\midrule
\multirow{5}{*}{Qwen2.5-7B}
& Dense BGE-M3                       & 0.13 & 0.08 & 0.32  \\
& Multi-HyDE + RR                    & 0.13 & 0.06 & 0.50  \\
& \textbf{Multi-HyDE + FT-RR (Ours)} & \textbf{0.14} & \textbf{0.08}
  & \textbf{0.52}  \\
& Oracle Doc                         & 0.13 & 0.08 & 0.44  \\
& Oracle Page                        & 0.17 & 0.10 & 0.70 \\
\midrule
\multirow{5}{*}{Llama3.1-8B}
& Dense BGE-M3                       & 0.12 & 0.08 & 0.36 \\
& Multi-HyDE + RR                    & 0.13 & 0.07 & 0.56 \\
& \textbf{Multi-HyDE + FT-RR (Ours)} & \textbf{0.13} & \textbf{0.10}
  & \textbf{0.52}  \\
& Oracle Doc                         & 0.13 & 0.08 & 0.46 \\
& Oracle Page                        & 0.17 & 0.13 & 0.78  \\
\bottomrule
\end{tabular}
\end{adjustbox}
\caption{Generative performance at $k=5$ for selected baselines, the proposed method
  (Multi-HyDE + FT-RR), and oracle conditions. NumMatch~(m) is evaluated on the 50
  metrics-generated questions.}
\label{tab:gen_overall}
\end{table}

A consistent pattern across both models is that improvements in answer quality are
primarily driven by retrieval rather than the choice of generator. The relative ordering
of methods is nearly identical for Qwen and Llama, and differences between the two models
are small compared to the gains obtained from improved retrieval. This suggests that,
in this setting, generation quality is largely constrained by the quality of the retrieved
context rather than model-specific capabilities.

The effect of improved retrieval is most clearly reflected in the numerical accuracy
metrics. As retrieval quality improves, NumMatch increases consistently across both
models, indicating that generation errors are largely tied to whether the correct
evidence is present in the context. In particular, the transition from dense retrieval
to the proposed method produces a substantial gain, suggesting that many of the errors
under baseline configurations arise from missing or incorrectly ranked pages rather
than limitations of the generator itself. This trend holds for both models, reinforcing
that retrieval quality is the dominant factor.

The fact that the proposed method matches or exceeds the oracle document condition
further highlights the importance of ranking. Restricting retrieval to the correct
document alone is insufficient if irrelevant pages are still prioritised within that
document. In contrast, the proposed method improves the ordering of evidence, resulting
in a more concentrated and useful context. The corresponding gains in NumMatch indicate
that generation quality benefits directly from this improved context, even without
changes to the generator.

At the same time, the gap to the oracle page condition provides a useful estimate of the
remaining generation error. Even when the correct pages are available, a non-trivial
fraction of answers remain incorrect, implying that generation failures persist
independently of retrieval. This is particularly evident in cases where the model must
extract or compute a value rather than directly copy it. In these settings, NumMatch
serves as a proxy for how effectively the generator can utilise the provided evidence,
and the observed gap suggests that improvements in reasoning or extraction are still
needed.

Finally, the relatively low ROUGE-L and BERTScore values across all configurations are
consistent with the nature of the task. Answers are often paraphrased or derived rather
than copied verbatim, so surface-level overlap metrics do not fully capture correctness.
In contrast, the behaviour of NumMatch across retrieval configurations provides a more
reliable signal of generation quality, and its alignment with retrieval improvements
confirms that most performance gains in this setting originate from better evidence
selection rather than changes in the generator.

% ──────────────────────────────────────────────────────────
\subsection{Generator Scaling and Comparison with Prior Work}
\label{sec:gen_scaling}
% ──────────────────────────────────────────────────────────

Table~\ref{tab:scaling} examines how answer quality changes as the generator is scaled
from 3B to 14B parameters within the Qwen~2.5 model family, with retrieval fixed to the
proposed method (Multi-HyDE + FT-RR). Results are shown under two conditions: top-$k$
retrieval, where the proposed pipeline supplies the context, and oracle context, where
the gold evidence pages are provided directly. Human accuracy, as described in Section ~\ref{sec:llm_eval} is a binary
correct/incorrect score conducted over all 150 questions under top-$k$ retrieval,
providing a broader view of answer quality that covers question types where numeric
matching is not applicable.

\begin{table}[ht]
\centering
\small
\setlength{\tabcolsep}{4pt} % reduce column spacing (default ~6pt)

\begin{tabular}{l ccc ccc}
\toprule
\multirow{2}{*}{\textbf{Metric}} 
  & \multicolumn{3}{c}{\textbf{Top-$k$ retrieval}} 
  & \multicolumn{3}{c}{\textbf{Oracle Page}} \\
\cmidrule(lr){2-4} \cmidrule(lr){5-7}
  & \textbf{3B} & \textbf{7B} & \textbf{14B} 
  & \textbf{3B} & \textbf{7B} & \textbf{14B} \\
\midrule
\textbf{NumMatch} 
  & 0.40 & 0.44 & 0.50 
  & 0.54 & 0.70 & 0.74 \\
\addlinespace
\textbf{Human Acc.} 
  & \makecell{55/150\\(36.67\%)} & \makecell{65/150\\(43.33\%)} & \makecell{76/150\\(50.67\%)} 
  & \makecell{74/150\\(49.33\%)} & \makecell{88/150\\(58.67\%)} & \makecell{109/150\\(72.67\%)} \\
\bottomrule
\end{tabular}

\caption{Generator performance as a function of model size, with retrieval fixed to
  Multi-HyDE + FT-RR. NumMatch is on 50 metrics-generated questions.}
\label{tab:scaling}
\end{table}


A clear scaling trend emerges across both metrics. Larger models consistently produce more accurate answers, and under top-$k$ (Multi-HyDE + FT-RR) retrieval both NumMatch and human accuracy improve steadily with model size, indicating that increased capacity leads to better numerical extraction, reasoning, and overall answer correctness. This confirms that scaling the generator alone provides meaningful gains even when retrieval is imperfect.

The effect of scaling becomes more pronounced when higher-quality context is available.
Under oracle page conditions, the improvement with model size is substantially larger
for both NumMatch and human accuracy. This indicates that larger models are not only
more capable in general, but are better able to utilise correct evidence when it is
present. In other words, scaling enhances the model’s ability to extract and reason over
relevant information, rather than compensating for missing or incorrect context.

Table~\ref{tab:sota} contextualises the proposed system within published results on the 150-question FinanceBench open-source subset, using human accuracy for comparability with prior work.


\begin{table}[ht]
\centering

\small
\begin{tabular}{llcc}
\toprule
\textbf{System} & \textbf{Generator} & \textbf{LLM Acc.} & \textbf{Manual Acc.} \\
\midrule
\citet{islam2023financebench}  & GPT-4-Turbo            & --   & 0.19 \\
\citet{jimeno2024financial}    & GPT-4                  & 0.33 & 0.37 \\
\citet{setty2024improving}     & GPT-4o                 & 0.26 & --   \\
\citet{wang2024finsage}        & Qwen2.5-72B-Instruct   & 0.50 & 0.57 \\
\citet{choe2025hirec}          & GPT-4o                 & 0.50 & --   \\
\midrule
Ours (Multi-HyDE + FT-RR)     & Qwen2.5-14B-Instruct   & 0.42 & 0.51 \\
\bottomrule
\end{tabular}
\caption{Comparison with published results on the 150-question FinanceBench subset.
  Dashes indicate metrics not reported by the original work.}
\label{tab:sota}
\end{table}

The proposed system with a 14B open-source generator achieves 51\% human accuracy and a 42\% LLM-as-a-Judge accuracy with an 81\% agreement rate between the two evaluations. This approximately matches FinSage's LLM-graded score (50\%) and narrowing the gap to FinSage's manual
evaluation (57\%) to 6 points. FinSage employs a multi-modal preprocessing pipeline, a
DPO-trained domain-specialised reranker, and a 72B generator. Our proposed system
relies on a fine-tuned cross-encoder trained exclusively on FinQA data and a fixed
open-source 14B generator with no proprietary components. The 7B variant (43.3\%)
already exceeds the closest comparable published system using human evaluation
\citep{jimeno2024financial} by 6 points, with substantially lower generator capacity.
These results confirm that improvements in page-level retrieval account for a large share
of the performance gain over the FinanceBench baseline, and that the remaining gap to
FinSage is plausibly attributable to their larger generator and DPO-trained reranker,
both of which address failure modes identified in Sections~\ref{sec:fail_metgen}
and~\ref{sec:gen_failure}.

% ──────────────────────────────────────────────────────────
\subsection{Performance by Question Type}
\label{sec:gen_qtype}
% ──────────────────────────────────────────────────────────

Table~\ref{tab:gen_qtype} and Figure~\ref{fig:gen_qtype} report answer quality by
question type for the proposed method (Multi-HyDE + FT-RR) and the dense baseline,
using Qwen-2.5-7B-Instruct. Figure~\ref{fig:gen_qtype} additionally shows human
accuracy for both the 7B and 14B models under top-$k$ retrieval and oracle context,
allowing a direct comparison of how model scale and retrieval quality interact across
question types.

\begin{table}[ht]
\centering
\small
\begin{adjustbox}{max width=\columnwidth}
\begin{tabular}{l l l c c c}
\toprule
\textbf{Retrieval} & \textbf{Model} & \textbf{Q-Type}
  & \textbf{ROUGE-L} & \textbf{BERT-F1} & \textbf{NumMatch} \\
\midrule
\multirow{3}{*}{Dense BGE-M3}
& \multirow{3}{*}{Qwen2.5-7B}
  & Metrics & 0.01  & -- & 0.32 \\
& & Domain  & 0.17  & 0.13    & --   \\
& & Novel   & 0.21  & 0.19    & --   \\
\midrule
\multirow{3}{*}{Multi-HyDE + FT-RR (Ours)}
& \multirow{3}{*}{Qwen2.5-7B}
  & Metrics & 0.01  & -- & 0.52 \\
& & Domain  & 0.19  & 0.14    & --   \\
& & Novel   & 0.22  & 0.20    & --   \\
\midrule
\multirow{3}{*}{Oracle Page}
& \multirow{3}{*}{Qwen2.5-7B}
  & Metrics & 0.02  & -- & 0.70 \\
& & Domain  & 0.22  & 0.17    & --   \\
& & Novel   & 0.25  & 0.24    & --   \\
\bottomrule
\end{tabular}
\end{adjustbox}
\caption{Generative performance by question type (Qwen2.5-7B). NumMatch is only
  defined for metrics-generated questions.}
\label{tab:gen_qtype}
\end{table}

\begin{figure}[ht]
    \centering
    \includegraphics[width=0.75\linewidth]{figures/gen_by_q_type.png}
    \caption{Human accuracy by question type for the 7B and 14B models under top-$k$
      retrieval (FT-RR) and oracle context.}
    \label{fig:gen_qtype}
\end{figure}

The NumMatch gains for metrics-generated questions directly track the retrieval gains
from Section~\ref{sec:qtype_retrieval}: the proposed method improves NumMatch~(m) from
0.32 to 0.52, a gain of 0.20 points attributable to retrieval improvement, while the
oracle page ceiling (0.70) identifies a further 0.18-point headroom that belongs to
generation failures analysed in Section~\ref{sec:gen_failure}.

Human accuracy in Figure~\ref{fig:gen_qtype} adds important nuance that NumMatch alone
does not capture. On metrics-generated questions, the 7B model reaches 46\% under
top-$k$ retrieval and 72\% under oracle context, while the 14B model reaches 54\% and
78\% respectively. The oracle gain for both models (roughly 26 points) is the largest
of any question type, confirming that retrieval quality is the dominant bottleneck here
and that both model sizes are capable of correctly answering a substantially larger
fraction of these questions given the right context.

Domain-relevant questions show a notably different scaling pattern. The 7B and 14B
models achieve identical human accuracy under top-$k$ retrieval (0.36 each), indicating
that the retrieval failure is so severe on this question type that model size provides
no additional benefit: both models receive essentially the same inadequate context and
produce similarly poor answers. Under oracle context, however, the 14B model reaches
0.64 against the 7B model's 0.44, a 20-point gap that does not appear under top-$k$
retrieval. This reveals that domain-relevant questions are simultaneously a retrieval
problem and a generation problem: the 14B model has substantially stronger reasoning
over narrative evidence when it is provided correctly, but retrieval failure masks this
capacity difference entirely under real operating conditions.

Novel-generated questions show the largest absolute human accuracy gain from model
scaling under top-$k$ retrieval: the 14B model (0.62) outperforms the 7B model (0.48)
by 14 points, a gap that is larger here than on any other question type. Under oracle
context the gap narrows slightly (0.76 vs.\ 0.60), suggesting that the 14B model's
advantage on novel questions is partly driven by better reasoning given imperfect
context, not only by better use of perfect evidence. The 7B model's oracle gain for
novel questions (0.48~$\to$~0.60) is smaller than for metrics questions
(0.46~$\to$~0.72), consistent with the retrieval analysis showing that novel questions
have a smaller remaining gap to the oracle document bound.

For domain-relevant and novel-generated questions, the proposed method also improves
ROUGE-L modestly over the dense baseline (domain: 0.17~$\to$~0.19; novel:
0.21~$\to$~0.22). These small gains are consistent with the retrieval analysis:
domain-relevant questions retain a large unresolved retrieval gap, and the generation
metrics reflect this difficulty rather than a generator limitation.

% ──────────────────────────────────────────────────────────
\subsection{Performance by Document Type}
\label{sec:gen_doctype}
% ──────────────────────────────────────────────────────────

Table~\ref{tab:gen_doc} and Figure~\ref{fig:gen_doctype} report generative performance
by SEC filing type for the proposed method (Multi-HyDE + FT-RR), the dense baseline,
and oracle page conditions, using Qwen-2.5-7B-Instruct. Figure~\ref{fig:gen_doctype}
shows human accuracy for both the 7B and 14B models under top-$k$ retrieval and oracle
context, providing a fuller picture of filing-type difficulty than automatic metrics
alone. As noted in the retrieval analysis, the non-10-K groups are small enough that figures should be read as directional rather than definitive. With that in mind, the results reveal a portfolio of qualitatively distinct failure modes rather than a single bottleneck, and model size interacts with those failures in ways that aggregate metrics alone would obscure.


\begin{table}[ht]
\centering
\small
\begin{adjustbox}{max width=\columnwidth}
\begin{tabular}{l l c c c}
\toprule
\textbf{Retrieval} & \textbf{Doc}
  & \textbf{ROUGE-L} & \textbf{BERT-F1} & \textbf{NumMatch} \\
\midrule
\multirow{4}{*}{Dense BGE-M3}
& 10-K  & 0.09 & 0.04 & 0.32 \\
& 10-Q  & 0.17 & 0.13 & 0.20 \\
& 8-K   & 0.33 & 0.31 & 0.44 \\
& Earn. & 0.23 & 0.25 & 0.50 \\
\midrule
\multirow{4}{*}{Multi-HyDE + FT-RR (Ours)}
& 10-K  & 0.11 & 0.04 & 0.43 \\
& 10-Q  & 0.16 & 0.12 & 0.20 \\
& 8-K   & 0.33 & 0.32 & 0.44 \\
& Earn. & 0.24 & 0.24 & 0.57 \\
\midrule
\multirow{4}{*}{Oracle Page}
& 10-K  & 0.13 & 0.04 & 0.55 \\
& 10-Q  & 0.19 & 0.19 & 0.27 \\
& 8-K   & 0.40 & 0.42 & 0.44 \\
& Earn. & 0.24 & 0.22 & 0.71 \\
\bottomrule
\end{tabular}
\end{adjustbox}
\caption{Generative performance by filing type (Qwen2.5-7B). NumMatch is computed over
  all questions within each filing-type group; questions with non-numeric reference
  answers contribute 0 to the average.}
\label{tab:gen_doc}
\end{table}

\begin{figure}[ht]
    \centering
    \includegraphics[width=0.75\linewidth]{figures/gen_q_by_doc_type.png}
    \caption{Human accuracy by filing type for the 7B and 14B models under top-$k$
      retrieval (FT-RR) and oracle context.}
    \label{fig:gen_doctype}
\end{figure}


Performance varies substantially across filing types, and these differences provide a
clear lens into where the RAG pipeline succeeds or fails. Rather than reflecting purely
model capability, the results reveal that each document type induces a distinct balance
between retrieval difficulty and generation difficulty.



On 10-K annual reports, fine-tuning delivers the largest NumMatch gain across filing
types (0.321~$\to$~0.429), consistent with the retrieval gains in
Section~\ref{sec:doctype_retrieval}. Human accuracy from Figure~\ref{fig:gen_doctype}
shows 7B at 43\% and 14B at 44\% under top-$k$ retrieval, nearly identical despite the
parameter difference. Under oracle context, however, the 14B model reaches 72\% against
the 7B's 57\%, a 15-point gap. This pattern is consistent with the question-type
finding: under imperfect retrieval both models receive similar noisy context and produce
similarly limited answers, but the 14B model exploits clean evidence substantially
better when it is available. The oracle page NumMatch ceiling for 10-K (0.545) and the
14B oracle human accuracy (72\%) together suggest that further retrieval improvement
would yield meaningful gains, particularly for the larger model.

On 8-K current reports, both NumMatch (0.444 across all conditions) and human accuracy
tell the same story. The 7B model reaches 67\% under both top-$k$ retrieval and oracle
context, confirming that evidence localisation is not the binding constraint: the
correct page is found regardless of retrieval method, and the remaining errors are
generation failures. The 14B model reaches 89\% under both conditions, showing that
the larger model resolves many of the generation-stage failures that the 7B model
cannot, without needing any retrieval improvement to do so.

On 10-Q quarterly reports, human accuracy reveals a striking pattern not visible in
NumMatch alone. The 7B model improves substantially from top-$k$ (40\%) to oracle
(73\%), a 33-point gain that is the largest oracle benefit of any filing type for the
7B model, indicating that retrieval is a severe bottleneck for this filing type.
The 14B model, by contrast, achieves 67\% under both top-$k$ and oracle conditions,
showing essentially no benefit from perfect retrieval. This divergence suggests that
the 14B model has a different failure mode on 10-Q filings: it is not bottlenecked
by retrieval quality to the same degree as the 7B model, but it also cannot improve
beyond 67\% even with perfect evidence, pointing to a generation-stage ceiling that
model scaling alone does not resolve within the evaluated range.

Earnings call transcripts show the largest model-scaling effect of any filing type
under top-$k$ retrieval: the 14B model (64\%) outperforms the 7B model (36\%) by 28
points, a gap far larger than on any other filing type. This suggests that the 14B
model is substantially better at handling the conversational structure of transcripts
even when retrieval is imperfect, which the 7B model cannot compensate for. Under
oracle context the gap narrows (71\% vs.\ 50\%), indicating that some of the 7B
model's difficulty is retrieval-driven, but the 14B model's advantage persists even
with perfect evidence. The NumMatch figures for earnings calls should still be
interpreted with caution given the small group size, but the human accuracy figures,
which cover all question types in the group, are more reliable and point to a genuine
capacity difference on conversational document formats.

% ──────────────────────────────────────────────────────────
\subsection{Generation Failure Analysis}
\label{sec:gen_failure}
% ──────────────────────────────────────────────────────────

Under oracle page conditions, NumMatch~(m) is 0.70, meaning roughly 30\% of numeric
errors persist even with perfect retrieval. Similarly, we achieve 72.67\% human accuracy in the Oracle framework. These residual failures fall into four
patterns.

First, noisy tabular context. When the correct page contains many numbers organised in a dense table, the generator
frequently reads off the wrong value even though the right one is present. It may
confuse rows or columns that are adjacent in the table, particularly when multiple
reporting periods appear side by side. This explains a pattern that appears across
several questions: ROUGE-L is reasonable, indicating that the retrieved content is
topically correct, but NumMatch is zero because the specific number extracted is wrong.

\textbf{Multi-step arithmetic failures.}
Some questions require the generator to combine two or more figures from the retrieved
context to compute a final answer, rather than reading a directly stated value. At 7B
and 3B scale, errors in this kind of multi-step computation are common even when the
correct evidence is provided. The 14B model narrows this gap substantially
(oracle context NumMatch 0.74), confirming that this type of reasoning scales more
reliably with model size than surface-form generation does.

\textbf{Positional bias.}
When a question requires evidence from two pages simultaneously, both must appear in the
retrieved context for the generator to produce the correct answer. When one of the
required pages is ranked low and consequently appears at the tail of the context window,
the generator tends to overlook it, a behaviour consistent with the well-documented
lost-in-the-middle phenomenon in long-context LLMs~\citep{liu2024lost}.
Example~\ref{ex:positional} illustrates this: both required pages are retrieved, but the
page containing one of the two required values is ranked last, causing the generator to
report that the information is unavailable despite it being present in the context.

\begin{failex}{Positional Bias --- Adobe FY2015 Operating Cash Flow Ratio}
\label{ex:positional}

\noindent\textbf{Question.}\enspace
What is the FY2015 operating cash flow ratio for Adobe? Operating cash flow ratio is
defined as cash from operations divided by total current liabilities. Round to two
decimal places.

\medskip
\noindent\textbf{Expected pages.}\enspace
p.\,58 (total current liabilities) and p.\,62 (cash from operations).

\noindent\textbf{Retrieved pages.}\enspace
\textbf{p.\,62} (1st),\enspace p.\,64 (2nd),\enspace p.\,60 (3rd),\enspace
p.\,71 (4th),\enspace \textbf{p.\,58} (5th).

\medskip
\noindent\textbf{Expected answer.}\enspace
$\$1{,}469{,}502 \div \$2{,}213{,}556 \approx \mathbf{0.66}$

\medskip
\noindent\textbf{Generated answer.}\enspace
\textit{``The net cash provided by operating activities for FY2015 is \$1,469,502
thousand. However, the total current liabilities are not directly provided in the given
context. We cannot compute the operating cash flow ratio.''}

\end{failex}
\vspace{1em}

Both required pages are retrieved, but the page containing total current liabilities
(p.\,58) is ranked 5th and falls at the tail of the context window. The generator
correctly reads the first required value from the top-ranked page but fails to attend
to the second value, producing a spurious ``not determinable'' response despite the
evidence being present.

\textbf{Domain knowledge deficit.}
Certain questions require the generator to select the correct inputs to a formula from
among several plausible candidates that appear on the same page. Without sufficient
domain knowledge to distinguish which inputs are appropriate, the generator may apply a
formula that is arithmetically valid but conceptually incorrect, producing a wrong answer
despite correct retrieval and sound computation. Example~\ref{ex:domain} illustrates
this: retrieval is correct, the arithmetic is correct, but the generator includes an
input that the question specification excludes.

\begin{failex}{Domain Knowledge Deficit --- Netflix FY2015 EBITDA Margin}
\label{ex:domain}

\noindent\textbf{Question.}\enspace
What is the FY2015 unadjusted EBITDA\% margin for Netflix? Calculate using unadjusted
operating income and D\&A from the cash flow statement.

\medskip
\noindent\textbf{Expected pages.}\enspace
p.\,39 (income statement) and p.\,41 (cash flow statement).

\noindent\textbf{Retrieved pages.}\enspace
p.\,39 (1st),\enspace p.\,41 (3rd) --- retrieval is correct.

\medskip
\noindent\textbf{Expected answer.}\enspace
$\mathbf{5.4\%}$\enspace
$\bigl(\$62{,}283 + \$305{,}826\bigr) \div \$6{,}779{,}511$

\medskip
\noindent\textbf{Generated answer.}\enspace
$\mathbf{6.6\%}$\enspace
$\bigl(\$62{,}283 + \$79{,}380 + \$305{,}826\bigr) \div \$6{,}779{,}511$

\end{failex}

\vspace{1em}
The cash flow statement contains two separate line items that could both be interpreted
as depreciation and amortisation inputs. The correct answer uses only one of them
(\$305{,}826K); the generator includes both, adding a second component (\$79{,}380K)
that the question specification does not intend. The arithmetic is correct given the
generator's interpretation, but the interpretation itself is wrong. This failure cannot
be resolved through better retrieval; it requires the generator to understand the
precise definition of the formula being applied.

These four failure patterns collectively explain the 0.30-point gap between oracle page
NumMatch (0.70) and perfect generation. Noisy tabular context and domain knowledge
deficits are failures of comprehension given correct evidence; positional bias is a
failure of attention over the context window; and arithmetic errors are a capacity
limitation that diminishes with model scale. All four are generation-stage phenomena
that improved retrieval alone cannot resolve.

% ══════════════════════════════════════════════════════════
\section{Benchmark Extension Results}
\label{sec:crossdomain}
% ══════════════════════════════════════════════════════════

This section presents initial results on the benchmark extension described in
Section~\ref{sec:method_benchmark}, which combines FinanceBench with a FinQA-derived
subset to evaluate retrieval methods in a larger, mixed-document setting. These results
should be read as a first evaluation pass rather than a definitive study: several known
issues remain and the generation component has not yet been evaluated on the FinQA
subset. The primary contribution is to characterise how each retrieval method behaves
when the evaluation corpus becomes larger and more diverse, and to identify the
directions that require further work.

% ──────────────────────────────────────────────────────────
\subsection{Extended Corpus Statistics}
\label{sec:extension_corpus}
% ──────────────────────────────────────────────────────────

The benchmark extension constructs a shared global index by combining the original
FinanceBench corpus with a PDF-aligned FinQA test subset. The FinQA subset consists of
530 questions drawn from the 1,147-question FinQA test set, filtered to cases where the
corresponding PDF filing is available locally. Gold page annotations were derived by
fuzzy-matching the FinQA evidence text against extracted PDF page text as described in
Section~\ref{sec:method_benchmark}, with match scores ranging from 0.71 to 1.00 and a
mean of 0.88. There is no document overlap between the two components: all retrieval
is performed against the global 206-document index, meaning that FinanceBench questions
compete against 122 unrelated FinQA filings and vice versa. Table~\ref{tab:extension_stats}
summarises the corpus statistics.

\begin{table}[ht]
\centering

\small
\begin{tabular}{lrrr}
\toprule
\textbf{Statistic} & \textbf{FinanceBench} & \textbf{FinQA (subset)} & \textbf{Global} \\
\midrule
Questions          & 150    & 503     & 653    \\
Unique documents   & 84     & 122     & 206    \\
Total pages (est.) & 11,948 & 18,400  & 30,400 \\
Total chunks       & 14,808 & 22,836  & 37,644 \\
Avg.\ pages / doc  & 142    & 151     & 147    \\
\bottomrule
\end{tabular}
\caption{Corpus statistics for the benchmark extension.}
\label{tab:extension_stats}
\end{table}

% ──────────────────────────────────────────────────────────
\subsection{Retrieval Results}
\label{sec:extension_results}
% ──────────────────────────────────────────────────────────

Table~\ref{tab:ext_global} reports full retrieval metrics for all methods evaluated
against the global 206-document index across all 653 questions. These numbers are not
directly comparable to Table~\ref{tab:retrieval}, which used the original 84-document
index: adding 122 additional distractor documents approximately doubles the candidate
pool, and every method scores lower on the FinanceBench subset as a result. 

\begin{table}[ht]
\centering
\scriptsize
\renewcommand{\arraystretch}{0.88}
\setlength{\tabcolsep}{3pt}
\begin{tabularx}{\textwidth}{Xcccc}
\toprule
\textbf{Method} & \textbf{DocRec@5} & \textbf{PageRec@5} & \textbf{BLEU@5} & \textbf{ROUGE-L@5} \\
\midrule
\multicolumn{5}{l}{\textit{Sparse}} \\
BM25                       & 0.81 & 0.35 & 0.22 & 0.39 \\
SPLADE                     & 0.79 & 0.21 & 0.18 & 0.32 \\
\midrule
\multicolumn{5}{l}{\textit{Dense}} \\
Dense BGE-M3               & 0.93 & 0.31 & 0.23 & 0.37 \\
BGE-M3 + RR                & 0.93 & 0.28 & 0.21 & 0.35 \\
\midrule
\multicolumn{5}{l}{\textit{Hybrid}} \\
RRF 50/50                  & 0.92 & 0.39 & 0.27 & 0.43 \\
\midrule
\multicolumn{5}{l}{\textit{Hierarchical and query reformulation}} \\
Parent-Child               & 0.91 & 0.33 & 0.05 & 0.16 \\
Query Expansion            & 0.93 & 0.31 & 0.23 & 0.37 \\
BGE-M3 + HyDE              & 0.84 & 0.29 & 0.22 & 0.36 \\
BGE-M3 + Multi-HyDE        & 0.85 & 0.30 & 0.22 & 0.36 \\
BGE-M3 + Multi-HyDE + RR   & 0.89 & 0.29 & 0.22 & 0.36 \\
\midrule
\multicolumn{5}{l}{\textit{Proposed}} \\
BGE-M3 + FT-RR             & 0.94 & 0.30 & 0.22 & 0.36 \\
Multi-HyDE + FT-RR (Ours)  & 0.92 & 0.33 & 0.24 & 0.38 \\
\midrule
Oracle document            & 1.00 & 0.42 & 0.28 & 0.43 \\
Oracle page                & 1.00 & 1.00 & 0.49 & 0.70 \\
\bottomrule
\end{tabularx}
\caption{Global retrieval performance over all 653 questions in the combined
  206-document index at $k=5$. Bold marks the best non-oracle result per column.}
\label{tab:ext_global}
\end{table}

Table~\ref{tab:ext_breakdown} provides a per-corpus breakdown of DocRec@5 and PageRec@5
for a selection of methods that best illustrate the cross-corpus differences. 

\begin{table}[ht]
\centering

\small
\begin{tabular}{lcccc}
\toprule
\multirow{2}{*}{\textbf{Method}}
  & \multicolumn{2}{c}{\textbf{FB subset}}
  & \multicolumn{2}{c}{\textbf{FinQA subset}} \\
\cmidrule(lr){2-3}\cmidrule(lr){4-5}
  & D@5 & P@5 & D@5 & P@5 \\
\midrule
BM25                      & 0.53 & 0.11 & 0.89 & 0.42 \\
Dense BGE-M3              & 0.93 & 0.34 & 0.93 & 0.30 \\
BGE-M3 + RR               & 0.91 & 0.43 & 0.94 & 0.24 \\
RRF 50/50                 & 0.87 & 0.28 & 0.93 & 0.43 \\
BGE-M3 + Multi-HyDE       & 0.85 & 0.45 & 0.86 & 0.26 \\
Multi-HyDE + FT-RR (Ours) & 0.91 & 0.46 & 0.92 & 0.29 \\
\midrule
Oracle document           & 1.00 & 0.53 & 1.00 & 0.39 \\
\bottomrule
\end{tabular}
\caption{DocRec@5 and PageRec@5 broken down by corpus subset for selected methods.
  FB = FinanceBench subset; FQ = FinQA subset; both evaluated within the global index.}
\label{tab:ext_breakdown}
\end{table}

% ──────────────────────────────────────────────────────────
\subsection{Analysis}
\label{sec:extension_analysis}
% ──────────────────────────────────────────────────────────

Table~\ref{tab:ext_breakdown} reveals several inversions between the two corpora that
are directly informative about what each retrieval method actually learns.

The most striking inversion is BM25, which achieves PageRec@5 of 0.110 on
FinanceBench but 0.419 on FinQA, the single best non-oracle result on that corpus.
The explanation lies in how the two benchmarks were constructed. FinQA questions were
written by annotators working directly from specific tables in filings; the questions
reproduce exact company names, fiscal year labels, and table header terms that appear
verbatim on the answer page, giving BM25 near-perfect lexical signal. FinanceBench
questions are concept-based and rarely share tokens with the answer page, making lexical
matching ineffective. Hybrid RRF 50/50 benefits directly from this: incorporating
BM25's strong lexical signal lifts PageRec@5 to 0.425 on FinQA while its
FinanceBench score (0.283) reflects the noise that BM25 introduces on concept-based
queries.

The opposite inversion occurs for Multi-HyDE, which is among the strongest methods on
FinanceBench (PageRec@5 of 0.447) but falls below the dense baseline on FinQA (0.258).
HyDE generates a hypothetical passage that would answer the query. On FinanceBench,
this hypothetical resembles the evidence page because questions ask about specific
reported values. On FinQA, many questions require numerical computation over a table;
the hypothetical answer is a single computed number rather than a passage resembling the
source table, so the generated hypothesis actively points the retriever away from the
structured evidence.

Dense BGE-M3 shows the smallest performance gap between the two corpora (0.343 vs
0.302), confirming it as the most stable backbone across query styles. The proposed
method (Multi-HyDE + FT-RR) is the strongest on the FinanceBench subset (0.462) but
does not outperform the plain dense baseline on FinQA (0.291 vs 0.302). The
cross-encoder reranker performs even worse on FinQA (BGE-M3 + RR: 0.243), below the
dense baseline. Two factors explain this: the cross-encoder was not trained on financial
tables and may down-rank table pages whose extracted text does not read like natural
language; and the fine-tuned reranker was trained on FinQA evidence formatted as
structured text fields from the original JSON annotations, while the evaluation index
stores raw PDF-extracted text where tables are linearised differently, creating a format
mismatch that limits its benefit on the corpus it was nominally trained on.




%%%%%%%%%%%%%%%%%%%%%%%%%%%%%%%%%%%%%%%%%%%%%%%%%%%%%
%%              CHAPTER 5: CONCLUSION               %%
%%%%%%%%%%%%%%%%%%%%%%%%%%%%%%%%%%%%%%%%%%%%%%%%%%%%%
\chapter{Conclusion}

\section{Summary of Contributions}

This thesis investigated retrieval-augmented generation for financial question answering
over long SEC filings, with retrieval as its central object of study. Five
contributions were made.

The first is an oracle-based evaluation framework that decomposes retrieval failure into
document-level and page-level components. By progressively restricting the candidate pool
to the gold filing and then to the annotated gold pages, the framework yields empirical upper
bounds at each granularity and a two-gap diagnostic: the cross-document retrieval gap
$\Delta_{\text{page}}$ measures how much page recall is lost once document identity is held
perfect, and the generation ceiling gap $\Delta_{\text{gen}}$ quantifies the maximum gain in
downstream answer accuracy that better page localisation alone could deliver. Applied to
FinanceBench, the framework shows that standard pipelines achieve reasonable document
recall, from 0.88 to 0.95 across baselines, but page recall lags far behind, ranging from 0.11
to 0.48. The generation ceiling gap of 0.26 numeric-match points between oracle-document
(0.44) and oracle-page (0.70) conditions confirms that within-document localisation, not
document discovery, is the binding bottleneck on end-to-end answer quality.

The second contribution is a hierarchical page-then-chunk retrieval pipeline centred on
a learned page scorer. A bi-encoder fine-tuned on FinanceBench page-level data via the
Multiple Negatives Ranking loss first selects the top-$P$ pages from the full corpus; standard
dense chunk retrieval is then restricted to those pages. The page scorer achieves
PageRec@5 of 0.55, closing the gap to the oracle-document ceiling of 0.60 from 0.14 points
(best baseline) to 0.05. On metrics-generated questions, where evidence resides on
structurally predictable financial statement pages, the scorer reaches PageRec@5 of 0.80,
exceeding the oracle-document bound of 0.68 because its cross-attention scoring explicitly
captures the alignment between a numerical query and a financial table. The same
hierarchical architecture is used to evaluate ColPali as an alternative Stage 1 page scorer,
replacing the text-based bi-encoder with a vision-language model that represents pages as
patch embeddings. In standalone mode, ColPali achieves PageRec@5 of only 0.01,
reflecting a severe domain mismatch with the dense numerical layouts of SEC filings. As a
reranker over BGE-M3 candidates, it actively degrades page recall from 0.44 to 0.29,
indicating that patch-level visual similarity is dominated by typography and layout signals
rather than the numerical content that distinguishes answer-bearing pages from structurally
similar neighbours. These results establish the conditions under which text-based and
visual page scoring succeed and fail on financial documents.

The third contribution is a domain fine-tuned cross-encoder reranker, trained on FinQA
financial page data via binary cross-entropy loss and evaluated zero-shot on the held-out
FinanceBench benchmark. Combined with Multi-HyDE query reformulation, this reranker
achieves PageRec@5 of 0.56 and improves numeric match on metrics-generated questions
from 0.32 for the dense baseline to 0.52, a 63\% relative gain. With a 14B generator,
the full pipeline reaches 51\% human accuracy on FinanceBench, matching the LLM-graded
score of FinSage, a system that uses a 72B generator and a DPO-trained domain-specialised
reranker, using only open-source components and a fixed prompt.

The fourth contribution is a systematic multi-dimensional evaluation that stratifies
retrieval and generation performance by question type and SEC filing type, and complements
aggregate metrics with qualitative failure analysis and human evaluation. The stratified
analysis reveals that the benefit of domain fine-tuning is concentrated in metrics-generated
and novel-generated questions, where evidence resides on financial statement pages well
represented in the FinQA training corpus, while domain-relevant questions remain the
hardest subgroup because FinQA contains virtually no MD\&A narrative or policy disclosure
evidence. The filing-type breakdown shows that the fine-tuned reranker degrades on
earnings call transcripts by 0.14 PageRec@5 points relative to the zero-shot baseline, a
clear out-of-distribution failure driven by the structural mismatch between conversational
transcripts and the regulatory tabular filings on which the model was trained. The
qualitative failure analysis identifies four recurring generation failure modes under oracle
retrieval: noisy tabular context, multi-step arithmetic errors, positional bias over long
context windows, and domain knowledge deficits in formula input selection. These patterns
collectively explain the 0.30-point gap between oracle-page numeric match (0.70) and
perfect generation.

The fifth contribution is a benchmark extension pipeline that aligns questions from the
LoFIN/HiREC benchmark to verified gold PDF pages in EDGAR filings via fuzzy evidence
matching. Of 530 FinQA rows processed, 503 (94.9\%) were aligned to a PDF page with
a match score at or above the acceptance threshold of 0.70. The SEC-QA portion of
LoFIN could not be included because its available records did not contain gold evidence
text suitable for page alignment, leaving PDF page supervision unavailable. The combined
206-document index yields a cross-corpus evaluation that reveals a sharp method inversion:
BM25, which fails severely on FinanceBench due to lexical mismatch (PageRec@5 of 0.11),
achieves PageRec@5 of 0.42 on FinQA because its questions were written directly from
filing tables and share verbatim tokens with evidence pages. Multi-HyDE shows the
opposite pattern, performing strongly on FinanceBench (0.45) but falling below the dense
baseline on FinQA (0.26) because a hypothetical numeric answer does not resemble the
source table that a FinQA question requires. These inversions demonstrate that no single
retrieval design dominates across both question styles, and motivate the continued
development of a larger and more diverse PDF-native benchmark.


\section{Limitations and Future Work}

Several limitations of the present work point to concrete directions for future research.

The evaluation is built on 150 FinanceBench questions with a strong skew toward 10-K
annual reports (75\% of questions). The poor performance on earnings call transcripts, where
the fine-tuned reranker degrades relative to the zero-shot baseline, indicates that the learned
page scorer does not generalise to filing types that were absent from its training distribution.
Addressing this requires either a more diverse training corpus that includes conversational
and transcript-style filings, or a mixture-of-experts design that routes queries to filing-type-specific scorers.

The page scorer is trained on 170 positive pairs derived from 150 FinanceBench
annotations, with five-fold document-level cross-validation over 84 documents. This is a
very small fine-tuning regime for a large multilingual encoder, and the reported
cross-validation results should be interpreted cautiously given the limited training signal.
Hard negative mining, drawing pages from within the same filing that are topically related
to the query but do not contain the answer, would provide substantially more informative
training signal than the in-batch negatives currently used. This is especially important for
temporal near-duplicates, where adjacent years' financial statements share identical
structure and the model must learn fine-grained temporal grounding from examples that
genuinely stress that distinction.

The generation component was deliberately kept simple: top-$k$ retrieved chunks are
concatenated into a fixed in-context prompt, with no generator fine-tuning, no learned
context ordering, and no task-specific reasoning scaffold. The failure analysis in Section 4.3.5
shows that multi-step arithmetic errors account for a meaningful fraction of residual
generation failures even under oracle retrieval. Future work should address these
generation-side bottlenecks directly, for example through instruction tuning on financial QA
data, stronger context ordering that promotes answer-bearing pages to the beginning of the
context window, tool-assisted arithmetic for computed-metric questions, and evaluation of
larger or more domain-specialised language models.

The benchmark extension currently includes only the FinQA-derived portion of
LoFIN/HiREC, and the generation component has not yet been evaluated on this subset.
Completing the extension requires integrating the SEC-QA portion, which demands either
a richer release of SEC-QA evidence text or an alternative alignment procedure that does
not rely on gold passage text. Beyond completing the existing datasets, richer question-type
labels and reasoning annotations should be added to the FinQA-aligned data, since many
of its questions are numerical in style but lack the metadata needed to stratify evaluation
the way FinanceBench allows. Synthetic question generation tied to verifiable PDF evidence
spans could accelerate this process, provided that generated questions are filtered for
answerability and cover the filing types, such as 10-Q quarterly reports, proxy statements,
and multi-document comparison questions, that are underrepresented in the current corpus.

ColPali's failure in standalone mode raises questions that the present evaluation does not
resolve. It remains unclear whether the failure stems primarily from resolution constraints
during inference, from the absence of SEC filing examples in ColPali's training distribution,
or from the fundamental difficulty of distinguishing dense numerical tables by visual
similarity alone. Domain-specific fine-tuning of the vision-language encoder, evaluation at
higher rendering resolution, and the multi-scale pyramid indexing approach of Fang et al.
[2025] are all promising directions that this thesis did not explore. A hybrid architecture
that applies ColPali scoring selectively on visually complex pages while falling back to the
text-based bi-encoder on narrative sections would test whether the two modalities are
complementary at a finer granularity than the page level.

Additional directions for future work include combining page scoring with query
augmentation in a jointly optimised pipeline, extending the oracle analysis and page scorer
framework to multi-document questions that require cross-filing evidence aggregation, and
exploring reinforcement learning for iterative retrieval that refines page selection based on
generation feedback.


\section{Broader Implications}

Beyond financial question answering, the contributions of this thesis bear on the broader
problem of retrieval in long, structured documents. The oracle-based evaluation framework
is domain-agnostic and can be applied wherever multi-level annotation is available,
including legal case files with section-level evidence labels, clinical records with note-type
metadata, and academic papers with section boundaries. The consistent finding that
within-document localisation is the dominant retrieval bottleneck challenges the common
practice of measuring and optimising RAG pipelines primarily at the document level. It
suggests that intermediate retrieval units, such as pages, sections, and chapters, deserve
more systematic attention as first-class retrieval targets rather than as preprocessing
artefacts.

The scaling analysis provides a practical guideline for practitioners: larger generators
amplify rather than compensate for retrieval failures. The gap between top-$k$ and oracle
human accuracy grows from 12.7 points at 3B to 22.0 points at 14B, meaning that
investing in generator scale without addressing page-level retrieval yields diminishing
returns. Retrieval improvement and generator improvement are complements, not substitutes,
and the oracle framework introduced here provides a principled way to determine which of
the two is the current bottleneck before allocating engineering effort.

The cross-corpus inversion between FinanceBench and FinQA further illustrates that
aggregate benchmark scores can mask fundamentally different retrieval regimes. A system
that appears strong on one benchmark may be exploiting properties specific to how that
benchmark was constructed rather than learning general retrieval competence. Evaluation
across multiple question sources, as the benchmark extension begins to enable, is necessary
to make generalisable claims about financial RAG systems.

The diagnostic framework and the fine-tuned page scorer introduced in this thesis are
intended as practical tools for practitioners building RAG systems over long, structured
documents, as much as they are contributions to the academic evaluation of financial QA.
The central message is simple: in long-document retrieval, finding the right page matters
as much as finding the right document, and the gap between these two levels of recall is
where most of the performance headroom currently lies.

\bibliographystyle{plainnat}
\def\bibname{R\'ef\'erences}
\bibliography{ref}

\appendix

\chapter{Prompt Templates}
\label{app:prompt}

\section{Generator Prompt Template}
\label{app:gprom}
The prompt template used for all generation experiments is reproduced below.


\begin{verbatim}
System: You are a financial analyst answering questions based on SEC filings.
Use ONLY the provided context. If the context does not contain the answer,
say "I cannot determine this from the provided information." 

Context: {retrieved_evidence}

Question: {query}

Answer:
\end{verbatim}


\section{HyDE and Multi-HyDE Prompt Template}
\label{fig:mHprompt}

\begin{verbatim}
System: You are a financial analyst reviewing SEC filings.
Given the following question, write a short passage (3–4 sentences)
that would directly answer the question, as if excerpted from a 
company's SEC filing. Include specific financial figures and 
terminology as they would appear in the filing.

Question: {question}

Passage from annual report:
\end{verbatim}

\chapter{FinQA Data}

\begin{figure}[H]
    \centering
    \includegraphics[width=0.9\linewidth]{figures/finqa_ex.png}
    \caption{FinQA example. From https://finqasite.github.io/}
    \label{fig:finqa_ex}
\end{figure}

\chapter{Document Page Examples}

\section{Temporal Confusion}
\label{app:temp_conf}

\begin{figure}[H]
    \centering
    \includegraphics[width=0.8\linewidth]{figures/3M_2018_10k_Gold.png}
    \caption{Temporal Confusion: Gold Page}
    \label{fig:3M_gold}
\end{figure}

\begin{figure}[ht]
    \centering
    \includegraphics[width=0.8\linewidth]{figures/3M_2022_10K_Retrieved.png}
    \caption{Temporal Confusion: Retrieved Page}
    \label{fig:3M_ret}
\end{figure}

\FloatBarrier

\section{Lexical Gap}
\label{app:lex_gap}

\begin{figure}
    \centering
    \includegraphics[width=0.8\linewidth]{figures/gold_tbvps_lexical.png}
    \caption{Lexical Gap: Gold Page}
    \label{fig:JPM_gold}
\end{figure}

\begin{figure}
    \centering
    \includegraphics[width=0.8\linewidth]{figures/JMP_ret.png}
    \caption{Lexical Gap: Retrieved Page}
    \label{fig:JPM_ret}
\end{figure}
\FloatBarrier

\section{Annotation Error}
\label{app:page_ex}
\begin{figure}[H]
    \centering
    \includegraphics[width=1\linewidth]{figures/annotation_gold.png}
    \caption{Annotation Limitation: Gold Page}
    \label{fig:annot_gold}
\end{figure}

\begin{figure}[H]
    \centering
    \includegraphics[width=1\linewidth]{figures/annot_ret.png}
    \caption{Annotation Limitation: Retrieved Page}
    \label{fig:annot_ret}
\end{figure}

\chapter{Additional Results}
\label{app:page_ablation}

Table~\ref{tab:page_ablation} reports retrieval performance of the learned page scorer
as a function of the number of pages $P$ selected in Stage 1. Performance plateaus at
$P=20$, which is used as the default throughout this thesis.

\begin{table}[ht]
\centering
\caption{Page recall and chunk-level similarity of the learned page scorer as a
  function of the Stage 1 page budget $P$.}
\label{tab:page_ablation}
\begin{tabular}{ccccc}
\toprule
$P$ & DocRec@5 & PageRec@5 & BLEU@5 & ROUGE-L@5 \\
\midrule
5  & 0.92 & 0.48 & 0.30 & 0.44 \\
10 & 0.95 & 0.48 & 0.31 & 0.44 \\
20 & 0.95 & 0.55 & 0.33 & 0.46 \\
\bottomrule
\end{tabular}
\end{table}


\section{Extraction Method Comparison}
\label{app:extraction-comparison}

To select a default PDF extractor for the FinanceBench corpus, we evaluated five
tools, PyMuPDF, pdfplumber, PyPDF2, pypdf, and Marker, under identical conditions
using the same BGE-M3 dense retrieval setup and reporting PageRec@5 and PageRec@10
across the 150 FinanceBench questions.

\begin{table}[ht]
\centering
\caption{Retrieval recall by PDF extraction method on FinanceBench (BGE-M3, no
reranking). Higher is better.}
\label{tab:extraction-comparison}
\begin{tabular}{lcccc}
\toprule
\textbf{Method} & \textbf{Rec@1} & \textbf{Rec@3} & \textbf{Rec@5} & \textbf{Rec@10} \\
\midrule
pdfplumber & 0.147 & 0.267 & 0.360 & 0.500 \\
pypdf      & 0.153 & 0.267 & 0.340 & 0.487 \\
PyMuPDF    & 0.147 & 0.260 & 0.333 & 0.493 \\
PyPDF2     & 0.133 & 0.267 & 0.313 & 0.473 \\
Marker     & 0.127 & 0.280 & 0.327 & 0.460 \\
\bottomrule
\end{tabular}
\end{table}

\chapter{Chunking Ablation}
\label{app:chunking_ablation}

To isolate the effect of chunking from other design choices, and to verify the 1024-token configuration adopted from Yepes et al. [2024], we conducted a preliminary ablation holding the retrieval method fixed at BGE-M3 dense retrieval ($k=5$). We varied chunk size and overlap, and compared them against a page-boundary-aligned variant and a table-aware variant.

Table \ref{tab:chunking_sweep} reports the available completed results. The table-aware 1024-token configuration successfully retrieves the correct document for most questions (DocRec@5 of 0.913) but yields a much lower page recall at $k=5$ (0.207) compared to the dense BGE-M3 baseline (0.350) reported in Chapter 4. 

This confirms that preserving table-like regions alone is insufficient to solve the page localization bottleneck. While chunking strategy demonstrably impacts text representation, the ablation supports the thesis's central architectural premise: better page selection via a hierarchical pipeline remains a vastly stronger lever for retrieval performance than hyper-optimizing local chunk boundaries.

\begin{table}[htbp]
\centering
\caption{Preliminary chunking sweep results. Retrieval metrics are reported at $k=5$ using BGE-M3.}
\label{tab:chunking_sweep}
\begin{tabular}{lcccccc}
\toprule
\textbf{Strategy} & \textbf{Size} & \textbf{Overlap} & \textbf{DocRec@5} & \textbf{PageRec@5} & \textbf{BLEU@5} & \textbf{ROUGE-L@5} \\
\midrule
Table-aware & 1024 & 128 & 0.913 & 0.207 & 0.168 & 0.265 \\
Recursive (Default) & 1024 & 128 & 0.940 & 0.350 & 0.250 & 0.350 \\
\bottomrule
\end{tabular}
\end{table}
\end{document}

\endinput
